\documentclass[11pt]{article}

\usepackage[margin=1.02in]{geometry}
\usepackage{amsmath,amssymb,amsthm,mathtools,mathrsfs}
\usepackage{bm}
\usepackage{enumitem}
\usepackage{booktabs,array,tabularx}
\usepackage{xcolor}
\usepackage{microtype}
\usepackage[colorlinks=true,linkcolor=blue,citecolor=blue,urlcolor=blue]{hyperref}
\usepackage[nameinlink,capitalise]{cleveref}
\numberwithin{equation}{section}
\setlength{\emergencystretch}{4em}
\raggedbottom
\hbadness=10000

% theorem environments
\newtheorem{theorem}{Theorem}[section]
\newtheorem{proposition}[theorem]{Proposition}
\newtheorem{lemma}[theorem]{Lemma}
\newtheorem{corollary}[theorem]{Corollary}

\newtheorem{example}[theorem]{Example}
\theoremstyle{definition}
\newtheorem{definition}[theorem]{Definition}
\newtheorem{construction}[theorem]{Construction}
\newtheorem{assumption}[theorem]{Assumption}
\theoremstyle{remark}
\newtheorem{remark}[theorem]{Remark}


% notation
\newcommand{\N}{\mathbb N}
\newcommand{\Z}{\mathbb Z}
\newcommand{\R}{\mathbb R}
\newcommand{\C}{\mathbb C}
\newcommand{\one}{\mathbf 1}
\newcommand{\id}{\mathrm{id}}
\newcommand{\dd}{\mathrm d}
\newcommand{\Tr}{\operatorname{Tr}}
\newcommand{\tr}{\operatorname{tr}}
\newcommand{\rank}{\operatorname{rank}}
\newcommand{\Ker}{\operatorname{Ker}}
\newcommand{\Ran}{\operatorname{Ran}}
\newcommand{\Span}{\operatorname{span}}
\newcommand{\End}{\operatorname{End}}
\newcommand{\Hom}{\operatorname{Hom}}
\newcommand{\supp}{\operatorname{supp}}
\newcommand{\spec}{\operatorname{spec}}
\newcommand{\diag}{\operatorname{diag}}
\newcommand{\Aut}{\operatorname{Aut}}
\newcommand{\U}{\mathrm U}
\newcommand{\SU}{\mathrm{SU}}
\newcommand{\Cl}{\mathrm{Cl}}
\newcommand{\HS}{\mathrm{HS}}
\newcommand{\norm}[1]{\lVert#1\rVert}
\newcommand{\abs}[1]{\lvert#1\rvert}
\newcommand{\ip}[2]{\left\langle#1,#2\right\rangle}
\newcommand{\A}{\mathcal A}
\newcommand{\Hh}{\mathcal H}
\newcommand{\M}{\mathcal M}
\newcommand{\Kk}{\mathcal K}
\newcommand{\Oalg}{\mathcal O}
\newcommand{\GT}{\boldsymbol{\mathbb T}_{\mathrm{Grand}}}
\newcommand{\GTX}{\boldsymbol{\mathbb T}_{\mathrm{Grand},X}}
\newcolumntype{L}[1]{>{\raggedright\arraybackslash}p{#1}}

% Presentation revision: original mathematical scope and source hypotheses retained.
% Author affiliation was not supplied; complete it before journal submission.
\title{\textbf{Action--Multiplicity Duality for Spacetime--Gauge Incidence}}

\author{Aurelien Pelissier$^{\dagger}$\\
{\small $^{\dagger}$Correspondence: \href{mailto:aurelien.pelissier.38@gmail.com}{aurelien.pelissier.38@gmail.com}}}
\date{}

\begin{document}
\maketitle


\begin{abstract}
We prove an exact action--multiplicity duality for finite gravity--matter systems in which spacetime, internal typing, and matter incidence are represented on a common carrier. The complete represented action and the residual matter multiplicity determine one another as mutual commutants. The result applies to arbitrary incidence maps, including singular and rank-deficient ones, and admits intrinsic support, quiver, rigidity, and cofinal-stability formulations.
%
On a certified structural branch, this duality reconstructs the Standard-Model internal algebra. Finite naturality and anomaly conditions then recover the global Standard-Model gauge group and hypercharge ratios. A source-minimal fermion analysis, together with a positive endpoint-allocation criterion, yields exactly three complete charged copies while excluding future-visible charged exotic and mirror sectors. The construction is finite and operator algebraic; it does not depend on a particular microscopic realization. Renewal Geometry provides one realization of the required joint carrier.
All selection statements retain explicit occurrence and compatibility hypotheses; absent source data are not supplied by algebraic completion.
\end{abstract}


\medskip
\noindent\textbf{Keywords.}
double-centralizer duality; matter multiplicity; Standard Model;
finite operator algebras; gauge-group reconstruction; generations;
commutant rigidity; incidence geometry; Renewal Geometry.


\section{Introduction}
\label{sec:introduction}
The two central theories of fundamental physics organize structure in
different ways.  General relativity identifies gravity with the geometry of
spacetime itself, whereas the Standard Model is formulated over an already
given spacetime and introduces a separate internal gauge structure, chiral
matter representations, a scalar sector, and Yukawa couplings.  Under the
usual assumptions on an ordinary relativistic $S$-matrix, the
Coleman--Mandula theorem sharply constrains attempts to combine continuous
spacetime and internal Lie symmetries into a larger nontrivial symmetry group
\cite{ColemanMandula1967}.

This separation does not imply that spacetime action, internal symmetry, and
matter multiplicity must be algebraically unrelated.  They may instead occur
as distinct structures on a common finite carrier and constrain one another
through their joint representation.  The relevant question is then not
whether the Lorentz and Standard-Model gauge groups should be unified as
subgroups of a larger ordinary symmetry group.  It is whether the complete
algebra of represented spacetime, internal, and matter-incidence operations
determines the multiplicity left invisible to that action, and conversely
whether this residual multiplicity determines the represented action.

The present paper formulates this question in finite operator-algebraic terms.
Its central result is an intrinsic coefficient-space reconstruction and
exact double-centralizer theorem for a source-complete joint packet.  Once an external Clifford action, a typed
internal action, and physical matter-incidence operators are represented on
one finite Hilbert carrier, the algebra generated by the complete visible
action and the residual typed matter multiplicity determine one another as
mutual commutants.  The theorem depends only on this represented finite
packet and not on a particular microscopic mechanism by which the packet was
constructed.


\paragraph{Relation to earlier geometric and algebraic approaches.}

Noncommutative geometry provides a natural point of comparison because it
places spacetime and internal particle structure in a common
operator-theoretic framework.  In spectral geometry, the data
$(\A,\Hh,D)$ encode geometric information
\cite{Connes1994,ConnesReality1995}, while the Connes--Lott and
almost-commutative constructions combine ordinary spin geometry with a
finite noncommutative sector and recover gauge, Higgs, fermionic, and
gravitational terms through the spectral action
\cite{ConnesLott1991,ConnesGravityMatter1996,ChamseddineConnes1997,
ChamseddineConnesMarcolli2007,ChamseddineConnes2008,
vanSuijlekom2025}.  Finite spectral triples admit a substantial
classification theory \cite{Krajewski1998,IochumSchuckerStephan2004,
vanDenDungenVanSuijlekom2012}, and modified spectral assumptions can lead to
larger internal symmetries such as Pati--Salam models
\cite{ChamseddineConnesVanSuijlekom2013}.  Lorentzian and twisted
formulations address the corresponding indefinite geometric structures
\cite{Barrett2007,DevastatoEtAl2018}.

Particularly relevant here are constructions in which internal information is
encoded by commutants or centralizers.  Morita formulations of the
noncommutative Standard Model relate the coordinate algebra to its spectral
Clifford algebra \cite{DAndreaDabrowski2016}, while the Hodge condition asks
for exact commutant relations between spectral Clifford actions
\cite{DabrowskiDAndreaSitarz2018,DabrowskiSitarz2019}.  More generally,
classical Howe duality organizes mutually centralizing actions through their
multiplicity structure \cite{Howe1989}, and reconstruction of internal gauge
symmetry from algebraic or representation-theoretic data also appears in
algebraic quantum field theory
\cite{DoplicherRoberts1989,DoplicherRoberts1990,Todorov2010}.

A complementary class of approaches seeks Standard-Model structure inside
larger Clifford, division, or exceptional algebras
\cite{TraylingBaylis2001,Daviau2017,Gresnigt2020,Krasnov2021,
BaezSchwahn2026,FureyRoadmap2025}.  These constructions typically identify
a distinguished centralizer, stabilizer, or subalgebra inside a prescribed
ambient algebra.

The logical structure considered here is different.  The primitive datum for
the main theorem is not a larger unifying algebra and is not a preassigned
finite spectral triple.  It is a \emph{joint represented packet}: an external
Clifford factor, a commuting typed internal algebra, and physical incidence
operators acting on the same carrier.  The algebra dual to the complete
action is also not the Standard-Model gauge algebra itself.  It is the
residual matter multiplicity that the complete represented action cannot
resolve.


\paragraph{Finite action--multiplicity duality.}

Let $\A_{{\rm Cl},X}\cong M_4(\C)$ be the represented external Clifford
factor, let $\A_{{\rm type},X}$ be a commuting finite-dimensional kinematic
typing algebra, and let $Y_{e,X}$ denote represented matter-incidence
operators.  The complete action algebra is
\begin{equation}
\boxed{
\A_{{\rm act},X}
=
C^*\!\left(
\A_{{\rm Cl},X},
\A_{{\rm type},X},
Y_{e,X},Y_{e,X}^*
\right).}
\label{eq:intro-action-algebra}
\end{equation}
Finite representation theory separates the carrier into external,
internal-representation, and residual multiplicity factors.  The operators
on the residual factors that intertwine the intrinsic coefficient spaces of
all represented incidence maps form the typed multiplicity algebra
$\M_{{\rm type},X}$.

The main finite theorem proves
\begin{equation}
\boxed{
\A_{{\rm act},X}'
=
\M_{{\rm type},X},
\qquad
\M_{{\rm type},X}'
=
\A_{{\rm act},X}.}
\label{eq:intro-main-duality}
\end{equation}
Thus the duality is between the complete
spacetime--gauge--incidence action and residual matter multiplicity.  It is
exact for arbitrary operator-Schmidt rank, including rectangular, singular,
rank-deficient, and zero incidence maps.  Intrinsic coefficient-space quivers
require no one-dimensional intertwiner assumption.  Support-polar data give a support-exact formulation without introducing inverse
incidence operators; noncommuting supports lead to a multiplicity quiver; and
on connected invertible branches the same residual algebra is described by
positive polar metrics and cycle holonomies.  This is a finite double-centralizer duality, without assuming a classical
reductive dual pair.

The incidence family itself has a finite minimality theory.  A submodular
rank function measures how much each represented incidence removes from the
residual commutant, yielding inclusion-minimal incidence skeletons.  This
makes it possible to distinguish structural incidence data from redundant
generators and to identify precisely when deleting one physical incidence
enlarges the multiplicity algebra.


\paragraph{One upstream realization: Renewal Geometry.}

Predictive-state approaches provide one route to finite structure without
introducing a manifold or field configuration at the outset: state is
characterized by the information required for future prediction rather than
by an independently postulated hidden coordinate
\cite{ShaliziCrutchfield2001}.  Multitime operational descriptions with
memory provide a related framework for intervention-dependent processes
\cite{ChiribellaEtAl2009,PollockEtAl2018}.  One predictive realization of the finite joint-carrier input is developed
through the Lorentzian and spectral Renewal Geometry constructions
\cite{PelissierSpacetime2026,PelissierSpectral2026}, while the
double-centralizer theorem itself is independent of that microscopic
construction.

In this realization, operational histories are quotiented by equality of
conditional future statistics.  The resulting minimum persistent predictive
carrier $Z_X^{\min}$, represented history algebra
$\mathcal A_X^{\mathrm{hist}}$, and positive word-Gram hierarchy are organized
into the \emph{Renewal Grand Tensor},
\begin{equation}
\boxed{
\GTX
=
\left(
\mathcal A_X^{\mathrm{hist}},
\{\mathbb K_{X,r}\}_{r\ge0},
Z_X^{\min}
\right).}
\label{eq:intro-grand-tensor}
\end{equation}
Here ``tensor'' denotes this finite predictive data structure rather than a
tensor field on a pre-existing manifold.  The predictive carrier, represented
history algebra, and word-Gram framework underlying this object are developed
in \cite{PelissierSpacetime2026,PelissierSpectral2026}.  The Lorentzian
construction \cite{PelissierSpacetime2026} extracts, on a certified branch, a
tetrahedral relational sector, calibrated geometric data, and a faithful
complex Clifford factor, while the spectral construction
\cite{PelissierSpectral2026} supplies a compatible history-algebraic and
first-order spectral realization.  These results furnish one route from the
Grand Tensor to the represented input of the present theory; they are not
hypotheses of the finite duality theorem.


\paragraph{Structural Standard-Model specialization.}

The general double-centralizer theorem does not select the Standard Model.
A second layer of the paper specializes to a tetrahedrally equivariant
source-complete branch and asks which residual multiplicity structures survive
a finite sequence of occurrence, connectivity, positivity, and incidence
tests.

On the certified branch the landed internal algebra is
\begin{equation}
\boxed{
\A_{\rm int}^{\rm word}
\cong
M_3(\C)\oplus M_2(\C)\oplus\C\oplus\C,}
\label{eq:intro-internal-algebra}
\end{equation}
and a canonical incidence line, selected by six finite naturality identities
and retained relative to a coherent neutral reference, selects
\begin{equation}
\boxed{
S(\U(3)\times\U(2))
\cong
\frac{\SU(3)\times\SU(2)\times\U(1)}{\Z_6}.}
\label{eq:intro-gauge-group}
\end{equation}
The determinant-source route therefore obtains the global quotient at the
level of represented incidence rather than from the Lie algebra alone
\cite{Hucks1991,Tong2017}.  Before that naturality test is applied, a finite
commutator residual certifies that the represented colour and two weak factor
actions occupy mutually commuting tensor slots on one carrier.  Once this
factor skeleton is fixed, the naturality operator and its determinant projector
are defined independently of any matter-source population; positive source mass
and vanishing naturality defect are separate occurrence and purity tests.
Independently, once the tensor types are fixed, anomaly cancellation determines
the same primitive central embedding and, together with the faithful matter
kernel, reconstructs the same global group without assuming determinant-source
occurrence.  It also gives the usual Standard-Model hypercharge ratios
\cite{GengMarshak1989,MinahanRamondWarner1990,
AlvarezGraciaBondiaMartin1995}.

At the process level, a Kraus-independent cross-sign Choi mass tests primitive
grading-changing occurrence before a coherent odd representative is chosen.
This active mass is unchanged by unread trace-preserving spectator extensions,
so hidden spectator oddness cannot substitute for matter occurrence on the
represented carrier.

The same structural branch contains a distinct rank-three endpoint residual.
After finite $*$-cyclic saturation of the physical fermion source, canonical
isotypic projectors audit the complete future-visible gauge-charged carrier.
On the zero-inventory branch four linear anomaly relations force the five
charged multiplicities to coincide.  A finite Schur-covariance test extracts a
unique multiplicity map from a represented endpoint-to-type ancestry
occurrence.  A positive allocation short then measures the excess of one
complete typed source beyond that rank-three endpoint image; its vanishing
yields exactly three copies of every charged matter type.  This already
gives an abstract common generation factor; common endpoint ancestry is needed
only when a particular source-determined relative generation frame is claimed.
A separate positive weak-copy Gram counts actual weak-doublet copies, while a
finite relation-space defect tests the allowed Yukawa and optional Majorana
blocks of a specified grading-odd generator.  These statements are structural
rather than numerical: they do not fix fermion masses or mixing parameters,
gauge couplings, Higgs parameters, or renormalization-group data.

Importantly, the Standard-Model branch is reached through finite alternatives
rather than by completing a desired answer.  Before a multiplicity of three
appears, an internal colour factor is algebraically impossible; if the word
hierarchy becomes flat before that occurrence, the obstruction is terminal.
Once the required multiplicity spaces exist, failure of the colour bridge,
weak router, determinant incidence, or other structural tests returns the
corresponding smaller algebra instead of supplying the missing structure by
assumption.


\paragraph{Rigidity, refinement, and variational compatibility.}

The paper also addresses whether the finite duality is stable rather than
accidental.  The commutant Laplacian gives a positive finite certificate whose
kernel is exactly the commutant, and the associated relative Howe Gram
quantifies the gap to additional multiplicity directions.  For polynomial or
analytic parameter families, its determinant describes the
symmetry-enhancement locus.  A positive finite gap is stable under controlled
perturbations.

A separate cofinal result treats regulator refinement.  A cutoff-uniform
commutator gap on the complement of a finitely transported protected algebra,
together with transported Mosco convergence, forces both the late-cutoff and
limiting commutant kernels to equal the protected multiplicity algebra; limiting
kernel exactness is therefore a conclusion rather than an independent
hypothesis.  Asymptotic compactness then upgrades this statement to
norm-resolvent, eigenvalue, and spectral-projection convergence.  A common
abstract algebraic identification additionally requires multiplicatively
compatible transports.

Finally, the common source quotient separates internal directions from an
already represented gravitational source.  This is logically separate from
the double-centralizer theorem and does not construct a gravity--matter
action or a continuum Einstein--Standard-Model theory.  The nonlinear
finite-to-continuum closure problem requires additional variational,
compactness, and PDE analysis.

The principal operator-algebraic result is therefore general at the level of
finite joint packets.  The tetrahedral Standard-Model branch and Renewal
Geometry realization provide progressively more specialized structures built
around that theorem rather than premises required for it.

\paragraph{Organization.}
Section~\ref{sec:joint-interface} specifies the joint packet and separates
represented input from its operational identification.
Section~\ref{sec:finite-duality} proves the exact action--multiplicity theorem,
then develops its principal consequences and incidence minimality.
Section~\ref{sec:internal-algebra} reconstructs the internal algebra on the
specified tetrahedral branch.
Section~\ref{sec:gauge-matter} reconstructs the global gauge group and matter
typing algebra.  Section~\ref{sec:charged-matter} treats the charged inventory,
endpoint allocation, and the resulting generation structure.
Section~\ref{sec:certificate} establishes finite rigidity and cofinal stability.
Section~\ref{sec:discussion} discusses interpretation and scope.
Appendices~\ref{app:commutant}--\ref{app:stability} retain the supporting
operator-algebraic, source-identification, representation-theoretic,
occurrence, and analytic details in the same article.


\section{Joint packets and the commutant obstruction}
\label{sec:joint-interface}

The finite duality proved below is formulated directly in terms of represented
operator-algebraic data.  No predictive reconstruction, spectral triple, or
microscopic stochastic law is part of the theorem statement.  Such structures
may provide realizations of the required finite packet, but the
double-centralizer argument uses only the common carrier, the external and
internal actions, and the represented incidence operators specified here.

\subsection{Source-complete joint packets}

\begin{definition}[Finite source-complete joint packet]
\label{def:joint-packet}
At cutoff $X$, a finite source-complete spacetime--gauge packet consists of:
\begin{enumerate}[label=\textup{(J\arabic*)},leftmargin=3em]
\item a finite Hilbert carrier $\Hh_X$ on which a faithfully represented
      external Clifford algebra
      $\A_{{\rm Cl},X}\cong M_4(\C)$ acts;
\item a commuting finite-dimensional kinematic typing $C^*$-algebra
      $\A_{{\rm type},X}$;
\item a central type decomposition
      \[
      \Hh_X
      \cong
      \bigoplus_{\lambda\in\Lambda_X}
      \C^4\otimes V_{\lambda,X}\otimes M_{\lambda,X},
      \]
      such that
      \[
      \A_{{\rm type},X}
      =
      \bigoplus_{\lambda}
      I_4\otimes\mathcal B(V_{\lambda,X})\otimes I_{M_{\lambda,X}};
      \]
\item a finite directed incidence graph
      $\Gamma_X=(\Lambda_X,E_X)$ and represented incidence operators
      \[
      Y_{e,X}:
      \C^4\otimes V_{s(e),X}\otimes M_{s(e),X}
      \longrightarrow
      \C^4\otimes V_{t(e),X}\otimes M_{t(e),X};
      \]
\item the mixed represented products needed to determine products and adjoints
      of $\A_{{\rm Cl},X}$, $\A_{{\rm type},X}$, and the $Y_{e,X}$ on the same
      carrier.
\end{enumerate}
All operators entering a claimed mixed relation are required to be represented
on this common carrier.  No inverse incidence map is part of the definition.
\end{definition}

\begin{definition}[Joint typed action algebra]
\label{def:action-algebra}
For a joint packet define
\begin{equation}
\A_{{\rm act},X}
:=
C^*\!\left(
\A_{{\rm Cl},X},
\A_{{\rm type},X},
Y_{e,X},Y_{e,X}^*:e\in E_X
\right).
\label{eq:joint-action-algebra}
\end{equation}
The notation ``action'' refers here to represented operator action on the
common finite carrier, not to a variational action functional.
\end{definition}

The common carrier is essential: replacing it by a direct sum of unrelated
external and internal packets removes the mixed incidence relations.  We first
identify the freedom left by the external and typing actions alone.


\subsection{External, internal, and multiplicity factors}
\label{sec:obstruction}

\begin{proposition}[Factor obstruction and typed decomposition]
\label{thm:trinity}
Let a unital copy of $M_n(\C)$ act on a finite Hilbert space $\Hh$.
For a multiplicity space $\mathcal K$,
\begin{equation}
 \Hh\cong\C^n\otimes\mathcal K,\qquad
 M_n(\C)'=I_n\otimes\mathcal B(\mathcal K).
 \label{eq:factor-commutant}
\end{equation}
Thus a reducible internal algebra with nontrivial centre cannot be the full
commutant of the external factor.  If $\A_{\rm type}$ is a commuting unital
finite-dimensional $C^*$-algebra, there are spaces $V_\lambda,M_\lambda$ with
\begin{align}
 \Hh&\cong\bigoplus_\lambda\C^n\otimes V_\lambda\otimes M_\lambda,
 \label{eq:trinity-H}\\
 \A_{\rm type}&=\bigoplus_\lambda I_n\otimes\mathcal B(V_\lambda)
                         \otimes I_{M_\lambda}.
 \label{eq:trinity-Aint}
\end{align}
Writing $\mathcal C=C^*(\A_{\rm Cl},\A_{\rm type})$, one has
\begin{align}
 \mathcal C&=\bigoplus_\lambda\mathcal B(\C^n\otimes V_\lambda)
                         \otimes I_{M_\lambda},\label{eq:trinity-C}\\
 \mathcal C'&=\bigoplus_\lambda I_{\C^n\otimes V_\lambda}
                         \otimes\mathcal B(M_\lambda).
 \label{eq:trinity-Cprime}
\end{align}
The residual spaces $M_\lambda$, not the internal representation spaces
$V_\lambda$, carry flavour and generation multiplicity.
\end{proposition}
\begin{proof}
Every finite unital representation of $M_n(\C)$ is a direct sum of its
defining representation, giving $\Hh=\C^n\otimes\mathcal K$ and the first
commutant identity by matrix units.  The commuting typing algebra acts on
$\mathcal K$.  Its finite-dimensional representation decomposition gives
$\mathcal K=\bigoplus_\lambda V_\lambda\otimes M_\lambda$ and
\eqref{eq:trinity-Aint}.  On each central block the two full matrix factors
generate $\mathcal B(\C^n\otimes V_\lambda)$; Schur's lemma gives
\eqref{eq:trinity-Cprime}.
\end{proof}

The preceding decomposition identifies everything invisible to the bare
external--typing action.  The remaining question is how the represented
incidences constrain the spaces $M_\lambda$.  Section~\ref{sec:finite-duality}
answers it without assuming that an incidence factors as a single tensor.


\subsection{Represented input and source identification}
\label{sec:operational-identification}

The joint packet is the input to the finite theorem, not its output.  A
microscopic realization must identify the operations acting on the system
carrier, including their mixed products.  Classical records alone need not
suffice: Proposition~\ref{thm:record-nonidentifiability} constructs two
normalized finite instruments with the same resolved records at every word
length, one-step effects, coefficient Gram, and nonzero coarse Choi spectrum,
but different system algebras and commutants.

There is also a positive reconstruction route.  For an isometry
$V=\sum_\nu C_\nu\otimes|\nu\rangle$, the complementary Heisenberg pullback
satisfies $\mathcal P^*(|z\rangle\langle w|)=C(z)^*C(w)$.  If occurring
coefficients include a nonzero identity pivot and nonzero multiples of a
system-algebra generating family, these pullbacks generate exactly that
system algebra (Proposition~\ref{thm:pullback-reconstruction}).
Appendix~\ref{app:source-identifiability} proves both statements and gives the
finite coherent reconstruction panel.

Renewal Geometry is one source of such represented data.  Its Lorentzian
branch supplies, under the hypotheses of the cited construction,
\begin{equation}
 \A_{{\rm Cl},X}=C^*(\gamma^0,\gamma^1,\gamma^2,\gamma^3)\cong M_4(\C),
 \label{eq:external-clifford}
\end{equation}
while its spectral branch gives a compatible history-algebraic realization
\cite{PelissierSpacetime2026,PelissierSpectral2026}.  Neither construction is a
premise of the finite theorem.  For historical realizations,
Appendix~\ref{sec:historical-operation-recognition} additionally distinguishes
executable system operations from invariant endomorphisms of a regular word
space, and specifies when mixed source Grams are computable from existing
word moments.  These identification questions do not interrupt the algebraic
argument once the joint packet is supplied.


\section{Exact action--multiplicity duality}
\label{sec:finite-duality}

Throughout this section the represented joint packet is fixed and the cutoff
index is suppressed.  The proof proceeds in three steps: extract the full
multiplicity coefficient space of each incidence, identify its intertwining
algebra, and apply finite bicommutant closure.


\subsection{Intrinsic coefficient spaces and the reconstruction theorem}

Put $K_\lambda=\C^4\otimes V_\lambda$ and
$\mathscr M=\bigoplus_\lambda M_\lambda$.  Let $Z_\lambda$ be the
vertex projections on $\mathscr M$, and let
\[
 \iota((R_\lambda)_\lambda)
 =\bigoplus_\lambda I_{K_\lambda}\otimes R_\lambda
\]
be the faithful representation of the block-diagonal multiplicity operators
on $\Hh$.  Commutants on $\Hh$ and on $\mathscr M$ are distinguished below;
when the carrier is unambiguous, the representation $\iota$ is suppressed.

\begin{definition}[Intrinsic incidence coefficient space]
\label{def:incidence-coefficients}
For an arbitrary represented edge
$Y_e:K_\lambda\otimes M_\lambda\to K_\mu\otimes M_\mu$, define
\begin{equation}
 \mathfrak B_e
 =\{(\varphi\otimes\id)(Y_e):
     \varphi\in\mathcal B(K_\lambda,K_\mu)^*\}
 \subseteq\mathcal B(M_\lambda,M_\mu).
 \label{eq:intrinsic-incidence-coefficients}
\end{equation}
This is a linear operator space, not a choice of a single residue.
Equivalently, for a Hilbert--Schmidt orthonormal basis
$(D_{e,a})_a$ of $\mathcal B(K_\lambda,K_\mu)$, write
\begin{equation}
 Y_e=\sum_a D_{e,a}\otimes F_{e,a},
 \qquad \mathfrak B_e=\Span\{F_{e,a}:a\}.
 \label{eq:general-incidence-expansion}
\end{equation}
The coefficient-space quiver has vertices $M_\lambda$, forward arrow
spaces $\mathfrak B_e$, and the corresponding adjoint arrow spaces.
\end{definition}

\begin{definition}[Typed multiplicity algebra]
\label{def:typed-multiplicity}
The endomorphism algebra of this adjoint-closed coefficient-space quiver is
\begin{equation}
 \begin{aligned}
 \M_{\rm type}=\{(R_\lambda)_\lambda:\ &R_\mu B=BR_\lambda,\quad
 R_\lambda B^*=B^*R_\mu\\[-.2em]
 &\text{for every }B\in\mathfrak B_e,
                   \ e:\lambda\to\mu\}.
 \end{aligned}
 \label{eq:typed-multiplicity}
\end{equation}
It is a unital $C^*$-algebra on $\mathscr M$, represented on the full carrier
by $\iota$.  Define also the multiplicity coefficient action algebra
\begin{equation}
 \Oalg_Y=C^*(Z_\lambda,B,B^*:B\in\mathfrak B_e,\ e\in E)
 \subseteq\mathcal B(\mathscr M).
 \label{eq:coefficient-action-algebra}
\end{equation}
\end{definition}

Commutation with the external and typing algebras has already reduced an
operator to the multiplicity factors.  The coefficient spaces now retain
exactly the additional information needed to commute with every physical
incidence and its adjoint.

\begin{theorem}[Arbitrary-incidence action--multiplicity reconstruction]
\label{thm:main-duality}
For every finite source-complete joint packet of
Definition~\ref{def:joint-packet}, with no restriction on the
operator-Schmidt rank of its incidences,
\begin{equation}
 \A_{\rm act}'{}_{\Hh}=\iota(\M_{\rm type}),\qquad
 \iota(\M_{\rm type})'{}_{\Hh}=\A_{\rm act}.
 \label{eq:main-duality}
\end{equation}
On the total multiplicity carrier,
\begin{equation}
 \Oalg_Y'{}_{\mathscr M}=\M_{\rm type}
   =\End\mathfrak Q_Y,\qquad
 \M_{\rm type}'{}_{\mathscr M}=\Oalg_Y,
 \label{eq:coefficient-quiver-duality}
\end{equation}
where $\mathfrak Q_Y$ is the adjoint-closed coefficient-space quiver.
Both algebras are independent of all coefficient bases and of the chosen
Wedderburn coordinates up to the canonical multiplicity unitaries of
Corollary~\ref{cor:typing-coordinate-canonicality}.  These conclusions include
rectangular, singular, rank-deficient, and zero incidences.

\end{theorem}

\begin{proof}
By Proposition~\ref{thm:trinity}, every operator commuting with
$\A_{\rm Cl}$ and $\A_{\rm type}$ is uniquely
$X=\iota((R_\lambda)_\lambda)$.  On an edge $e:\lambda\to\mu$,
\[
 XY_e-Y_eX
 =\sum_aD_{e,a}\otimes
          (R_\mu F_{e,a}-F_{e,a}R_\lambda).
\]
Linear independence of the $D_{e,a}$ makes this commutator vanish exactly
when the forward coefficient equations hold.  Applying the same argument
to $Y_e^*$ gives the adjoint equations in
\eqref{eq:typed-multiplicity}.  This proves the first identity in
\eqref{eq:main-duality}; finite bicommutant closure proves the second.

On $\mathscr M$, commutation with every $Z_\lambda$ is equivalent to block
diagonality.  Commutation with all coefficient arrows and their adjoints is
then exactly \eqref{eq:typed-multiplicity}, proving the first identity in
\eqref{eq:coefficient-quiver-duality}.  Finite bicommutant closure again
proves the reverse identity.  A change of carrier basis changes the
coefficient list by an invertible linear transformation and leaves its
span, hence every defining equation, unchanged.  The same is true of a
change of spanning family within an arrow space.  None of these arguments
uses an inverse coefficient or a positive-rank assumption.
\end{proof}

This is the algebraic core of the paper.  The constructive content is the
intrinsic coefficient-space description of the first commutant; the reverse
identity follows from finite-dimensional bicommutant closure.  The result
recovers a represented algebra, not a preferred incidence list, Clifford
marking, or numerical Yukawa matrix.

The following example shows why retaining all coefficient directions matters.
A rank-one replacement can enlarge the commutant even when it preserves a
nonzero part of the same physical incidence.

\begin{example}[Two coefficients impose an additional multiplicity relation]
\label{ex:two-coefficient-incidence}
Take two types with $K_1=K_2=\C^4$ and $M_1=M_2=\C^2$, and the base
algebra $M_4(\C)\oplus M_4(\C)$ acting trivially on the multiplicities.
On the single physical edge put
\[
 Y=E_{11}\otimes I_2+E_{22}\otimes D,
 \qquad D=\diag(1,2).
\]
Its intrinsic coefficient space is $\Span\{I_2,D\}$.  The identity
coefficient forces $R_2=R_1$, and the second coefficient forces
$[R_1,D]=0$.  Hence the residual algebra is the common diagonal $\C^2$.
Keeping only the identity coefficient would instead leave a common
$M_2(\C)$ commutant.  This example is a source-complete joint packet:
the external Clifford factor is the same $M_4$ on both types, and the
type projections separate them.  It shows why the unrestricted theorem
must retain an entire coefficient space rather than assume rank-one
factorization of every physical edge.
\end{example}


\subsection{Canonicality and reciprocal structure}

\begin{corollary}[Canonical coefficient action and reciprocal blocks]
\label{cor:reciprocal-wedderburn}
The coefficient-action algebra $\Oalg_Y$ is the smallest unital
$C^*$-algebra on $\mathscr M$ containing the vertex projections and all
multiplicity slices.  For any unital finite-dimensional algebra
$\mathcal C\subseteq\mathcal B(\mathscr M)$,
\begin{equation}
 \mathcal C'=\M_{\rm type}\quad\Longleftrightarrow\quad\mathcal C=\Oalg_Y.
 \label{eq:coefficient-action-universal}
\end{equation}
A change of Wedderburn coordinates acts by multiplicity unitaries
$W_\lambda$, with
$\widetilde{\mathfrak B}_e=W_{t(e)}\mathfrak B_eW_{s(e)}^*$.
Thus the dual pair is intrinsic up to these unitary identifications.

On the full carrier there is a simultaneous decomposition
\begin{equation}
 \Hh_X\cong\bigoplus_\alpha H^{\rm act}_{\alpha,X}
                             \otimes H^{\rm mult}_{\alpha,X},
 \label{eq:wedderburn-H}
\end{equation}
with
\begin{align}
 \A_{{\rm act},X}&\cong\bigoplus_\alpha
   \mathcal B(H^{\rm act}_{\alpha,X})\otimes I,
   \label{eq:wedderburn-action}\\
 \M_{{\rm type},X}&\cong\bigoplus_\alpha
   I\otimes\mathcal B(H^{\rm mult}_{\alpha,X}).
   \label{eq:wedderburn-multiplicity}
\end{align}
Consequently
\begin{equation}
 Z(\A_{{\rm act},X})=Z(\M_{{\rm type},X})
                  =\A_{{\rm act},X}\cap\M_{{\rm type},X}.
 \label{eq:common-centre}
\end{equation}
The common central projections label the same sectors on both sides.
\end{corollary}
\begin{proof}
Slicing any realization of $Y_e$ through a coefficient subspace shows that
it contains $\mathfrak B_e$; expansion in a carrier basis gives the converse.
The algebraic minimality follows by generation, and uniqueness from the
commutant follows from Theorem~\ref{thm:main-duality}.  Coordinate changes
are absorbed by the full family of carrier linear functionals and by
multiplicity unitaries, as detailed in
Corollary~\ref{cor:typing-coordinate-canonicality}.
The finite-dimensional structure theorem applied to $\A_{{\rm act},X}$ gives
the simultaneous blocks and their common centre.
\end{proof}

The coefficient spaces also have a canonical positive realization.  This is
useful later because it separates exact algebraic information from the
stability of its numerical presentation.

\begin{proposition}[Slice Gram and operator--Schmidt spectrum]
\label{prop:incidence-slice-Gram}
For one represented incidence $Y$, write $\mathfrak B_Y$ for its intrinsic coefficient space.  Thus, for
$Y\in\mathcal B(K_s,K_t)\otimes\mathcal B(M_s,M_t)$, equip both operator
spaces with their Hilbert--Schmidt products.  For
$D\in\mathcal B(K_s,K_t)$ put
$\varphi_D(A)=\Tr(D^*A)$ and define
\[
 \mathcal S_Y(\overline D):=(\varphi_D\otimes\id)(Y),
 \qquad
 H_Y:=\mathcal S_Y\mathcal S_Y^*\succeq0.
\]
Then
\begin{equation}
 \Ran H_Y=\Ran\mathcal S_Y=\mathfrak B_Y.
 \label{eq:incidence-slice-Gram-range}
\end{equation}
If
$Y=\sum_{k=1}^r\sigma_kD_k\otimes F_k$ is an operator--Schmidt decomposition
with orthonormal coefficient families, then
\begin{equation}
 H_YF_k=\sigma_k^2F_k,\qquad
 \rank H_Y=r,\qquad
 \Tr H_Y=\norm{Y}_{\HS}^2.
 \label{eq:incidence-slice-Gram-spectrum}
\end{equation}
Thus the intrinsic coefficient space is the support of a canonical positive
operator and its positive spectrum is the squared operator--Schmidt spectrum.
\end{proposition}

\begin{proof}
In any Hilbert--Schmidt orthonormal carrier basis $(D_a)$ write
$Y=\sum_aD_a\otimes F_a$.  Then
$\mathcal S_Y(\overline{D_a})=F_a$, so its range is exactly the intrinsic
slice space.  Finite-dimensional range equality
$\Ran(SS^*)=\Ran S$ gives the first identity.  In a Schmidt basis,
$\mathcal S_Y(\overline{D_k})=\sigma_kF_k$, which proves the spectral and
rank statements.  Finally
$\|\mathcal S_Y\|_{\HS}^2=\sum_a\|F_a\|_{\HS}^2=\|Y\|_{\HS}^2$.
\end{proof}

For an orthonormal expansion $Y_e=\sum_aD_{e,a}\otimes F_{e,a}$, the
same calculation as in the main proof gives
\begin{equation}
 \|[\iota(R),Y_e]\|_{\HS}^2
 =\sum_a\|R_{t(e)}F_{e,a}-F_{e,a}R_{s(e)}\|_{\HS}^2,
 \label{eq:coefficient-commutator-form}
\end{equation}
and $\|\iota(R)\|_{\HS}^2=
\sum_\lambda(\dim K_\lambda)\|R_\lambda\|_{\HS}^2$.
The adjoint incidence adds the adjoint constraints.  These vertex weights
enter the spectral problem, although they do not change its kernel.


\subsection{Support-polar and quiver formulations}

Choose a finite bank $(F_e)$ spanning the intrinsic arrow spaces, with
adjoints.  This is a presentation choice: coefficient recombination can
change individual support projections, but not the generated coefficient
algebra.  For $F_e=U_eP_e$, retain the polar partial isometry and the positive
metric rather than an inverse of the incidence.

\begin{definition}[Support-polar action algebra]
\label{def:support-polar}
Define
\begin{equation}
\Oalg_F^{\rm sup}
=
C^*(Z_\lambda,F_e^*F_e,U_e,U_e^*:\lambda,e)
\subseteq\mathcal B(\mathscr M).
\label{eq:support-polar}
\end{equation}
\end{definition}

\begin{proposition}[Exact support-polar duality]
\label{thm:support-polar}
For arbitrary rectangular, singular, rank-deficient, or zero incidence maps,
\begin{equation}
\M_{\rm type}
=
(\Oalg_F^{\rm sup})',
\qquad
\Oalg_F^{\rm sup}
=
\M_{\rm type}'.
\label{eq:support-polar-duality}
\end{equation}
Moreover,
\begin{equation}
C^*(Z_\lambda,F_e,F_e^*:\lambda,e)
=
\Oalg_F^{\rm sup}.
\label{eq:incidence-support-polar-equality}
\end{equation}
Thus a zero singular value remains a genuine support boundary and is never
replaced by an inverse residue.  Proposition~\ref{prop:singular-polar-energy}
quantifies leakage across this boundary, and
Proposition~\ref{thm:certified-support-stratum} gives a finite support certificate
when a source-side rank ceiling is available.
\end{proposition}

\begin{proof}
Commutation with the vertex projections makes an operator block diagonal.
Lemma~\ref{lem:polar-edge} shows that commutation with
$F_e^*F_e,U_e,U_e^*$ is exactly the pair of incidence intertwining equations
in Definition~\ref{def:typed-multiplicity}.  This gives the first equality,
and finite bicommutant closure gives the second.

For \eqref{eq:incidence-support-polar-equality}, $F_e^*F_e$ belongs to the edge
algebra.  Since its spectrum is finite, interpolate the function
$g(0)=0$, $g(t)=t^{-1/2}$ for $t>0$ on that spectrum by a polynomial.
Then
\[
U_e=F_eg(F_e^*F_e)
\]
belongs to the incidence algebra.  Conversely
$F_e=U_e(F_e^*F_e)^{1/2}$.
\end{proof}

The singular-edge equations and their support-leakage energy are proved in
Appendix~\ref{app:polar-details}.  When supports do not commute, their
isotypic decomposition gives a refined operator-space quiver with
\begin{equation}
 C^*(\mathcal S_F,\bm T)'\cong\End\mathfrak Q_F(\bm T).
 \label{eq:quiver-main-summary}
\end{equation}
Its full definition and commutator form are in
Proposition~\ref{thm:quiver}.  On a connected bank of invertible arrows,
Proposition~\ref{thm:polar-holonomy} instead transports all multiplicities to
a root fibre: the residual algebra is the commutant of the transported
positive polar metrics and cycle holonomies.  These are formulations of the
same duality, not additional conditions on the general theorem.


\subsection{Incidence essentiality and minimal skeletons}
\label{sec:incidence-skeleton}

Fix the non-incidence action
$\mathcal A_0=C^*(\mathcal A_{\rm Cl},\mathcal A_{\rm type})$ and a finite
incidence bank $E$.  For $S\subseteq E$ set
\[
 \M(S)=C^*(\mathcal A_0,Y_e,Y_e^*:e\in S)',
 \qquad \M=\M(E),
\]
and define the residual test space
$V=\mathcal A_0'\cap\M^\perp$, with $n=\dim_\C V$.  For $X\in V$ let
\begin{equation}
 D_eX=2^{-1/2}([Y_e,X],[Y_e^*,X]),
 \qquad G_e=D_e^*D_e,
 \qquad G_S=\sum_{e\in S}G_e.
 \label{eq:incidence-forms}
\end{equation}
The full bank is exact precisely when $G_E\succ0$ on $V$.

\begin{proposition}[Incidence rank polymatroid]
\label{thm:incidence-polymatroid}
Let $R_e=\Ran D_e^*\subseteq V$ and
\begin{equation}
 f(S)=\dim_\C\!\left(\sum_{e\in S}R_e\right).
 \label{eq:incidence-rank}
\end{equation}
Then
\begin{equation}
 f(S)=\rank G_S
     =n-\dim_\C(\M(S)\cap V),
 \qquad
 \M(S)=\M\oplus\Ker G_S
 \label{eq:incidence-kernel}
\end{equation}
as orthogonal vector spaces.  The function $f$ is normalized, monotone and
submodular:
\begin{equation}
 f(S)+f(T)\ge f(S\cup T)+f(S\cap T).
 \label{eq:incidence-submodular}
\end{equation}
An incidence $e\in E$ is essential for the full duality exactly when
\begin{equation}
 f(E\setminus\{e\})<n
 \label{eq:essential-incidence}
\end{equation}
or, equivalently, when deleting $e$ strictly enlarges the commutant.
\end{proposition}

\begin{proof}
The quadratic form of $G_S$ is
$\sum_{e\in S}\|D_eX\|^2$, hence
\[
 \Ker G_S=\bigcap_{e\in S}\Ker D_e
 =\left(\sum_{e\in S}\Ran D_e^*\right)^\perp.
\]
Every element of $\M(S)$ already lies in $\mathcal A_0'$.  Orthogonal
projection onto $\M$ therefore leaves a residual vector in $V$, and the
incidence constraints on that residual are exactly $G_SX=0$.  This proves
\eqref{eq:incidence-kernel}.  If $R(S)=\sum_{e\in S}R_e$, then
$R(S)+R(T)=R(S\cup T)$ and
$R(S\cap T)\subseteq R(S)\cap R(T)$.  The dimension identity for two
subspaces yields submodularity.  The deletion criterion is the kernel formula
with $S=E\setminus\{e\}$.
\end{proof}

\begin{corollary}[Finite source-minimal incidence skeleton]
\label{cor:minimal-incidence-skeleton}
Every finite exact incidence bank contains an inclusion-minimal subset
$S\subseteq E$ with $\M(S)=\M(E)$.  Such a subset is obtained by repeatedly
deleting nonessential incidences and has at most $n$ members.  In general
these minimal subsets need not be unique and need not have equal cardinality;
the correct grouped rank object is the polymatroid $f$, not a matroid on the
incidence labels.
\end{corollary}

\begin{proof}
Finite reverse deletion terminates.  In an inclusion-minimal sufficient bank,
adding its members one by one strictly increases the span of the corresponding
constraint ranges, so at most $n$ additions are possible.  The remaining
claims follow from the grouped nature of the subspaces $R_e$; an explicit
unequal-cardinality example is given in Appendix~\ref{app:incidence-details}.
\end{proof}

The grouped rank is a polymatroid rather than, in general, a matroid on the
physical incidence labels.  Appendix~\ref{app:incidence-details} gives the
unequal-cardinality example, a positive determinant refinement, and exact
simultaneous-deletion margins.  The ordering and typing-refinement
consequences of the duality are collected in
Appendix~\ref{app:coefficient-consequences}.

The general theory is now complete.  We next impose additional finite
source and representation hypotheses to identify a particular internal
algebra; those hypotheses are not required for Theorem~\ref{thm:main-duality}.


\section{Reconstruction of the structural internal algebra}
\label{sec:internal-algebra}

The general duality determines the multiplicity left by any represented
joint packet.  It does not select a gauge theory.  We now consider the
specified tetrahedrally equivariant branch and ask when its internal word
algebra saturates a particular relative commutant.  The argument separates
three questions: which source directions survive removal of geometry, which
multiplicity spaces occur on the physical carrier, and which represented
operations connect those spaces.


\subsection{Geometry shorting and the endpoint residual}

The oriented $A_3$ root carrier used in this specialization is the tetrahedral
relational carrier reconstructed in \cite{PelissierSpacetime2026}.  The
representation decomposition and source-shorting analysis below are carried
out here on the resulting finite joint packet.

Let $\Phi$ be the twelve oriented $A_3$ roots associated with the tetrahedral
cell and let
\[
P=\R^\Phi,
\qquad
P_0=\left\{f\in P:\sum_{\alpha\in\Phi}f(\alpha)=0\right\}.
\]
Let $\mathscr T=\{c,\ell\}$ be two accepted categorical types and set
\[
t_0=2^{-1/2}(\mathbf e_c+\mathbf e_\ell),
\qquad
t_1=2^{-1/2}(\mathbf e_c-\mathbf e_\ell).
\]
After removal of the overall accepted mean, define
\begin{equation}
E_{\rm acc}^0
=
(P_0\otimes t_0)
\oplus
(\R\one_\Phi\otimes t_1)
\oplus
(P_0\otimes t_1).
\label{eq:accepted-residual}
\end{equation}
The three summands have dimensions $11,1,11$.

Write
\[
E_{\rm ST}^{\rm cat}:=P_0\otimes t_0,
\qquad
E_{\rm int}^{\rm cat}
:=(\R\one_\Phi\otimes t_1)\oplus(P_0\otimes t_1).
\]
Let
\[
S_{\rm ST}:E_{\rm ST}\to\Hh_{\rm acc},
\qquad
S_{\rm int}:E_{\rm int}^{\rm cat}\to\Hh_{\rm acc}
\]
denote respectively the complete represented spacetime source and the accepted
internal categorical source on one common carrier.  The spacetime source may
contain, in addition to $E_{\rm ST}^{\rm cat}$, independently reconstructed
waiting-time, calibration, bracket, shift, or geometric-circulation directions.
Put
\[
G_{\rm ST}=S_{\rm ST}^*S_{\rm ST},
\qquad
G_{\rm int}=S_{\rm int}^*S_{\rm int},
\qquad
C=S_{\rm int}^*S_{\rm ST},
\]
and let $P_{\rm ST}$ be the orthogonal projection onto $\Ran S_{\rm ST}$.

\begin{proposition}[Universal geometry short and active residual]
\label{thm:geometry-short}
For every source-complete common realization, shorting the accepted internal
source by the complete represented spacetime source gives
\begin{equation}
G_{{\rm int}\mid{\rm ST}}
=
G_{\rm int}-CG_{\rm ST}^{\dagger}C^*
=
S_{\rm int}^*(I-P_{\rm ST})S_{\rm int}
\succeq0.
\label{eq:geometry-short-general}
\end{equation}
The source-minimal internal categorical carrier is therefore
\begin{equation}
E_{{\rm int}\mid{\rm ST}}^{\min}
=
E_{\rm int}^{\rm cat}/\Ker G_{{\rm int}\mid{\rm ST}},
\qquad
\dim E_{{\rm int}\mid{\rm ST}}^{\min}
=
\rank G_{{\rm int}\mid{\rm ST}}.
\label{eq:geometry-short-quotient}
\end{equation}

On the explicit orthogonal active realization, the accepted contrast synthesis
has scalar Gram
\[
S_{\rm acc}^*S_{\rm acc}=\vartheta I_{E_{\rm acc}^0},
\qquad \vartheta>0,
\]
and the additional spacetime scores are orthogonal to the accepted categorical
source.  Hence the mixed block vanishes on $E_{\rm int}^{\rm cat}$ and
\begin{equation}
G_{{\rm int}\mid{\rm ST}}=\vartheta I_{12}.
\label{eq:geometry-short}
\end{equation}
On the tetrahedrally equivariant branch, the categorical residual carries the
multiplicity-free representation
\begin{equation}
E_{\rm int}^{\rm cat}\otimes\C
\cong
\C_{\rm type}^{+}
\oplus W_+
\oplus(V_2)_+
\oplus W_-
\oplus(W\otimes{\rm sgn})_-.
\label{eq:residual-S4}
\end{equation}
Here $\dim W=3$ and $\dim V_2=2$.  The odd standard block $W_-$ is the
internal endpoint module, while the odd sign-twisted block is the distinct loop
module.  If the restriction of $G_{{\rm int}\mid{\rm ST}}$ to the endpoint
block is strictly positive, that block survives the source-minimal quotient at
complex rank three.
\end{proposition}

\begin{proof}
The short is the Gram of $(I-P_{\rm ST})S_{\rm int}$, so positivity and the
quotient/rank identity follow.  The stated orthogonal realization makes the
mixed block zero and leaves $\vartheta I_{12}$.  The reversal-even oriented-root
sector is $\mathbf1\oplus W\oplus V_2$, and the reversal-odd sector is
$W\oplus(W\otimes{\rm sgn})$; Appendix~\ref{app:S4} proves this decomposition.
Equivariance makes the inequivalent endpoint and loop summands reducing for
the shorted Gram.  Strict positivity on the endpoint block therefore
preserves its three dimensions in the source quotient.
\end{proof}

The endpoint residual is defined after the complete represented spacetime
source has been removed.  It is not a second name for the spatial carrier.
Its later interpretation as a generation source will require a represented
map into an actual charged multiplicity and an allocation test.


\subsection{Compatibility with the Clifford factor}
\label{sec:clifford-tetrahedral-compatibility}

The algebra $\A_{\rm Cl}\cong M_4(\C)$ and the tetrahedral relabeling
algebra $\A_{{\rm tet},r}=C^*(\rho_r(S_4))$ are distinct represented
objects.  The latter acts on a reduced physical carrier and need not be a
Clifford factor.  Its commutant can supply internal operations for a joint
packet only after compatibility with the represented Clifford action has
been established.  The following finite criterion specifies that handoff.

\begin{proposition}[Clifford--tetrahedral compatibility and internal lift]
\label{thm:clifford-tetrahedral-compatibility}
Let $(e_{ij})_{1\le i,j\le4}$ be matrix units for a unital represented
$\A_{\rm Cl}$ on $\Hh$, and let $P_r$ be a reducing carrier projection:
$[P_r,e_{ij}]=0$.  Suppose on $P_r\Hh$ there are a genuine unitary
representation $v_r:S_4\to\U(P_r\Hh)$ and a finite historical operator
bank $(b_j)$ such that
\begin{equation}
 [v_r(g),e_{ij}|_{P_r\Hh}]=0,\qquad
 [b_j,e_{ik}|_{P_r\Hh}]=[b_j,v_r(g)]=0.
 \label{eq:clifford-tetrahedral-compatibility}
\end{equation}
Put $\mathcal V_r=e_{11}P_r\Hh$.  There is a unitary
\begin{equation}
 \mathcal U_r:\C^4\otimes\mathcal V_r\longrightarrow P_r\Hh,
 \qquad \mathcal U_r(e_i\otimes\xi)=e_{i1}\xi,
 \label{eq:clifford-reduction-unitary}
\end{equation}
under which
\begin{align}
 \A_{\rm Cl}|_{P_r\Hh}&=M_4(\C)\otimes I,\nonumber\\
 v_r(g)&=I_4\otimes\rho_r(g),\qquad
 b_j=I_4\otimes\widetilde b_j,\nonumber\\
 C^*(\widetilde b_j)&\subseteq\A_{{\rm tet},r}'.
 \label{eq:clifford-tetrahedral-lift}
\end{align}
If the reduced bank satisfies the hypotheses of
Theorem~\ref{thm:internal-seed-saturation}, its lift is a represented copy
of $M_3(\C)\oplus M_2(\C)\oplus\C\oplus\C$ in
$\A_{\rm Cl}'\cap v_r(S_4)'$.  Conversely, any such tensor-factor
realization satisfies \eqref{eq:clifford-tetrahedral-compatibility}.
\end{proposition}

\begin{proof}
The matrix-unit identities give
\[
 \langle e_{i1}\xi,e_{j1}\eta\rangle
 =\delta_{ij}\langle\xi,\eta\rangle,\qquad
 \sum_i e_{i1}e_{1i}=I.
\]
Thus \eqref{eq:clifford-reduction-unitary} is unitary.  An operator commuting
with all $e_{ij}$ is $I_4\otimes\widetilde b$ in this factorization, with
$\widetilde b$ its restriction to $e_{11}P_r\Hh$.  Apply this to $v_r(g)$
and $b_j$.  Their remaining commutators become
$[\widetilde b_j,\rho_r(g)]=0$, and the genuine group law descends to
$\rho_r$.  Relative-commutant saturation on $\mathcal V_r$ gives the
asserted internal lift.  The converse is immediate in tensor coordinates.
\end{proof}

\noindent\textit{Finite verification and nontrivial spin action.} All commutation conditions above are finite identities.  For example
$\sum_{i,k}\norm{[b,e_{ik}]}_{\HS}^2=0$ certifies Clifford commutation
exactly.  Commutation with tetrahedral relabeling alone does not imply it:
a trivial tetrahedral action commutes with every $b$, including a nonscalar
Clifford matrix unit.  If an original relabeling $u_r(g)$ normalizes rather
than commutes with the Clifford factor, a specified coherent spin
implementation $s_r(g)\in\A_{\rm Cl}$ must first be removed:
$v_r(g)=s_r(g)^*u_r(g)$.  The commuting residual must satisfy the genuine
$S_4$ group law to use the ordinary multiplicity census below.  A projective
residual instead requires its corresponding projective representation
census.  No such compatibility is inferred from a dimension count.


The word-generated internal seed is also distinct from the kinematic typing
algebra.  The seed supplies matrix-factor actions; their tensorial matter
representation and type projections generate $\A_{\rm type}$ on the matter
carrier.  Proposition~\ref{prop:typed-matter-lift} makes this handoff explicit.


\subsection{Multiplicity birth and terminating obstructions}
\label{sec:multiplicity-birth}

The external tetrahedral action is inherited from the Lorentzian relational
carrier of \cite{PelissierSpacetime2026}, while the internal gauge action
reconstructed below must be kept on the opposite side of the centralizer.
Let $\mathcal V_r$ be the future-minimal
reachable--observable residual \emph{physical carrier} spanned by represented
source words of length at most $r$, after the complete spacetime and
physical-null shorts.  The source letters are required to descend to this fixed quotient, its
word spans are nested, and its tetrahedral action preserves the spans.
Compatibility with the Clifford carrier is the separate finite condition of
Proposition~\ref{thm:clifford-tetrahedral-compatibility}.  This is not the regular word Hilbert space
$L^2(\A,\tau_{\rm reg})$: by Corollary~\ref{cor:word-multiplicity-firewall},
large invariant multiplicities of the latter do not by themselves create
physical internal factors.  Decompose the external $S_4$ action as
\begin{equation}
\mathcal V_r
\cong
\bigoplus_{\lambda\in\widehat{S_4}}
V_\lambda\otimes M_{\lambda,r},
\qquad
m_{\lambda,r}:=\dim_\C M_{\lambda,r}.
\label{eq:word-depth-isotypic}
\end{equation}
Since $S_4$ has five irreducible complex representations, we record the
\emph{complete} multiplicity census in the fixed order
\begin{equation}
\bm m_r
:=
(m_{\mathbf1,r},m_{{\rm sgn},r},m_{W,r},
  m_{W^{\rm sgn},r},m_{V_2,r}),
\qquad W^{\rm sgn}:=W\otimes{\rm sgn}.
\label{eq:complete-S4-census}
\end{equation}
No unlisted isotypic sector is understood to have zero multiplicity unless

\begin{proposition}[Complete census, colour eligibility, and flat termination]
\label{thm:active-isotypic}
At every finite physical word depth,
\begin{align}
 \A_{{\rm tet},r}&=C^*(\rho_r(S_4))
   \cong\bigoplus_\lambda M_{\dim V_\lambda}(\C)\otimes I_{M_{\lambda,r}},
 \label{eq:external-isotypic}\\
 \A_{{\rm tet},r}'&\cong\bigoplus_\lambda
   I_{V_\lambda}\otimes M_{m_{\lambda,r}}(\C).
 \label{eq:external-isotypic-commutant}
\end{align}
An internal simple factor $M_n(\C)$ on the $\lambda$-multiplicity space
requires $m_{\lambda,r}\ge n$.  The categorical degree-one census is
$(1,0,2,1,1)$, with
\begin{equation}
 \A_{{\rm tet},1}'\cong\C\oplus M_2(\C)\oplus\C\oplus\C.
 \label{eq:degree-one-commutant}
\end{equation}
One additional external-trivial private line changes the census to
$(2,0,2,1,1)$ and the commutant to
\begin{equation}
 M_2(\C)\oplus M_2(\C)\oplus\C\oplus\C.
 \label{eq:degree-one-private-commutant}
\end{equation}
Neither carrier supports a commuting internal colour $M_3(\C)$.

For the proposed assignment of the colour, weak, and scalar blocks to the
trivial, standard, sign-twisted-standard, and two-dimensional sectors,
respectively, exact exhaustion of the relative commutant requires precisely
\begin{equation}
 \bm m_r=(3,0,2,1,1).
 \label{eq:SM-complete-census}
\end{equation}
In particular, a nonzero sign-isotypic entry produces an additional summand
and cannot be omitted.

Let $K_r$ be the labelled physical source-word Gram.  If
\begin{equation}
 \rank K_r=\rank K_{r+1},\label{eq:word-flatness}
\end{equation}
then $\mathcal V_s=\mathcal V_r$ for all $s\ge r$, and every isotypic
multiplicity is frozen.  Writing
\begin{equation}
 m_{0,r}=\dim_\C\Hom_{S_4}(\mathbf1,\mathcal V_r),\qquad
 r_{\rm col}^{\rm birth}=\min\{r:m_{0,r}\ge3\},
 \label{eq:colour-birth-depth}
\end{equation}
with the value $\infty$ when no birth occurs before flatness, a flat depth
with $m_{0,r}\le2$ permanently excludes an external-trivial internal colour
factor.  More generally, failure of $m_{\lambda,r}\ge n$ at a flat depth is
terminal for an internal $M_n(\C)$ on that multiplicity space.
\end{proposition}
\begin{proof}
The isotypic decomposition, Burnside's theorem on each irreducible external
factor, and Schur's lemma give the two algebras.  A nonzero representation of
$M_n(\C)$ has dimension at least $n$, giving the multiplicity obstruction.
The degree-one representation is
$\C_{\rm type}\oplus(W\otimes\C^2)\oplus V_2\oplus W^{\rm sgn}$,
which gives the two displayed censuses.  Independent isotypic blocks give
\eqref{eq:SM-complete-census} for the specified assignment.
Finally, $\mathcal V_{r+1}$ is spanned by $\mathcal V_r$ and its one-letter
images.  Equality of Gram ranks makes those spaces equal, so every source
letter preserves $\mathcal V_r$.  Induction proves stabilization and the
terminal multiplicity obstructions.
\end{proof}

Multiplicity three only creates a colour-eligible carrier.  Full matrix
operations appear only when the represented word algebra connects it.
The next tests separate this algebra-generation question from occurrence.


\subsection{Internal generation and relative-commutant saturation}
\label{sec:internal-landing}

\begin{proposition}[Colour bridge, weak router, and central separation]
\label{thm:colour-bridge}
On an external-trivial three-space $T\oplus H_{\rm priv}$, with
$\dim T=1$ and $\dim H_{\rm priv}=2$, assume the represented private support
words generate $M_2(H_{\rm priv})$.  For an actual word $u$, put
\begin{equation}
 \omega_{\rm col}(u)=\|p_Hup_T\|_{\HS}^2+\|p_Tup_H\|_{\HS}^2.
 \label{eq:colour-bridge}
\end{equation}
Then
\begin{equation}
 C^*(\C p_T\oplus M_2(H_{\rm priv}),u)=M_3(\C)
 \quad\Longleftrightarrow\quad\omega_{\rm col}(u)>0.
 \label{eq:colour-M3}
\end{equation}
At zero bridge the algebra is exactly $\C\oplus M_2(\C)$.
On the standard-isotypic multiplicity plane, let $p_{\rm t},p_{\rm h}$
project onto unit vectors $t,h$.  Then
\begin{equation}
 0<|\langle t,h\rangle|<1
 \quad\Longleftrightarrow\quad C^*(p_{\rm t},p_{\rm h})=M_2(\C).
 \label{eq:weak-M2}
\end{equation}
A single rank-one projection generates only $\C^2$.

Finally, a unital subalgebra of
$M_3(\C)\oplus M_2(\C)\oplus\C\oplus\C$ surjecting onto each factor is
either the full $15$-dimensional product or the $14$-dimensional algebra
\begin{equation}
 \{(A,B,z,z):A\in M_3(\C),\ B\in M_2(\C),\ z\in\C\}.
 \label{eq:central-locked-14}
\end{equation}
For a generating bank $(b_j)$ and the two scalar characters $\chi_1,\chi_2$,
the full product occurs exactly when
\begin{equation}
 \eta_{\rm cen}=\sum_j|\chi_1(b_j)-\chi_2(b_j)|^2>0.
 \label{eq:central-separation-margin}
\end{equation}
Independent represented scalar central supports suffice.
\end{proposition}
\begin{proof}
A nonzero line--plane corner, its adjoint, and the private matrix units
generate all rectangular matrix units, hence $M_3$; absent both corners the
two blocks reduce the algebra.  Two distinct nonorthogonal rank-one
projections in dimension two have no common reducing line, so Burnside's
theorem gives $M_2$.  Finally, a subdirect product of simple matrix algebras
can identify only isomorphic factors
(Proposition~\ref{thm:subdirect-classification}).  Only the two scalar factors
can therefore be locked, and a difference of their characters on a generator
is exactly the condition excluding that identification.
\end{proof}

A bridge can itself be a composite historical word.
Appendix~\ref{app:reciprocal-return} proves the finite reciprocal-return
criterion: nonorthogonality of the two mediator cyclic multiplicity spaces
is equivalent to a positive return mass, while zero mass gives a reducing
projection excluding every longer excursion in the declared route bank.
This alternative supplies a bridge rather than assuming an extra direct one.

Define $\A_{{\rm int},r}^{\rm word}$ to be the $C^*$-algebra generated by
all represented physical word compressions at depth $r$ that commute with the
external tetrahedral action.

\begin{theorem}[Internal-algebra reconstruction by relative-commutant saturation]
\label{thm:internal-seed-saturation}
Suppose that at one finite word depth the complete external multiplicity
census is
\begin{equation}
\bm m_r=(3,0,2,1,1)
\label{eq:SM-multiplicity-census}
\end{equation}
in the ordering \eqref{eq:complete-S4-census},
that the external-trivial three-space is coherently realized as
$T\oplus H_{\rm priv}$, and that:
\begin{enumerate}[label=\textup{(I\arabic*)},leftmargin=3em]
\item the private support words generate $M_2(H_{\rm priv})$ and either one
      represented word has $\omega_{\rm col}>0$, or an admitted reciprocal
      mediator bank satisfies $\mathfrak m_{\rm ret}>0$ in
      Proposition~\ref{thm:reciprocal-return};
\item the standard multiplicity plane contains two distinct nonorthogonal
      represented rank-one projections;
\item the $V_2$ and $W^{\rm sgn}$ external blocks have multiplicity one;
\item the two scalar sectors are centrally separated, equivalently the
      generated algebra contains their independent central supports or passes
      $\eta_{\rm cen}>0$ in Proposition~\ref{thm:colour-bridge}.
\end{enumerate}
Then
\begin{equation}
\A_{{\rm int},r}^{\rm word}
=
\A_{{\rm tet},r}'
\cong
M_3(\C)\oplus M_2(\C)\oplus\C\oplus\C.
\label{eq:active-internal-algebra}
\end{equation}
Its complex dimension is $15$, its centre has dimension four, and the two
nonabelian simple factors are intrinsically distinguished by their matrix
sizes.  Hence the labels ``colour'' and ``weak'' are canonical up to inner
automorphism once the multiplicity packet has landed.
\end{theorem}

\begin{proof}
Condition~\textup{(I1)} and Proposition~\ref{thm:colour-bridge} generate the
$M_3$ block; when the reciprocal route alternative is used,
Proposition~\ref{thm:reciprocal-return} supplies the required represented bridge.
Condition~\textup{(I2)} and Proposition~\ref{thm:colour-bridge} generate the $M_2$ block.
The multiplicity-one sectors give surjective scalar compressions, and
condition~\textup{(I4)} excludes their diagonal identification by
Proposition~\ref{thm:colour-bridge}.  The zero sign entry in
\eqref{eq:SM-multiplicity-census}, together with
Proposition~\ref{thm:active-isotypic}, excludes every additional
sign-isotypic commutant block.  Thus the generated internal algebra is the
full product and exhausts the complete relative commutant.
\end{proof}

The saturation theorem is conditional on the complete census and on the
occurrence of the stated connecting words.  At that census it also gives
an explicit classification of the smaller branches rather than a completion
of missing structure.

\begin{proposition}[Exact internal-algebra deficit alternatives]
\label{thm:internal-deficit-alternatives}
On a coherent carrier with complete census $\bm m_r=(3,0,2,1,1)$, assume the colour
three-space initially carries $M_2(\C)\oplus\C$, the weak multiplicity plane
carries one rank-one projection, and the displayed colour, weak and two scalar
central supports occur independently in the represented internal algebra.
Then exactly the following dimensions occur:
\begin{center}
\begin{tabular}{L{0.47\textwidth}c}
\toprule
represented internal words & $\dim_\C\A_{\rm int}^{\rm word}$\\\midrule
neither colour bridge nor weak router & $9$\\
weak router but no colour bridge & $11$\\
colour bridge but no weak router & $13$\\
both bridge and router & $15$\\
\bottomrule
\end{tabular}
\end{center}
The terminal value $15$ is equivalent to saturation of the complete relative
commutant in Theorem~\ref{thm:internal-seed-saturation}.  Each failure is
therefore an exact finite source-word rank deficit rather than a discarded
branch.  If independent central supports have not yet been established,
Proposition~\ref{thm:colour-bridge} must be applied before componentwise
dimensions are added; in particular, the terminal blockwise-surjective branch
may still be the $14$-dimensional scalar-locked algebra
\eqref{eq:central-locked-14}.
\end{proposition}

\begin{proof}
Without the bridge the colour block is $M_2\oplus\C$, of dimension five; with
the bridge it is $M_3$, of dimension nine.  One weak rank-one projection
generates $\C^2$, of dimension two; a noncommuting second projection generates
$M_2$, of dimension four.  The two multiplicity-one sectors add two scalar
dimensions, giving $9,11,13,15$.
\end{proof}

For fixed multiplicity carriers and diagonal minimal projections,
nonabelian router saturation is open and dense in the affine router class;
its failure is a finite union of disconnected-support linear subspaces
(Proposition~\ref{prop:router-failure-locus}).  This genericity is relative to
that admitted class and asserts no microscopic occurrence law.

The algebra just reconstructed is the complex relative commutant of the
reduced tetrahedral action, lifted through Clifford compatibility.  It is
neither the commutant of the complete matter-incidence action nor the real
finite coordinate algebra of an almost-commutative spectral triple.
The next section determines the physical gauge group from represented
incidence or anomaly data, and then its actual matter typing algebra.


\section{Gauge group and matter typing}
\label{sec:gauge-matter}

Landing the internal matrix factors does not yet fix their relative tensor
position, the central gauge subgroup, or the matter inventory.  We treat
these questions in that order.  The determinant-source and anomaly routes
to the global group are logically independent; the charged-copy count then
requires additional source-completeness and endpoint-allocation hypotheses.


\subsection{Tensor alignment and finite naturality}
\label{sec:determinant-naturality}

The global form of the Standard-Model gauge group is not fixed by its Lie
algebra alone \cite{Hucks1991,Tong2017}.  We first identify a determinant
incidence by finite equations on the landed matrix factors, then distinguish
its operator ray from the phase-sensitive datum that selects a subgroup.
Throughout, $C\cong\C^3$ and $W_2\cong\C^2$ are the actual internal
multiplicity carriers after Theorem~\ref{thm:internal-seed-saturation}.
Their actions on the indicated source and target legs must be represented on
the same identified typed carrier.  They are not external tetrahedral
Burnside blocks.

The relative orientation of these factor actions is additional represented
data.  Let $\mathcal K_{\det}$ be the carrier supporting the colour input and
the two retained weak input slots, and let
$\mathcal A_C\cong M_3(\C)$ and
$\mathcal A_{W,1},\mathcal A_{W,2}\cong M_2(\C)$ be their represented unital
factor algebras on $\mathcal K_{\det}$.  Choose Hermitian spanning families
$(C_\alpha)$, $(W_\beta^{(1)})$, and $(W_\gamma^{(2)})$ and put
\begin{equation}
\begin{aligned}
\Delta_{\otimes}:={}&
\sum_{\alpha,\beta}\norm{[C_\alpha,W_\beta^{(1)}]}_{\HS}^2
+\sum_{\alpha,\gamma}\norm{[C_\alpha,W_\gamma^{(2)}]}_{\HS}^2\\
&+\sum_{\beta,\gamma}\norm{[W_\beta^{(1)},W_\gamma^{(2)}]}_{\HS}^2.
\end{aligned}
\label{eq:tensor-alignment-defect}
\end{equation}

\begin{proposition}[Same-carrier tensor-factor alignment]
\label{thm:tensor-factor-alignment}
The condition
\begin{equation}
\Delta_{\otimes}=0
\label{eq:tensor-alignment-zero}
\end{equation}
holds if and only if the three represented factor algebras commute pairwise.
On this branch multiplication induces an injective unital $*$-homomorphism
\begin{equation}
M_3(\C)\otimes M_2(\C)\otimes M_2(\C)
\hookrightarrow\mathcal B(\mathcal K_{\det}),
\label{eq:tensor-factor-product}
\end{equation}
and the generated algebra is unitarily equivalent to
$M_{12}(\C)\otimes I_M$ for a finite multiplicity space $M$.  Thus the
intrinsic tensor action $C\otimes W_2\otimes W_2$ is represented on the same
carrier, while any residual multiplicity remains a spectator for this tensor
factorization.
\end{proposition}

\begin{proof}
Vanishing of \eqref{eq:tensor-alignment-defect} is equivalent to commutation of
all elements of the chosen spanning families and hence, by linearity,
adjoints and multiplication, to pairwise commutation of the generated factor
algebras.  Multiplication then defines a unital $*$-homomorphism from their
finite tensor product, which is the simple algebra $M_{12}(\C)$.  A unital
homomorphism from a simple full matrix algebra has zero kernel, so the product
representation is faithful.  The finite-dimensional representation theorem
then gives the displayed spatial form.
\end{proof}

A sufficient reconstruction of these commuting tensor actions from mixed
pair and triple source occurrences is proved in
Proposition~\ref{thm:fusion-native-alignment}.  Separate factor marginals are
not sufficient (Proposition~\ref{cth:factor-marginals-no-alignment}).  Once
alignment is established, the matrix units below are coordinates of already
represented actions, not additional physical controls.

Set
\begin{equation}
 \mathscr U=\Hom(C\otimes W_2\otimes W_2,\Lambda^2C^*),
 \qquad \dim_\C\mathscr U=36.
 \label{eq:naturality-port}
\end{equation}

Source and target typing alone do not determine an incidence line:
Appendix~\ref{app:naturality} computes the normalization kernel of the
unrestricted rectangular envelope.  The required additional information is
supplied by finite identities for the represented colour and weak actions.

For $X\in\End(C)$ and $Y\in\End(W_2)$ define
\begin{align}
 d\rho_C^{\rm in}(X)&=X\otimes I\otimes I,\nonumber\\
 d\rho_C^{\rm out}(X)(\varphi\wedge\psi)
   &=-(\varphi\circ X)\wedge\psi-\varphi\wedge(\psi\circ X),\nonumber\\
 d\rho_W^{\rm in}(Y)&=I_C\otimes Y\otimes I+I_C\otimes I\otimes Y,
 \label{eq:naturality-actions}
\end{align}
and the linear defect maps on $\mathscr U$
\begin{equation}
 \mathcal R_X(T)=d\rho_C^{\rm out}(X)T-Td\rho_C^{\rm in}(X),
 \qquad
 \mathcal S_Y(T)=-T d\rho_W^{\rm in}(Y).
 \label{eq:naturality-defects}
\end{equation}
Choose matrix-unit systems for the represented factors and put
\begin{equation}
 \mathcal G_C=\{E_{12},E_{21},E_{23},E_{32}\},\qquad
 \mathcal G_W=\{F_{12},F_{21}\}.
 \label{eq:six-naturality-generators}
\end{equation}
No central generator or tensor-slot interchange is included in this bank.

\begin{theorem}[Determinant-line reconstruction and positive naturality certificate]
\label{thm:six-naturality}
The common kernel
\begin{equation}
 \mathsf S_{\rm nat}=\bigcap_{X\in\mathcal G_C}\Ker\mathcal R_X
                  \cap\bigcap_{Y\in\mathcal G_W}\Ker\mathcal S_Y
 \label{eq:naturality-kernel}
\end{equation}
is one-dimensional and canonically isomorphic to
$(\det C\otimes\det W_2)^*$.  Its nonzero elements are
$\SU(C)\times\SU(W_2)$-equivariant, vanish on
$C\otimes\operatorname{Sym}^2W_2$, and obey
\begin{equation}
 T(I_C\otimes\tau_W)=-T,\qquad
 T\longmapsto e^{-i(3\alpha+2\beta)}T
 \label{eq:naturality-swap-character}
\end{equation}
under weak-slot interchange and central phases.  The colour equations alone
have kernel dimension four, and the weak equations alone dimension nine.

Choose $T_0$ on this line with $T_0^*T_0=P_{\rm alt}$, where $P_{\rm alt}$
projects onto $C\otimes\Lambda^2W_2$.  Then $\|T_0\|_{\HS}^2=3$ and
$P_{\det}=|T_0\rangle\!\rangle\langle\!\langle T_0|/3$ is its orthogonal
line projector.  For standard normalized matrix units and orthonormal wedge
bases, the positive defect operator
\begin{equation}
 H_{\rm nat}=\sum_{X\in\mathcal G_C}\mathcal R_X^*\mathcal R_X
           +\sum_{Y\in\mathcal G_W}\mathcal S_Y^*\mathcal S_Y
 \label{eq:naturality-operator}
\end{equation}
satisfies
\begin{equation}
 \Ker H_{\rm nat}=\mathsf S_{\rm nat},\qquad
 H_{\rm nat}\succeq2(I-P_{\det}),\qquad
 \lambda_{\min}^+(H_{\rm nat})=2.
 \label{eq:naturality-gap}
\end{equation}
For any positive Choi operator $J_D$ on this port, without a Kraus-rank
assumption, put $d_{\rm nat}=\Tr(H_{\rm nat}J_D)$.  Then
\begin{equation}
 \Tr[(I-P_{\det})J_D]\le\tfrac12d_{\rm nat},
 \label{eq:naturality-leakage}
\end{equation}
and
\begin{equation}
 d_{\rm nat}=0\quad\Longleftrightarrow\quad
 J_D=x|T_0\rangle\!\rangle\langle\!\langle T_0|,\qquad x\ge0.
 \label{eq:naturality-ray}
\end{equation}
The additional condition $\Tr J_D>0$ is equivalent to $x>0$.
The spectrum, kernel, and simultaneously transported certificates are
unchanged by unitary changes of matrix-unit systems.
\end{theorem}
\begin{proof}
The four colour units generate $\mathfrak{sl}_3(\C)$ and the two weak units
$\mathfrak{sl}_2(\C)$ under brackets and linear combinations.  Intertwining
the generators therefore intertwines these Lie algebras.
The weak decomposition
$W_2\otimes W_2=\operatorname{Sym}^2W_2\oplus\Lambda^2W_2$
has only one trivial summand, the alternating line.  Since the target is
weak-trivial, the weak identities leave $\Hom(C,\Lambda^2C^*)$, of dimension
nine.  The volume identification makes its two colour modules irreducible
and isomorphic for $\mathfrak{sl}_3$, so Schur's lemma leaves one line.
Colour identities alone leave four weak covector coefficients.
Connectedness exponentiates the identities to the compact special-unitary
groups.  The target weight $-2\alpha$ and inverse source weight
$-\alpha-2\beta$ give \eqref{eq:naturality-swap-character}.

The canonical line identification is given by
\[
 (\Theta_\tau(c\otimes w_1\otimes w_2))(c_1,c_2)
 =\tau(c\wedge c_1\wedge c_2\otimes w_1\wedge w_2).
\]
It is nonzero and hence an isomorphism between the two lines.
The sum-of-squares definition gives the kernel of $H_{\rm nat}$;
Appendix~\ref{app:naturality-spectrum} diagonalizes its colour and weak
summands and proves the sharp gap.  Tracing this bound against $J_D\succeq0$
gives \eqref{eq:naturality-leakage}.  Finally,
$\Tr(H_{\rm nat}J_D)=\|H_{\rm nat}^{1/2}J_D^{1/2}\|_{\HS}^2$ vanishes
exactly when $J_D$ is supported on the one-dimensional kernel.  Positivity
gives the stated ray and its trace $3x$.  Unitary changes of coordinates
conjugate the representations and the defect operator simultaneously.
\end{proof}

The factor actions determine the selector before a source populates its
port.  Zero defect tests purity, whereas positive mass tests occurrence;
small positive defect bounds leakage but does not imply exact selection.
For positive branch refinements, zero total defect forces zero defect on
each branch.  Proposition~\ref{thm:determinant-skeleton-population} gives the
two-sided spectral bounds and branchwise occurrence estimates.


\subsection{The global group: determinant and anomaly routes}

\begin{definition}[Phase-sensitive determinant seed]
A determinant seed is a fixed nonzero functional, normalized to a unitary
isomorphism,
\begin{equation}
 \tau:\det C\otimes\det W_2\xrightarrow{\sim}\C,
 \label{eq:determinant-seed}
\end{equation}
where the target scalar line is fixed pointwise.  Such a tensor can be
supplied directly as a represented incidence or recovered relative to a
marked coherent neutral reference.  Changing its overall normalization
phase does not change its stabilizer.  Retaining only the operator ray or
its Choi projector is a weaker datum and is not a phase-sensitive seed.
\end{definition}

\begin{theorem}[Gauge-group reconstruction and the phase-sensitive criterion]
\label{thm:gauge-group}
The unitary stabilizer of a phase-sensitive determinant seed is
\begin{equation}
 \Aut_{\U}(C,W_2,\tau)=S(\U(3)\times\U(2))
 \cong\frac{\SU(3)\times\SU(2)\times\U(1)}{\Z_6}.
 \label{eq:SM-gauge-group}
\end{equation}
Equivalently, let $\chi(A,B)=\det A\det B$ and let the group act by $1$ on
a marked neutral amplitude space $\mathscr V_0$ and by $\chi^{-1}$ on
$\mathsf S_{\rm nat}$.  For a positive packet
\[
 J=\begin{pmatrix}J_{00}&J_{0D}\\J_{0D}^*&J_{DD}\end{pmatrix},\qquad
 h_{\det}=\|J_{0D}\|_{\HS}^2,
\]
one has
\begin{equation}
 h_{\det}>0\quad\Longleftrightarrow\quad
 \operatorname{Stab}_{\U(3)\times\U(2)}(J)=\ker\chi.
 \label{eq:operational-determinant-group}
\end{equation}
If $h_{\det}=0$, this marked packet has the full stabilizer
$\U(3)\times\U(2)$, even when $J_{DD}$ is a nonzero exact naturality ray.
For $\chi(g)=e^{i\phi}$,
\begin{equation}
 \|\mathcal R(g)J\mathcal R(g)^*-J\|_{\HS}^2
       =8h_{\det}\sin^2(\phi/2).
 \label{eq:determinant-phase-rigidity}
\end{equation}
A larger marked packet has this full stabilizer only when its other retained
components are $\ker\chi$-invariant; otherwise their stabilizers must be
intersected.
\end{theorem}
\begin{proof}
A pair $(A,B)$ preserves $\tau$ exactly when $\det A\det B=1$.
The homomorphism
$\Phi(g_3,g_2,z)=(z^{-2}g_3,z^3g_2)$ maps onto this subgroup: for $(A,B)$
in it, choose $z^6=\det B$, and put $g_3=z^2A$, $g_2=z^{-3}B$.
Its kernel is
$\{(z^2I_3,z^{-3}I_2,z):z^6=1\}\cong\Z_6$.
On the marked amplitude packet, the diagonal blocks are invariant and
conjugation multiplies $J_{0D}$ by $\chi$.  A nonzero cross block is fixed
exactly when $\chi=1$.  The two orthogonal cross blocks contribute
$2|e^{i\phi}-1|^2\|J_{0D}\|_{\HS}^2$, proving the quantitative formula.
Additional marked data restrict a stabilizer by intersection.
\end{proof}

A rank-one Choi packet is phase blind.  Neither its normalization nor its
positive mass creates the coherent neutral--incidence cross block.
The normalization and minimal-arity statements are retained in
Appendix~\ref{app:anomaly}.  We next obtain the same group by a different
route, using a specified chiral tensor packet and anomaly freedom rather
than determinant-source occurrence.

Let the infinitesimal central action on $C\oplus W_2$ be
\[
iaI_C\oplus ibI_{W_2}.
\]
Define the tensor types
\begin{equation}
Q=C\otimes W_2,\qquad
u=\Lambda^2C^*,\qquad
d=C,\qquad
L=W_2^*,\qquad
e=\Lambda^2W_2^*,\qquad
H=W_2.
\label{eq:SM-types}
\end{equation}
Their central weights are
\begin{equation}
(a+b,-2a,a,-b,-2b,b).
\label{eq:central-weights}
\end{equation}

\begin{proposition}[Anomaly-forced central weights]
\label{thm:hypercharge}
Write right-handed fermions as left-handed conjugates when computing chiral
anomalies.  The local
$\SU(3)^2\U(1)$,
$\SU(2)^2\U(1)$,
gravitational--$\U(1)$,
and cubic $\U(1)$ anomalies vanish simultaneously if and only if
\begin{equation}
3a+2b=0.
\label{eq:anomaly-condition}
\end{equation}
Up to overall reversal of the $\U(1)$ orientation, the primitive integral
solution is
\[
(a,b)=(-2,3).
\]
With $y=6Y$, the resulting physical weights are
\begin{equation}
y(Q,u,d,L,e,H)=(1,4,-2,-3,-6,3).
\label{eq:hypercharges}
\end{equation}
If a neutral right-handed singlet is retained, $y(\nu)=0$.
\end{proposition}

\begin{proof}
After charge conjugation of the right-handed fields,
\begin{align*}
\mathcal A_{\SU(3)^2Y}
&=
\mathcal A_{\SU(2)^2Y}
=
\frac12(3a+2b),\\
\mathcal A_{{\rm grav}^2Y}
&=
3(3a+2b),\\
\mathcal A_{Y^3}
&=
3(3a+2b)(3a^2+2b^2).
\end{align*}
The pure colour anomaly cancels as $2-1-1$, and the weak doublet count is even.
For a nontrivial abelian action $3a^2+2b^2>0$, so simultaneous local anomaly
cancellation is equivalent to \eqref{eq:anomaly-condition}.  The primitive
coprime integral solution is $(-2,3)$ up to sign; substitution gives
\eqref{eq:hypercharges}.
\end{proof}

For the global quotient one must also determine the common kernel of the
represented matter action.  We display it in left-handed convention so
that the group kernel and the chiral anomaly convention can be compared
directly.

\begin{proposition}[Faithful matter realization of the determinant quotient]
\label{prop:faithful-SM-quotient}
Normalize the abelian charge by $y=6Y$ and write all fermions as left-handed
Weyl fields.  On one generation the represented matter packet is
\[
Q:(\mathbf3,\mathbf2)_1,
\qquad
u^c:(\overline{\mathbf3},\mathbf1)_{-4},
\qquad
d^c:(\overline{\mathbf3},\mathbf1)_2,
\]
\[
L:(\mathbf1,\mathbf2)_{-3},
\qquad
e^c:(\mathbf1,\mathbf1)_6,
\qquad
H:(\mathbf1,\mathbf2)_3,
\qquad
[\nu^c:(\mathbf1,\mathbf1)_0].
\]
The common kernel of these representations of
$\widetilde G=\SU(3)\times\SU(2)\times\U(1)_y$ is exactly
\begin{equation}
K
=
\left\langle
\left(e^{2\pi i/3}I_3,-I_2,e^{i\pi/3}\right)
\right\rangle
\cong\Z_6.
\label{eq:faithful-SM-kernel}
\end{equation}
Consequently the represented matter packet is faithful precisely for
$\widetilde G/K$, which is the same global group selected intrinsically by the
determinant seed in Theorem~\ref{thm:gauge-group}.
\end{proposition}

\begin{proof}
A common-kernel element must be central, so write it as
$(\omega_3^aI_3,(-1)^bI_2,e^{i\theta})$.  The charged singlet $e^c$ gives
$\theta=m\pi/3$ modulo $2\pi$; the lepton doublet gives
$b\equiv m\pmod2$; and the colour antitriplet $d^c$ gives
$a\equiv m\pmod3$.  Substitution in the remaining representations shows that
precisely the six powers of the generator in
\eqref{eq:faithful-SM-kernel} act trivially.  The optional neutral singlet is
trivial under all factors and adds no further restriction.
\end{proof}

\begin{corollary}[Faithful global group from anomaly freedom]
\label{cor:anomaly-faithful-group}
For the represented chiral tensor packet, primitive local anomaly freedom fixes
$(a,b)=(-2,3)$ up to simultaneous sign.  Together with the faithful matter
kernel of Proposition~\ref{prop:faithful-SM-quotient}, this gives
\[
 G_{\rm phys}\cong
 S(\U(3)\times\U(2))
 \cong
 (\SU(3)\times\SU(2)\times\U(1))/\Z_6.
\]
This reconstruction does not assume occurrence or phase coherence of a
historical determinant-incidence source.  When such a source is represented,
Theorem~\ref{thm:gauge-group} and
Theorem~\ref{thm:gauge-group} give the independent
fixed-tensor and marked-source stabilizer statements.
\end{corollary}

\begin{proof}
Proposition~\ref{thm:hypercharge} fixes the primitive embedding
$(a,b)=(-2,3)$ up to simultaneous sign, and
Proposition~\ref{prop:faithful-SM-quotient} gives its matter kernel $\Z_6$.
The homomorphism $\Phi$ in the proof of Theorem~\ref{thm:gauge-group}
identifies the quotient with $S(\U(3)\times\U(2))$ as a group-theoretic
identity; using that identity requires no determinant-source occurrence.
\end{proof}

\begin{proposition}[Kinematic typing algebra on the represented matter carrier]
\label{prop:typed-matter-lift}
Suppose the fermion modules of Proposition~\ref{prop:faithful-SM-quotient},
with multiplicity spaces $N_f$, occur on one Clifford-compatible carrier
\[
 \Hh_{\rm mat}\cong\C^4\otimes
       \bigoplus_{f\in\mathcal F}V_f\otimes N_f,
 \qquad \mathcal F=\{Q,u^c,d^c,L,e^c\}\quad[\text{and }\nu^c].
\]
Let $p_f$ be the actual type projections and $r(g)$ the represented
$G_{\det}$ action.  Then
\begin{equation}
 \A_{\rm type}:=C^*(r(g),p_f:g\in G_{\det},f\in\mathcal F)
 =\bigoplus_f I_4\otimes\mathcal B(V_f)\otimes I_{N_f}.
 \label{eq:typed-matter-lift}
\end{equation}
It commutes with $\A_{\rm Cl}$ and supplies the kinematic algebra in
Definition~\ref{def:joint-packet}.  Without the neutral singlet, its
abstract block sizes are $(6,3,3,2,1)$, rather than the seed sizes
$(3,2,1,1)$.  Any represented incidence family on these same sectors,
including non-factorizing incidences, then gives the joint packet of
Theorem~\ref{thm:main-duality}.
\end{proposition}

\begin{proof}
Every displayed matter representation is irreducible.  On each type block,
Schur's lemma gives the scalar commutant on $V_f$, and finite bicommutant
closure gives $\mathcal B(V_f)$.  The type projections separate the blocks,
while the represented multiplicity factors are untouched.  Clifford
compatibility makes the spinor factor the identity for this algebra.  The
stated block dimensions are those of the tensor representations.  The
incidence conclusion is precisely the unrestricted coefficient-space
construction; the hypothesis that these types occur on the carrier is not
inferred from the abstract seed algebra alone.
\end{proof}

Thus the seed algebra, its physical gauge group, and the matter typing
algebra have different roles.  In particular the two scalar seed summands
are not two additional physical $\U(1)$ groups, and the matter block sizes
are $(6,3,3,2,1)$ rather than $(3,2,1,1)$.  With this represented action
fixed, we can now ask which charged sectors are actually sourced or reached.


\section{Charged inventory and generations}
\label{sec:charged-matter}

The preceding section fixes the represented gauge action and its matter
types.  It does not yet determine the reached charged inventory or the
number of copies.  We first remove unsourced ambient directions, then
combine the chiral inventory and linear anomaly tests with a faithful
endpoint occurrence and its positive allocation residual.


\subsection{Source-minimal charged inventory}
\label{sec:source-minimal-charged-inventory}

All counts below refer to the complete retained physical fermion source and
its finite cyclic closure, not to an arbitrary ambient completion.  The
external Clifford factor is suppressed, so its dimension is not included in
the charged multiplicities.  Completeness is a premise of the source packet;
when an independently complete atlas is available, the positive residual of
Proposition~\ref{prop:fermion-source-completeness} tests it explicitly.

Let $\Hh_F^{\rm amb}$ be a finite represented carrier and let
\begin{equation}
 \mathcal A_F=C^*(c_1,\ldots,c_m)\subseteq\mathcal B(\Hh_F^{\rm amb})
 \label{eq:fermion-loaded-algebra}
\end{equation}
be the $C^*$-algebra generated by the represented gauge action, grading,
full-cell generator, matter incidences, and every other physical letter used by
the structural claim.  Let
$B_F:E_F\to\Hh_F^{\rm amb}$ synthesize the complete retained physical fermion
source after the declared null short.  Define
\begin{equation}
 \mathcal K_0=\Ran B_F,\qquad
 \mathcal K_{r+1}=\mathcal K_r+
 \sum_{j=1}^m(c_j\mathcal K_r+c_j^*\mathcal K_r).
 \label{eq:fermion-cyclic-filtration}
\end{equation}

\begin{proposition}[Terminating source-minimal fermion saturation]
\label{thm:fermion-source-saturation}
The filtration \eqref{eq:fermion-cyclic-filtration} stabilizes after finitely
many strict increases at
\begin{equation}
 \mathcal K_F^{\min}=\mathcal A_F\Ran B_F.
 \label{eq:fermion-minimal-carrier}
\end{equation}
At the first equality $\mathcal K_{r+1}=\mathcal K_r$, the stabilized space
reduces $\mathcal A_F$.  If $S_r$ synthesizes $\mathcal K_r$ and $N_r$
synthesizes all columns $c_jS_r,c_j^*S_r$, write
\[
 \begin{pmatrix}G_r&C_r\\C_r^*&D_r\end{pmatrix}
 =\begin{pmatrix}S_r^*\\N_r^*\end{pmatrix}
  \begin{pmatrix}S_r&N_r\end{pmatrix}
\]
and put
\begin{equation}
 \mathbb L_r=D_r-C_r^*G_r^\dagger C_r
 =N_r^*(I-P_r)N_r\succeq0.
 \label{eq:fermion-cyclic-innovation}
\end{equation}
Then
\begin{equation}
 \rank\mathbb L_r
 =\dim\mathcal K_{r+1}-\dim\mathcal K_r,
 \label{eq:fermion-cyclic-rank}
\end{equation}
so $\mathbb L_r=0$ is an exact terminating saturation test.  Compression to
$\mathcal K_F^{\min}$ preserves every loaded source-word Gram, while its
orthogonal complement is invisible to all such source correlations.
\end{proposition}

\begin{proof}
The spaces are nested and every strict inclusion increases finite dimension.
At the first equality all $c_j$ and $c_j^*$ preserve the stabilized space,
which is therefore reducing and is the smallest reducing subspace containing
$\Ran B_F$.  The Schur complement in
\eqref{eq:fermion-cyclic-innovation} is the Gram of the component of the next
columns orthogonal to $\mathcal K_r$; its rank is therefore exactly the new
cyclic dimension.  Every loaded word applied to $\Ran B_F$ lies in the cyclic
space, proving the final assertion.
\end{proof}

Fix on $\mathcal K_F^{\min}$ the represented gauge action, its self-adjoint
central generator $Y$ normalized by $y=6Y$, and a fermionic chirality grading
$\Gamma_F=\Gamma_F^*=\Gamma_F^{-1}$ commuting with the gauge action.  Choose a
finite family of represented nonabelian infinitesimal generators
$X_1,\ldots,X_s$ which, together with $Y$, generate the finite gauge-action
algebra.  For a target irreducible gauge type
$f\in\{Q,u,d,L,e\}$ from \eqref{eq:SM-types}, with central weight $y_f$ and
physical chirality $\epsilon_f\in\{\pm1\}$, define on
$\Hom(V_f,\mathcal K_F^{\min})$
\begin{equation}
 \mathbb T_{f,\epsilon}
 =\sum_{j=1}^s\mathcal R_{f,j}^*\mathcal R_{f,j}
  +\mathcal Y_f^*\mathcal Y_f
  +\mathcal G_{f,\epsilon}^*\mathcal G_{f,\epsilon},
 \label{eq:canonical-type-tester}
\end{equation}
where
\[
 \mathcal R_{f,j}(T)=X_jT-T\rho_f(X_j),\qquad
 \mathcal Y_f(T)=YT-y_fT,\qquad
 \mathcal G_{f,\epsilon}(T)=\Gamma_FT-\epsilon T.
\]
Let $\mathcal I_{f,\epsilon}=\Ker\mathbb T_{f,\epsilon}$ and equip it with
$\langle S,T\rangle_{\rm mult}=d_f^{-1}\Tr(S^*T)$,
$d_f=\dim V_f$.

\begin{proposition}[Canonical chiral projectors and the charged inventory]
\label{thm:canonical-chiral-projectors}
For an orthonormal basis $T_1,\ldots,T_n$ of $\mathcal I_{f,\epsilon}$,
\begin{equation}
 T_a^*T_b=\delta_{ab}I_{V_f},\qquad
 P_{f,\epsilon}=\sum_aT_aT_a^*,\qquad
 n_{f,\epsilon}=\dim\mathcal I_{f,\epsilon}
               =\frac{\rank P_{f,\epsilon}}{d_f}.
 \label{eq:canonical-chiral-projector}
\end{equation}
The projection is basis independent and its range is the complete physical
$(f,\epsilon)$-isotypic subspace.  In particular the evaluation map
$\operatorname{ev}_f:V_f\otimes\mathcal I_{f,\epsilon_f}
\to P_{f,\epsilon_f}\mathcal K_F^{\min}$, $v\otimes T\mapsto Tv$, is unitary.

Let $P_0^F$ project onto the fully gauge-trivial subspace, set
$P_{\rm gch}^F=P_F-P_0^F$, and define
\begin{align}
 P_{\rm SM}^F&=\sum_fP_{f,\epsilon_f},\qquad
 P_{\rm mir}^F=\sum_fP_{f,-\epsilon_f},\nonumber\\
 P_{\rm ex}^F&=P_{\rm gch}^F-\sum_f(P_{f,+}+P_{f,-}).
 \label{eq:future-visible-inventory-projectors}
\end{align}
These are orthogonal projections with
\begin{equation}
 P_{\rm gch}^F=P_{\rm SM}^F\oplus P_{\rm mir}^F\oplus P_{\rm ex}^F,
 \label{eq:future-visible-inventory-split}
\end{equation}
and hence
\begin{equation}
 \Delta_{\rm inv}^F:=\Tr(P_{\rm gch}^F-P_{\rm SM}^F)
     =\rank P_{\rm mir}^F+\rank P_{\rm ex}^F.
 \label{eq:future-visible-inventory-residual}
\end{equation}
Thus $\Delta_{\rm inv}^F=0$ exactly when the complete future-visible
nontrivial gauge carrier consists of the target chiral types.  A nonzero
residual is at least one, and its projection has operator norm one.
\end{proposition}
\begin{proof}
The tester is a sum of positive squared defects, so its kernel is the
intertwiner space for the stated gauge type, charge, and chirality.
Schur's lemma gives $T_a^*T_b=c_{ab}I$; the normalized multiplicity inner
product gives $c_{ab}=\delta_{ab}$.  The resulting isometries have orthogonal
ranges and exhaust the isotypic component, proving both the projector
formula and unitary evaluation.  The gauge action and commuting chirality
grading give an orthogonal decomposition into these isotypic spaces.  The
remaining nontrivial gauge part is $P_{\rm ex}^F$, and the trace of each
finite projection is its rank.  Cyclic saturation excludes decoupled ambient
spectators but includes every sourced or reached gauge submodule together
with its gauge orbit.
\end{proof}

The fully gauge-trivial sector, rather than the entire zero-hypercharge
sector, is removed.  A nonabelian multiplet of zero hypercharge is still
gauge charged; Appendix~\ref{app:anomaly} gives an explicit exotic pair that
would be missed by a hypercharge-only audit.

Write
\begin{equation}
 n_f:=n_{f,\epsilon_f},\qquad f\in\{Q,u,d,L,e\},
 \label{eq:charged-multiplicities}
\end{equation}
and, with the determinant-normalized integer charge table of
\eqref{eq:hypercharges}, define the four normalized linear anomaly rows
\begin{align}
 r_3&=2n_Q-n_u-n_d,\nonumber\\
 r_{3Y}&=n_Q-2n_u+n_d,\nonumber\\
 r_{2Y}&=n_Q-n_L,\nonumber\\
 r_{\rm gY}&=n_Q-2n_u+n_d-n_L+n_e,
 \label{eq:linear-anomaly-rows}
\end{align}
and
\begin{equation}
 \Delta_{\rm lin}
 =r_3^2+r_{3Y}^2+r_{2Y}^2+r_{\rm gY}^2.
 \label{eq:linear-anomaly-residual}
\end{equation}

\begin{proposition}[Linear anomaly rigidity of the charged multiplicities]
\label{thm:linear-anomaly-rigidity}
For nonnegative integer multiplicities with $n_Q>0$,
\begin{equation}
 \Delta_{\rm lin}=0
 \iff n_Q=n_u=n_d=n_L=n_e=:g>0.
 \label{eq:linear-anomaly-common-count}
\end{equation}
On this branch the cubic $U(1)^3$ anomaly vanishes automatically and the
number of left-handed weak doublets is $4g$, hence even.  Moreover
\begin{equation}
 \Delta_{\rm lin}\ne0\Longrightarrow\Delta_{\rm lin}\ge1.
 \label{eq:linear-anomaly-integer-gap}
\end{equation}
Thus, after the central charge table is fixed, the cubic anomaly is not an
independent multiplicity constraint.
If also $\Delta_{\rm inv}^F=0$ and one charged multiplicity is independently
certified as three, the complete charged carrier is
$3Q\oplus3u\oplus3d\oplus3L\oplus3e$, with no future-visible charged
mirror or exotic sector.  Gauge-trivial singlets remain unconstrained.
\end{proposition}

\begin{proof}
The equations $r_3=r_{3Y}=0$ give
$n_d=2n_Q-n_u$ and then $3(n_Q-n_u)=0$, so
$n_u=n_d=n_Q$.  The equation $r_{2Y}=0$ gives $n_L=n_Q$, and
$r_{\rm gY}=0$ then gives $n_e=n_Q$.  The converse is immediate.  For a common
value $g$, the cubic anomaly is $g$ times the already vanishing one-generation
cubic anomaly of Proposition~\ref{thm:hypercharge}; the weak-doublet count is
$3g+g=4g$.  Finally a nonzero sum of squares of integers is at least one.
\end{proof}

The anomaly relations make all charged multiplicities equal, but do not
select the common integer.  The endpoint source supplies a possible
rank-three anchor.  We must still verify its occurrence in a charged type
and test whether its image exhausts that type's complete multiplicity.


\subsection{Endpoint ancestry and the charged-copy count}
\label{sec:endpoint-allocation}

On the positive endpoint branch of
Proposition~\ref{thm:geometry-short}, denote the supported complex endpoint
space by $G_{\rm gen}$.  Then
\begin{equation}
 \dim_\C G_{\rm gen}=3.\label{eq:three-generations}
\end{equation}
The odd sign-twisted loop triplet is a different source sector and is not
included.  Its allocation to matter would require separate ancestry data.
For a populated type set $N_f=\mathcal I_{f,\epsilon_f}$, with the unitary
evaluation map above.  Lemma~\ref{lem:complete-typed-source} obtains a
surjective synthesis $S_f:E_f\to N_f$ by contracting a terminal fermion
source-word bank against $V_f$.  Redundant source columns are allowed.

\begin{proposition}[Gauge-covariant endpoint-ancestry compiler]
\label{thm:endpoint-ancestry-compiler}
Fix a populated charged type $f$, its canonical multiplicity space
$N_f=\mathcal I_{f,\epsilon_f}$, and the supported endpoint carrier
$G_{\rm gen}$ of Proposition~\ref{thm:geometry-short}.  Let
\begin{equation}
\Xi_f:V_f\otimes G_{\rm gen}
\longrightarrow P_{f,\epsilon_f}\mathcal K_F^{\min}
\label{eq:endpoint-ancestry-occurrence}
\end{equation}
be a represented same-carrier endpoint-ancestry occurrence, and put
\begin{equation}
\widehat\Xi_f:=\operatorname{ev}_f^*\Xi_f:
V_f\otimes G_{\rm gen}\longrightarrow V_f\otimes N_f.
\label{eq:endpoint-ancestry-evaluation}
\end{equation}
Using the finite gauge generators $X_1,\ldots,X_s$, the central generator
$Y$, and the chirality grading from
Proposition~\ref{thm:canonical-chiral-projectors}, define the positive covariance
defect
\begin{equation}
\begin{aligned}
\Delta_f^{\rm anc}(\Xi_f):={}&
\sum_{j=1}^s
\left\|X_j\Xi_f-
\Xi_f(\rho_f(X_j)\otimes I_{G_{\rm gen}})\right\|_{\HS}^2\\
&+\left\|Y\Xi_f-y_f\Xi_f\right\|_{\HS}^2
 +\left\|\Gamma_F\Xi_f-\epsilon_f\Xi_f\right\|_{\HS}^2.
\end{aligned}
\label{eq:endpoint-ancestry-defect}
\end{equation}
Then
\begin{equation}
\Delta_f^{\rm anc}(\Xi_f)=0
\quad\Longleftrightarrow\quad
\widehat\Xi_f=I_{V_f}\otimes T_f
\text{ for a unique }T_f:G_{\rm gen}\to N_f.
\label{eq:endpoint-ancestry-factorization}
\end{equation}
On the zero-defect branch the multiplicity map is recovered canonically by
\begin{equation}
T_f=\frac1{d_f}\Tr_{V_f}(\widehat\Xi_f),
\qquad d_f=\dim V_f,
\label{eq:endpoint-ancestry-partial-trace}
\end{equation}
where $\Tr_{V_f}$ denotes contraction of the input and output $V_f$ legs,
and obeys
\begin{equation}
\Xi_f^*\Xi_f=I_{V_f}\otimes T_f^*T_f,
\qquad
\|\Xi_f\|_{\HS}^2=d_f\|T_f\|_{\HS}^2.
\label{eq:endpoint-ancestry-Gram}
\end{equation}
Consequently endpoint faithfulness is the finite positive condition
$T_f^*T_f\succ0$, equivalently strict positivity of
$\Xi_f^*\Xi_f$ on $V_f\otimes G_{\rm gen}$.  Thus, once the represented
cross-source occurrence \eqref{eq:endpoint-ancestry-occurrence} is supplied,
the endpoint map used in the allocation theorem below is not an arbitrary
rank-three embedding: it is a unique coefficient extracted from the physical
same-carrier operator.
\end{proposition}

\begin{proof}
Zero defect means that $\Xi_f$ intertwines the represented finite
gauge-action algebra with the stated charge and chirality.  Transport by
the unitary evaluation map gives an intertwiner from
$V_f\otimes G_{\rm gen}$ to $V_f\otimes N_f$.  Irreducibility of $V_f$
and Schur's lemma give the unique coefficient $T_f$; conversely any
$I_{V_f}\otimes T_f$ satisfies these equations.  Contracting the $V_f$ legs
gives the partial-trace formula, and taking the Gram and trace gives the
remaining identities and the equivalence of strict positivity.
\end{proof}

This compiler removes a coordinate choice, not an occurrence hypothesis:
$\Xi_f$ is a represented cross-source operator.  Old word moments cannot
supply it when it is absent from the historical algebra
(Proposition~\ref{thm:external-shadow-no-go}).  Once it occurs faithfully,
the following theorem tests surjectivity without assuming it.

\begin{theorem}[Endpoint allocation and the exact charged-copy count]
\label{thm:endpoint-allocation}
Fix one populated charged type $f$ and a complete represented typed synthesis
$S_f:E_f\to N_f$, for example \eqref{eq:complete-typed-source-synthesis}.
Let $T_f:G_{\rm gen}\to N_f$ be the canonical coefficient of a represented
ancestry occurrence supplied by Proposition~\ref{thm:endpoint-ancestry-compiler},
or more generally a specified represented endpoint subsource, satisfying
\begin{equation}
 A_f:=T_f^*T_f\succ0,\qquad
 P_{fG}:=T_fA_f^{-1}T_f^*.
 \label{eq:endpoint-subsource-projector}
\end{equation}
No surjectivity of $T_f$ is assumed.  Put
\begin{equation}
 G_f=S_f^*S_f,\qquad C_f=T_f^*S_f,\qquad
 \mathbb A_{f\mid G}:=G_f-C_f^*A_f^{-1}C_f.
 \label{eq:endpoint-allocation-short}
\end{equation}
Then
\begin{equation}
 \mathbb A_{f\mid G}
        =S_f^*(I-P_{fG})S_f\succeq0,\qquad
 a_{f\mid G}:=\Tr\mathbb A_{f\mid G}
        =\|(I-P_{fG})S_f\|_{\HS}^2,
 \label{eq:endpoint-allocation-positive}
\end{equation}
and
\begin{equation}
 n_f=\dim(E_f/\Ker G_f)=\rank G_f
       =3+\rank\mathbb A_{f\mid G}.
 \label{eq:endpoint-generation-excess}
\end{equation}
In particular,
\begin{equation}
 a_{f\mid G}=0
 \iff\mathbb A_{f\mid G}=0
 \iff\Ran T_f=N_f
 \iff n_f=3.
 \label{eq:endpoint-allocation-zero}
\end{equation}
The quotient carries the inner product induced by $G_f$, and $S_f$ descends
to a unitary from that quotient onto $N_f$.
If, in addition, $\Delta_{\rm inv}^F=\Delta_{\rm lin}=0$, then
\begin{equation}
 n_Q=n_u=n_d=n_L=n_e=3+\rank\mathbb A_{f\mid G}.
 \label{eq:allocation-anomaly-excess}
\end{equation}
Consequently $a_{f\mid G}=0$ gives exactly three copies of every charged
target type and no future-visible charged exotic or mirror sector.
Fully gauge-trivial singlets require separate occurrence and multiplicity data.
\end{theorem}

\begin{proof}
Faithfulness of $T_f$ and $\dim G_{\rm gen}=3$ make $P_{fG}$ the orthogonal
projection onto a three-dimensional subspace of $N_f$.  Substitution proves
\eqref{eq:endpoint-allocation-positive}.  Since $S_f$ is onto,
\[
 N_f=\Ran T_f\oplus(I-P_{fG})N_f,\qquad
 \Ran((I-P_{fG})S_f)=(I-P_{fG})N_f.
\]
For $V=(I-P_{fG})S_f$, the identity $\rank(V^*V)=\rank V$ therefore gives
$\rank\mathbb A_{f\mid G}=n_f-3$.  Also
$\Ker G_f=\Ker S_f$, which proves the quotient and rank statements.
Positivity makes zero trace equivalent to the zero residual.  Its vanishing
is equivalent to $\Ran S_f\subseteq\Ran T_f$; completeness of $S_f$ and
$\rank T_f=3$ give the remaining equivalences.  Finally,
$\langle e,G_fe'\rangle=\langle S_fe,S_fe'\rangle$ proves the assertion
about the quotient inner product.
The four linear anomaly rows force every charged multiplicity to equal
$n_f=3+\rank\mathbb A_{f\mid G}$ by
Proposition~\ref{thm:linear-anomaly-rigidity}, while the zero-inventory test
excludes every other nontrivial gauge-charged sector.
\end{proof}

A small positive allocation trace is a leakage bound, not an exact
three-copy conclusion.  Appendix~\ref{app:source-allocation-certificates}
proves the rank-separating threshold under an independently controlled
source frame bound and gives counterexamples without such control.
The same appendix records the source-native and same-source coefficient
formulas, retaining their occurrence hypotheses.


\subsection{Generation factors and relative frames}

\begin{proposition}[Abstract replication and source-determined generation frames]
\label{thm:common-generation-factor}
Suppose all populated charged multiplicities $N_f$ have dimension $g$.
For any $g$-space $G$ and arbitrary unitaries $W_f:G\to N_f$,
\begin{equation}
 \bigoplus_fV_f\otimes N_f\cong\left(\bigoplus_fV_f\right)\otimes G,
 \label{eq:abstract-global-generation-factor}
\end{equation}
and a coefficient arrow $B:N_f\to N_h$ becomes $W_h^*BW_f\in M_g(\C)$.
Thus the three-copy conclusion already supplies an abstract common
$\C^3$ factor, without common endpoint ancestry.

For each connected incidence component $C$, suppose instead that a
specified rank-three endpoint source $G_C$ occurs with maps
$S_f:G_C\to N_f$ satisfying $S_f^*S_f\succ0$ and $\Ran S_f=N_f$.
Their polar maps $U_f=S_f(S_f^*S_f)^{-1/2}$ are unitaries and give
\begin{equation}
 \mathcal H_{\rm mat}\cong\bigoplus_{C\in\pi_0(\Gamma)}
       \left(\bigoplus_{f\in C}V_f\right)\otimes G_C,
 \qquad\dim G_C=3.
 \label{eq:componentwise-generation-factorization}
\end{equation}
The gauge and typing algebras act trivially on $G_C$, and every arrow
inside $C$ becomes $U_h^*BU_f$.  For the usual quark and lepton components
this gives separate $G_q$ and $G_\ell$, when both component sources occur.

If a single represented endpoint source $G_{\rm gen}$ supplies all types,
these polar maps determine
\begin{equation}
 \mathcal H_{\rm mat}\cong\mathcal H_F^{(1)}\otimes G_{\rm gen},\qquad
 \mathcal H_F^{(1)}=\bigoplus_fV_f,\quad\dim G_{\rm gen}=3.
 \label{eq:common-generation-factorization}
\end{equation}
An arrow preserves the full common generation algebra exactly when
\begin{equation}
 U_h^*BU_f\in\C I_{G_{\rm gen}},\label{eq:generation-neutral-incidence}
\end{equation}
equivalently when
\begin{equation}
 \sum_{i,j=1}^3\|BU_fE_{ij}U_f^*-U_hE_{ij}U_h^*B\|_{\HS}^2=0.
 \label{eq:generation-compatibility-defect}
\end{equation}
Arbitrary Yukawa coefficients remain ordinary matrices in this frame;
their surviving commutant is determined by Theorem~\ref{thm:main-duality}.

More generally, for equal-dimensional component fibres, if every retained
gauge, typing, source, and incidence operator is block diagonal across
components, the unselected relative coordinate frames form
\begin{equation}
 U(g)^{|\pi_0(\Gamma)|}/U(g)_{\rm diag}.
 \label{eq:relative-generation-frame-space}
\end{equation}
For the quark and lepton components with $g=3$, this is one relative $U(3)$.
No datum internal to the disconnected packet selects a point of this space.
\end{proposition}
\begin{proof}
The direct sum of $I_{V_f}\otimes W_f^*$ gives the first factorization and
transports all arrows.  Faithful surjective endpoint maps have unitary polar
factors, giving the componentwise and common-source identifications.
In common coordinates an arrow intertwines every generation matrix exactly
when it commutes with all matrix units, hence is scalar; the sum of squared
defects tests that condition.  To compare disconnected components, choose
$J_C:G\to G_C$.  Replacing $J_C$ by $J_CU_C$ conjugates the complete
represented tuple by
\begin{equation}
 \mathcal U_{\bm U}=\bigoplus_C I_{\mathcal V_C}\otimes U_C^*.
 \label{eq:relative-generation-coordinate-change}
\end{equation}
Words, Grams, and commutants transform covariantly, while spectra, singular
values, trace polynomials, and scalar occurrence residuals are unchanged.
A common right multiplication changes only the basis of the reference
space, leaving the quotient \eqref{eq:relative-generation-frame-space}.
There is no represented intercomponent relation to select a relative frame.
\end{proof}


\subsection{Grading change, weak copies, and finite Dirac relations}

These remaining tests concern occurrence and the structure of a specified
operator, not the charged-copy count.  For a protected grading
$J=P_+-P_-$ and an active completely positive map $\Phi$, the odd-Choi mass is
\[
 \mu_J(\Phi)=\Tr[P_-\Phi(P_+)]+\Tr[P_+\Phi(P_-)].
\]
It is Kraus independent and vanishes exactly when the Choi support consists
of grading-even operators.  Appendix~\ref{app:odd-occurrence} proves this
criterion, its finite-alphabet stopping alternative, and its invariance
under unread trace-preserving spectator lifts.  Positivity certifies active
grading change; it does not by itself select a weak-doublet copy.

\begin{proposition}[Positive weak-copy census]
\label{thm:weak-copy-census}
Let $W_2\cong\C^2$ be the represented weak defining module, let $M_H$ be a
finite multiplicity space, and let $J_H\succeq0$ be the complete actual
positive weak-scalar occurrence packet on the retained sector
$W_2\otimes M_H$.  Define
\begin{equation}
 B_H:=\Tr_{W_2}J_H\succeq0,
 \label{eq:weak-copy-Gram}
\end{equation}
and let its gauge-cyclic occurrence carrier be
\begin{equation}
 \mathcal K_H^{\rm occ}
 :=\Span\{(g\otimes I)v:g\in\SU(W_2),\ v\in\Ran J_H\}.
 \label{eq:weak-gauge-cyclic-carrier}
\end{equation}
For normalized Haar measure,
\begin{equation}
 \overline J_H:=\int_{\SU(W_2)}
 (g\otimes I)J_H(g^*\otimes I)\,\dd g
 =\frac12I_{W_2}\otimes B_H,
 \label{eq:weak-copy-twirl}
\end{equation}
and
\begin{equation}
 \mathcal K_H^{\rm occ}=\Ran\overline J_H
       =W_2\otimes\Ran B_H.
 \label{eq:weak-gauge-cyclic-support}
\end{equation}
Thus the source-minimal gauge-cyclic number of reached weak-doublet copies is
\begin{equation}
 g_H^{\rm occ}=\rank B_H.
 \label{eq:weak-copy-count}
\end{equation}
In particular, there is exactly one such doublet if and only if
\begin{equation}
 m_H:=\Tr B_H>0,\qquad
 \mathfrak p_H:=m_H^2-\Tr(B_H^2)=0.
 \label{eq:one-weak-doublet-test}
\end{equation}
The criterion is independent of Kraus coordinates, phases, and the choice of
basis of $M_H$.
\end{proposition}

\begin{proof}
Schur averaging gives $\overline J_H=I_{W_2}\otimes B_H/2$.
The kernel of a positive group average is the intersection of the kernels
of its conjugates, so its range is precisely the gauge-cyclic occurrence
carrier.  Hence the copy count is $\rank B_H$.  If $\lambda_i\ge0$ are the
eigenvalues of $B_H$, then
$m_H^2-\Tr(B_H^2)=2\sum_{i<j}\lambda_i\lambda_j$; together with $m_H>0$,
its vanishing is equivalent to rank one.  All constructions are independent
of coefficient coordinates.  The kernel-of-the-average argument is written
out in Appendix~\ref{app:weak-cyclic-example}.
\end{proof}

The count concerns the complete retained packet's gauge-cyclic support,
not an unused ambient multiplicity.  The rank of $J_H$ itself need not be
the copy count: the same appendix gives a rank-one packet with two weak copies.

Let $A$ be a source-saturated self-adjoint full-cell generator on the common
carrier and let
\begin{equation}
 \widetilde D_F=P_-AP_+ + P_+AP_-
 =\frac12(A-JAJ)
 \label{eq:finite-dirac}
\end{equation}
be its grading-odd component.

\begin{definition}[Finite-Dirac relation space]
\label{def:finite-dirac-relation-space}
Let $\mathfrak D_F$ be the finite-dimensional $*$-closed linear subspace of
grading-odd operators whose typed blocks are supported only on the represented
fermionic Yukawa incidence relations determined by the matter packet; when the
neutral and charge-conjugate sectors are represented, the declared Majorana
relation is included as well.  Let
$\Pi_F^{\rm D}$ be the Hilbert--Schmidt orthogonal projection onto
$\mathfrak D_F$ and define
\begin{equation}
 D_F^{\rm can}=\Pi_F^{\rm D}\widetilde D_F,
 \qquad
 \Delta_F^{\rm rel}
 =\|(I-\Pi_F^{\rm D})\widetilde D_F\|_{\HS}^2.
 \label{eq:finite-dirac-relation-defect}
\end{equation}
\end{definition}

\begin{proposition}[Fixed-generator finite-Dirac closure]
\label{prop:finite-dirac}
For the specified self-adjoint generator $A$,
\begin{equation}
 \widetilde D_F=D_F^{\rm can}
 \iff \Delta_F^{\rm rel}=0.
 \label{eq:finite-dirac-exact-relation}
\end{equation}
Moreover
\begin{equation}
 \operatorname{dist}_{\HS}(\widetilde D_F,\mathfrak D_F)
 =\sqrt{\Delta_F^{\rm rel}}.
 \label{eq:finite-dirac-distance}
\end{equation}
If $\Delta_F^{\rm rel}=0$ and $\|D_F^{\rm can}\|_{\HS}>0$, the complete odd
component of this generator is a nonzero admitted finite-Dirac operator.  In
any generation factorization of
Proposition~\ref{thm:common-generation-factor}, its typed coefficient
blocks are ordinary generation matrices.  If the stronger common-ancestry
conditions of Proposition~\ref{thm:common-generation-factor} hold, those matrices
are expressed in the corresponding source-determined common frame.
Odd-provenance completeness is not required for these fixed-generator
conclusions; when imposed, it additionally excludes independent future-visible
odd source directions outside the declared odd synthesis.
\end{proposition}

\begin{proof}
Orthogonal decomposition relative to $\mathfrak D_F$ gives
\[
 \widetilde D_F
 =\Pi_F^{\rm D}\widetilde D_F
 +(I-\Pi_F^{\rm D})\widetilde D_F.
\]
The second summand has squared Hilbert--Schmidt norm
$\Delta_F^{\rm rel}$, proving both displayed statements.  The remaining
claims follow from the typed decomposition and the generation-coordinate
identifications already established above.
\end{proof}

Odd-provenance completeness, when imposed, means that every future-visible
grading-odd action of the saturated carrier lies in its represented odd
source span.  This stronger exhaustion condition is not needed for the
fixed-generator projection test above.

\begin{corollary}[Structural synthesis]
\label{thm:structural-SM}
Suppose a same-source, Clifford-compatible packet satisfies the census and
internal-landing hypotheses of Theorem~\ref{thm:internal-seed-saturation}.
Suppose also that its global group is obtained either from the aligned,
phase-sensitive determinant data of Theorem~\ref{thm:gauge-group}, or from
the represented chiral tensor packet, primitive anomaly freedom, and faithful
kernel of Corollary~\ref{cor:anomaly-faithful-group}.
On its complete source-minimal fermion carrier, assume the faithful endpoint
occurrence and complete typed synthesis of Theorem~\ref{thm:endpoint-allocation},
with $\Delta_{\rm inv}^F=\Delta_{\rm lin}=a_{f_0\mid G}=0$ for one charged
type $f_0$.  Then
\begin{align}
 \A_{\rm int}^{\rm word}&\cong M_3(\C)\oplus M_2(\C)\oplus\C\oplus\C,
 \nonumber\\
 G_{\rm phys}&\cong S(\U(3)\times\U(2))
       \cong(\SU(3)\times\SU(2)\times\U(1))/\Z_6,\nonumber\\
 \mathcal K_{F,\rm gch}^{\min}&\cong3Q\oplus3u\oplus3d\oplus3L\oplus3e.
 \label{eq:structural-synthesis}
\end{align}
The canonical matter action supplies the typing algebra, so all its represented
incidences, of arbitrary coefficient rank, satisfy the general duality.
The independent weak-copy and fixed-generator tests above give respectively
one reached weak doublet when $m_H>0$, $\mathfrak p_H=0$, and an admitted
nonzero finite-Dirac operator when $\Delta_F^{\rm rel}=0$,
$\|D_F^{\rm can}\|_{\HS}>0$.
\end{corollary}
\begin{proof}
Combine the internal saturation and Clifford lift, the applicable global-group
route, the endpoint-allocation count, and the two independent occurrence and
operator tests.  Proposition~\ref{prop:typed-matter-lift} then supplies the
typing algebra for Theorem~\ref{thm:main-duality}.
\end{proof}

The synthesis retains the distinction between classification and occurrence:
a failed test returns its smaller algebra, missing source direction, or
nonzero residual.  It determines no numerical Yukawa or Majorana entries,
flavour hierarchy, CKM or PMNS parameters, gauge or Higgs couplings,
renormalization-group data, or fermionic measure phase.


\section{Rigidity and refinement}
\label{sec:certificate}

The preceding sections determine exact finite algebras.  We now ask how
that exactness can be certified and maintained.  One positive commutator
operator governs three distinct statements: recognition at a fixed packet,
perturbative stability with a protected algebra, and cofinal stability on
varying screened spaces.  None of these statements supplies a missing
physical generator or a missing source occurrence.


\subsection{A positive commutant certificate}

\begin{theorem}[Finite commutant certificate and coercivity]
\label{thm:howe-certificate}
Let $c_1,\ldots,c_s$ span a unital $*$-algebra
$\mathcal C\subseteq\mathcal B(\Hh)$.  On Hilbert--Schmidt operator space,
define the commutant Laplacian
\begin{equation}
 \mathscr L_{\mathcal C}=\sum_j\operatorname{ad}_{c_j}^*
                         \operatorname{ad}_{c_j},\qquad
 \operatorname{ad}_c(X)=[c,X].
 \label{eq:commutant-laplacian}
\end{equation}
Then $\mathscr L_{\mathcal C}\succeq0$ and
\begin{equation}
 \Ker\mathscr L_{\mathcal C}=\mathcal C'.
 \label{eq:commutant-kernel}
\end{equation}
If $\lambda_{\mathcal C}>0$ is its least positive eigenvalue,
\begin{equation}
 \|X-P_{\mathcal C'}X\|_{\HS}^2
 \le\lambda_{\mathcal C}^{-1}\sum_j\|[c_j,X]\|_{\HS}^2.
 \label{eq:commutant-coercivity}
\end{equation}
For a proposed protected algebra $\M\subseteq\mathcal C'$ and an
orthonormal basis $(E_a)$ of $\M^\perp$, define the relative Howe Gram
\begin{equation}
 \mathbb G_{\rm Howe}(\bm c\mid\M)_{ab}
   =\sum_j\langle[c_j,E_a],[c_j,E_b]\rangle_{\HS}.
 \label{eq:howe-gram}
\end{equation}
Then
\begin{equation}
 \mathbb G_{\rm Howe}(\bm c\mid\M)\succ0
 \quad\Longleftrightarrow\quad\mathcal C'=\M.
 \label{eq:howe-certificate}
\end{equation}
On this branch its least eigenvalue, when the complement is nonzero, is the
first positive eigenvalue of $\mathscr L_{\mathcal C}$.
\end{theorem}
\begin{proof}
The Hilbert--Schmidt adjoint of $\operatorname{ad}_c$ is
$\operatorname{ad}_{c^*}$, so the quadratic form is
$\sum_j\|[c_j,X]\|_{\HS}^2$.  Its zero set is exactly the simultaneous
commutant.  The spectral theorem gives coercivity on the orthogonal
complement of that kernel.  The relative Howe Gram is precisely the
compression to $\M^\perp$; since $\M$ is already in the kernel, strict
positivity there is equivalent to absence of every additional commutant
direction.
\end{proof}

For typed incidences, the type-central part of this certificate is a weighted
graph Laplacian.  Its connected-component kernel and its Schur separation
from noncentral flavour directions are proved in
Appendix~\ref{app:incidence-flavour-stability}.  The full test remains an
operator-algebraic one, not a dimensional match.


\subsection{Enhancement loci and finite perturbations}
\label{sec:howe-discriminant}

The positive certificate also determines the symmetry-enhancement locus of a
parameter family.  Let $\Theta\subset\R^p$ be an open parameter domain, keep a
fixed represented protected algebra $\M\subseteq\mathcal B(\Hh)$, and assume
$[c_j(\theta),\M]=0$ for every $j$ and $\theta$.  On
$\Kk:=\M^\perp$ define the stacked commutator map
\[
 D(\theta)X=([c_1(\theta),X],\ldots,[c_s(\theta),X]),
 \qquad G(\theta)=D(\theta)^*D(\theta).
\]
In an orthonormal basis of $\Kk$, $G(\theta)$ is the relative Howe Gram.
Put $r=\dim_\C\Kk$ and
\begin{equation}
 \Delta_{\M}(\theta)=\det G(\theta).
 \label{eq:howe-discriminant}
\end{equation}

\begin{proposition}[Howe enhancement discriminant and generic rigidity]
\label{thm:howe-discriminant}
Assume every matrix entry of every $c_j(\theta)$ is a complex polynomial in
the real coordinates $\theta$, of degree at most $d$.  Then
$\Delta_{\M}$ is a nonnegative real polynomial of degree at most $2dr$ and
admits the Cauchy--Binet representation
\begin{equation}
 \Delta_{\M}(\theta)
 =\sum_{|I|=r}|\det D_I(\theta)|^2.
 \label{eq:howe-minor-sos}
\end{equation}
Moreover,
\begin{equation}
 \{\theta:\Delta_{\M}(\theta)=0\}
 =\{\theta:C^*(\bm c(\theta))'\supsetneq\M\}.
 \label{eq:howe-enhancement-locus}
\end{equation}
If the proposed commutant is exact at one parameter value, then
$\Delta_{\M}$ is not the zero polynomial.  Hence its enhancement locus has
empty interior and Lebesgue measure zero, while exact duality holds on an open
dense full-measure subset of every open admissible parameter region.  The same
open-dense conclusion holds for a real-analytic family on each connected
analytic component containing one exact point.
\end{proposition}

\begin{proof}
The entries of $D(\theta)$ are polynomial of degree at most $d$.
Cauchy--Binet applied to $D^*D$ gives
\eqref{eq:howe-minor-sos}; each squared modulus is a sum of squares of real
polynomials, proving nonnegativity and the degree bound.  Since $G\succeq0$,
its determinant is positive exactly when $G\succ0$, which is equivalent to
$C^*(\bm c(\theta))'=\M$ by Theorem~\ref{thm:howe-certificate}.  A nonzero
real polynomial has a zero set of Lebesgue measure zero and empty interior;
continuity makes the complement open.  For a connected real-analytic family,
the identity theorem gives the corresponding empty-interior conclusion once
the determinant is nonzero at one point.
\end{proof}

\begin{corollary}[Determinantal nullity filtration]
\label{cor:howe-nullity}
For $1\le k\le r$, the locus on which the commutant acquires at least $k$
additional complex dimensions is the common zero set of all
$(r-k+1)\times(r-k+1)$ minors of $D(\theta)$.  Additional commutant dimension
is therefore upper semicontinuous in every continuous family.
\end{corollary}

The genericity is relative to a specified admissible family with a protected
algebra; it does not imply that a microscopic source realizes that family.
A finite quantitative statement is available around any exact anchor.

\begin{proposition}[Finite-anchor stability and generator-scale control]
\label{thm:finite-anchor}
Transport a family of tuples $\bm c_X=(c_{1,X},\ldots,c_{s,X})$ to one
common finite screened carrier, and suppose a fixed algebra $\M$ commutes
with every tuple.  Set
\[
 C_X=\left(\sum_j\|c_{j,X}\|_{\rm op}^2\right)^{1/2},\qquad
 \varepsilon_{X,Y}=\left(\sum_j\|c_{j,X}-c_{j,Y}\|_{\rm op}^2\right)^{1/2}.
\]
Assume $\mathbb G_{\rm Howe}(\bm c_{X_0}\mid\M)\succeq\gamma_0I$,
$\gamma_0>0$.  At every later $X$ for which
\begin{equation}
 4(C_X+C_{X_0})\varepsilon_{X,X_0}<\gamma_0,
 \label{eq:anchor-tail}
\end{equation}
one has
\begin{equation}
 C^*(\bm c_X)'=\M,\qquad
 \lambda_{\min}^+(\mathscr L_X)
 \ge\gamma_0-4(C_X+C_{X_0})\varepsilon_{X,X_0}.
 \label{eq:anchor-lock}
\end{equation}
Alternatively, the condition
$2\varepsilon_{X,X_0}<\sqrt{\gamma_0}$ gives
\begin{equation}
 C^*(\bm c_X)'=\M,\qquad
 \lambda_{\min}^+(\mathscr L_X)
 \ge(\sqrt{\gamma_0}-2\varepsilon_{X,X_0})^2,
 \label{eq:generator-scale-gap}
\end{equation}
without an upper bound on the full generator norm.  If the corresponding
condition holds for every later $X$, the conclusion holds throughout that tail.
\end{proposition}
\begin{proof}
For the stacked commutator maps,
$\|\partial_{\bm c_X}-\partial_{\bm c_Y}\|_{2\to2}\le2\varepsilon_{X,Y}$
and $\|\partial_{\bm c_X}\|_{2\to2}\le2C_X$.
Thus $\|\mathscr L_X-\mathscr L_Y\|\le
4(C_X+C_Y)\varepsilon_{X,Y}$, giving the first bound on $\M^\perp$.
Directly on that complement, the triangle inequality gives
$\|\partial_{\bm c_X}A\|\ge
(\sqrt{\gamma_0}-2\varepsilon_{X,X_0})\|A\|$, giving the second.
The protected algebra stays in the kernel, so either positive lower bound
makes that kernel exact.
\end{proof}

Coefficient-space and singular-support stability require a separate
singular-value gap.  Appendix~\ref{app:coefficient-stability} retains the
operator--Schmidt cluster estimate and the one-sided rank certificate, and
Appendix~\ref{app:polar-details} gives support leakage.  In particular,
small measured tails do not certify exact absence without an independent
upper-rank premise.


\subsection{Cofinal kernel exactness}
\label{sec:compact-mosco}

A cofinal family need not remain in a norm neighbourhood on a fixed carrier.
We therefore formulate a second stability mechanism on transported screened
Hilbert spaces.  Kernel exactness will use Mosco convergence, finite protected
transport, and a uniform complement gap; compactness will be added only for
the stronger spectral conclusion.

At a finite cutoff the screened operator space is the full
Hilbert--Schmidt space of the finite physical screen on which the displayed
generators act.  If instead a proper invariant operator test space is used,
all kernel statements are relative to that space and do not identify an
unseen commutant on a larger carrier.

Let $\mathscr H$ be a separable Hilbert space and let
$J_X:\mathscr H_X\to\mathscr H$ be isometric embeddings of the screened
Hilbert--Schmidt operator spaces, with $P_X=J_XJ_X^*$ the corresponding
orthogonal projections.  The limiting space is embedded by $J_\infty$, and we
assume the transport is compatible in the sense that
$P_X\to P_\infty$ strongly on $\mathscr H$.  For a densely defined closed
nonnegative form $q_X$ on $\mathscr H_X$, denote
its self-adjoint operator by $\mathscr L_X$ and its transported resolvent by
\begin{equation}
R_X:=J_X(\mathscr L_X+I)^{-1}J_X^*\in\mathcal B(\mathscr H).
\label{eq:transported-resolvent}
\end{equation}

\begin{definition}[Transported Mosco convergence]
\label{def:transported-mosco}
The forms $q_X$ converge to $q_\infty$ in the transported Mosco sense when:
\begin{enumerate}[label=\textup{(M\arabic*)},leftmargin=3em]
\item whenever $J_XA_X\rightharpoonup B$ weakly in $\mathscr H$, the strong
      convergence $P_X\to P_\infty$ gives $B\in J_\infty\mathscr H_\infty$;
      writing $B=J_\infty A$,
\[
 q_\infty(A)\le\liminf_X q_X(A_X);
\]
\item for every $A\in D(q_\infty)$ there are $A_X\in D(q_X)$ such that
$J_XA_X\to J_\infty A$ strongly and
\[
 q_X(A_X)\longrightarrow q_\infty(A).
\]
\end{enumerate}
\end{definition}

\begin{lemma}[Finite protected basis transport]
\label{lem:protected-basis-transport}
Let $\M_X\subseteq\mathscr H_X$ and $\M_\infty\subseteq\mathscr H_\infty$
have the same finite complex dimension $m$.  Choose Hilbert--Schmidt
orthonormal bases $e_{1,X},\ldots,e_{m,X}$ and
$e_{1,\infty},\ldots,e_{m,\infty}$ and suppose
\begin{equation}
 \eta_X:=\max_a\|J_Xe_{a,X}-J_\infty e_{a,\infty}\|\longrightarrow0.
 \label{eq:protected-basis-transport}
\end{equation}
Then the transported orthogonal projections
$\widehat P_X=J_XP_{\M_X}J_X^*$ satisfy
\begin{equation}
 \|\widehat P_X-\widehat P_\infty\|_{\rm op}
 \le2m\eta_X\longrightarrow0.
 \label{eq:protected-projection-transport}
\end{equation}
If the $\M_X$ are images of one fixed finite $C^*$-algebra under specified
unital $*$-isomorphisms, the same data also identify multiplication and
adjoints across cutoffs.
\end{lemma}

\begin{proof}
Write the transported projections as sums of rank-one projections onto the
transported orthonormal basis vectors.  For unit vectors $u,v$,
$\||u\rangle\langle u|-|v\rangle\langle v|\|_{\rm op}\le2\|u-v\|$.
Summing over $m$ basis vectors proves the estimate.  The final assertion is
immediate from the specified $*$-isomorphisms.
\end{proof}

\begin{theorem}[Derived-kernel cofinal action--multiplicity duality]
\label{thm:cofinal-duality}
Let
\begin{equation}
 q_X(A)=\sum_j\norm{[c_{j,X},A]}_{\HS}^2
 \label{eq:cofinal-form}
\end{equation}
be the closed commutator forms of a cofinal family of screened support-polar
or multiplicity-quiver generators, with associated commutant Laplacians
$\mathscr L_X$.  Assume:
\begin{enumerate}[label=\textup{(D\arabic*)},leftmargin=3em]
\item $q_X\to q_\infty$ in transported Mosco sense;
\item for every $X$ there is a protected multiplicity algebra
      $\M_X\subseteq\Ker\mathscr L_X$ of the same finite complex dimension
      $m$, together with a specified $m$-dimensional $\M_\infty$ satisfying
      the basis transport condition
      \eqref{eq:protected-basis-transport};
\item for all sufficiently late cutoffs there is one constant $\gamma_*>0$
      such that
      \begin{equation}
       q_X(A)\ge\gamma_*
       \norm{(I-P_{\M_X})A}_{\HS}^2
       \qquad(A\in D(q_X)).
       \label{eq:uniform-cofinal-gap}
      \end{equation}
\end{enumerate}
Then for every sufficiently late cutoff
\begin{equation}
 \Ker\mathscr L_X=\M_X,\qquad C^*(\bm c_X)'=\M_X,
 \label{eq:cofinal-kernel}
\end{equation}
where the commutant identity is understood on the complete represented
Hilbert--Schmidt test space.  Moreover the limiting form satisfies
\begin{equation}
 q_\infty(A)\ge\gamma_*
 \norm{(I-P_{\M_\infty})A}_{\HS}^2,
 \label{eq:cofinal-limit-gap}
\end{equation}
and therefore
\begin{equation}
 \Ker\mathscr L_\infty=\M_\infty.
 \label{eq:cofinal-limit-kernel}
\end{equation}
Thus limiting kernel exactness is derived from finite protected transport and
a uniform late-cutoff gap; it is not an independent cofinal hypothesis.
Neither asymptotic form compactness nor compact resolvent is required for these
kernel statements.
\end{theorem}

\begin{proof}
At a finite cutoff, $A\in\Ker\mathscr L_X$ gives $q_X(A)=0$; the gap
\eqref{eq:uniform-cofinal-gap} forces $A\in\M_X$, and the reverse inclusion is
assumed.  The commutant identity is then
Theorem~\ref{thm:howe-certificate}.

Lemma~\ref{lem:protected-basis-transport} gives norm convergence of the
transported protected projections.  For every transported basis vector,
$e_{a,X}\in\M_X\subseteq\Ker\mathscr L_X$ gives $q_X(e_{a,X})=0$; the Mosco
liminf inequality and the strong convergence in
\eqref{eq:protected-basis-transport} therefore imply
$q_\infty(e_{a,\infty})=0$.  Hence
$\M_\infty\subseteq\Ker\mathscr L_\infty$.

For $A\in D(q_\infty)$ choose a Mosco recovery sequence $A_X$ with
$J_XA_X\to J_\infty A$ strongly and
$q_X(A_X)\to q_\infty(A)$.  Projection convergence implies
\[
 J_XP_{\M_X}A_X\longrightarrow
 J_\infty P_{\M_\infty}A,
\]
and hence
\[
 \|(I-P_{\M_X})A_X\|
 \longrightarrow\|(I-P_{\M_\infty})A\|.
\]
Passing to the limit in \eqref{eq:uniform-cofinal-gap} proves
\eqref{eq:cofinal-limit-gap}.  Its zero set is exactly $\M_\infty$, proving
\eqref{eq:cofinal-limit-kernel}.
\end{proof}


\subsection{Compact spectral upgrade}

\begin{definition}[Asymptotic form compactness]
\label{def:asymptotic-compactness}
The family $(q_X)$ is asymptotically compact if every sequence
$A_X\in D(q_X)$ satisfying
\begin{equation}
\sup_X\bigl(\norm{A_X}^2+q_X(A_X)\bigr)<\infty
\label{eq:asymptotic-compactness}
\end{equation}
has a subsequence for which $J_XA_X$ converges strongly in $\mathscr H$.
The same requirement is imposed on the limiting form.  In particular,
$(\mathscr L_\infty+I)^{-1}$ is compact on the limiting screened space.
\end{definition}

\begin{corollary}[Compact spectral upgrade]
\label{cor:cofinal-spectral-upgrade}
Assume Theorem~\ref{thm:cofinal-duality} and, in addition, asymptotic form
compactness in the sense of Definition~\ref{def:asymptotic-compactness}.  If
the limiting screened space has a nonzero complement to $\M_\infty$, then
\begin{equation}
 \gamma_\infty:=\lambda_{m+1}(\mathscr L_\infty)\ge\gamma_*>0,
 \label{eq:limit-gap}
\end{equation}
the transported resolvents converge in operator norm,
\begin{equation}
 \norm{R_X-R_\infty}_{\rm op}\longrightarrow0,
 \label{eq:cofinal-norm-resolvent}
\end{equation}
and
\begin{equation}
 \lambda_{m+1}(\mathscr L_X)\longrightarrow\gamma_\infty,\qquad
 \norm{\Pi_X-\Pi_\infty}_{\rm op}\longrightarrow0.
 \label{eq:cofinal-spectral-upgrade}
\end{equation}
Here $\Pi_X$ is the transported kernel projection.  Consequently, for every
$0<\gamma_0<\gamma_\infty$ and all sufficiently late $X$,
\begin{equation}
 q_X(A)\ge\gamma_0\norm{(I-P_{\M_X})A}_{\HS}^2.
 \label{eq:cofinal-coercivity}
\end{equation}
\end{corollary}

\begin{proof}
Theorem~\ref{thm:cofinal-duality} supplies exact limiting kernel dimension $m$
and the positive lower bound $\gamma_*$.  Proposition~\ref{prop:compact-mosco-resolvent}
gives norm-resolvent convergence.  Compact self-adjoint spectral convergence
then gives convergence of the first positive eigenvalue and of the isolated
zero-cluster projection.  The final estimate follows from the spectral theorem.
\end{proof}

Appendix~\ref{app:compact-mosco} proves the compact-Mosco resolvent mechanism
used in the corollary, including the varying-space bookkeeping
\cite{KuwaeShioya2003}.  It also records the distinct spectral locking
criterion that assumes an exact limiting kernel, and an escaping-mode
counterexample showing why protected transport cannot be replaced by a
dimension count.

The two stability mechanisms address different hypotheses.  The finite
anchor controls a norm perturbation of generators; the cofinal theorem
controls varying domains through forms.  Neither constructs an interacting
continuum field theory.  Identifying the protected spaces as one fixed
abstract $C^*$-algebra additionally requires multiplicatively compatible
transports, for example specified unital $*$-isomorphisms.  A particular
spectral screen construction may supply the analytic hypotheses, but is not
itself a premise of the theorem \cite{PelissierSpectral2026}.


\section{Interpretation and discussion}
\label{sec:discussion}

The main result of this paper is a finite representation-theoretic statement
about the relation between what acts on a matter carrier and what remains
invisible to that action.  Once an external Clifford action, a typed internal
action, and physical matter-incidence operators are represented on a common
finite carrier, the residual matter multiplicity is no longer merely a count
of repeated irreducible representations.  It becomes an intrinsic algebraic
object: the commutant of the complete represented action.  Conversely, finite
bicommutant closure reconstructs that action from the multiplicity algebra.

This observation separates two levels that are often conflated.  The
double-centralizer theorem is general at the level of finite joint packets and
does not single out any particular gauge theory.  The Standard Model appears
only after additional structural conditions are imposed on a particular
finite branch.  Likewise, the predictive construction discussed below gives
one origin for the required joint packet, but is not part of the algebraic
proof.  The discussion is organized according to these three levels:
the general action--multiplicity duality, its Standard-Model specialization,
and the relation between the abstract finite result and possible microscopic
or continuum realizations.


\subsection{What the duality determines}

The complete represented action and its residual multiplicity have the
same central decomposition and act on opposite matrix factors within each
sector.  Incidence is the information that restricts the otherwise large
commutant of the external--typing action.  The reconstruction is exact even
at singular incidence maps: support-polar data retain the support boundary,
and quiver and root metric--holonomy descriptions give alternative
presentations of the same algebra.

The commutant does not recover a preferred physical generator bank or its
numerical markings.  Incidence minimality is therefore relative to the
specified bank: the rank polymatroid identifies sufficient subsets and
which deletions enlarge the commutant.  Likewise, stability concerns a
protected algebra within an admitted family.  Its hypotheses do not provide
occurrence of the generators whose commutant is being studied.


\subsection{The structural branch and generation multiplicity}

The Standard-Model specialization has three distinct algebraic stages.
The tetrahedral multiplicity census and actual internal words determine a
complex seed algebra; its represented determinant structure or its anomaly-free
chiral tensor action determines the physical global group; that group's
matter representation generates the kinematic typing algebra on the common
Clifford-compatible carrier.  The external tetrahedral irreducible spaces
are not identified with colour and weak multiplicity spaces, and the seed
is not the commutant of the full matter-incidence action.

This also distinguishes the construction from the usual almost-commutative
spectral geometry, where spacetime spin geometry and a finite internal
geometry enter a spectral triple and its spectral action
\cite{ConnesGravityMatter1996,ChamseddineConnes1997,
ChamseddineConnesMarcolli2007,ChamseddineConnes2008,
vanDenDungenVanSuijlekom2012}.  The comparison is between resulting
structures, not an identification of the underlying finite geometries.

The determinant and anomaly routes agree on the global quotient but make
different source assumptions.  A determinant ray is phase blind; a fixed
tensor or a nonzero neutral--incidence cross block imposes its stabilizer.
The anomaly route instead assumes the displayed chiral tensor packet and
primitive charge normalization.  The kernel of its matter representation
then determines the faithful quotient.  Neither route infers a missing
matter representation from the seed algebra alone.

For generations, endpoint rank, charged replication, and frame synchronization
are separate conclusions.  A positive rank-three endpoint residual must
first occur faithfully in a complete charged multiplicity source.  Its
allocation short counts the excess beyond that image, while the inventory
and anomaly tests control all charged types.  Zero excess then proves three
complete charged copies.  Equal multiplicities already give an abstract
common factor; common ancestry additionally supplies a physical relative
frame.  Without a cross-component datum, the quark--lepton relative frame is
a presentation freedom, not an extra observable.  None of these statements
fixes numerical masses, mixing parameters, or couplings.


\subsection{Source identification and continuum scope}

The operational inverse problem precedes the commutant problem.  A classical
record may leave the acted-on algebra undetermined, whereas an appropriate
operator-valued complementary pullback reconstructs it.  Historical word
moments determine mixed source data only when the relevant operations are
already represented and word native.  The module-recognition criterion
prevents invariant multiplicities of a regular word space from being
mistaken for executable internal factors on the physical carrier.  Similarly,
unread spectator dynamics cannot supply active matter occurrence.

Renewal Geometry provides one upstream route through its predictive carrier,
history algebra, and word-Gram hierarchy
\cite{PelissierSpacetime2026,PelissierSpectral2026}.  The finite theorems
remain applicable to unrelated realizations that supply the same represented
input.  The common-source Schur quotient in Appendix~\ref{app:coupled-source}
separates internal variation directions from an already represented
gravitational bank, but it is not a variational action.  Nonlinear matter
stress convergence, control of physical variations and constraints, and
classical or quantum continuum closure are additional analytic problems.

\paragraph{Scope and outlook.}

The results established here are finite and structural.  The
double-centralizer theorem applies to general source-complete joint packets,
whereas the Standard-Model classification is obtained only on the
tetrahedrally equivariant branch considered above.  The finite criteria
identify multiplicity birth, incidence connectivity, determinant structure,
generation support, and rigidity of the resulting commutant, but they do not
imply that a broad class of microscopic dynamics must realize these
conditions.  Nor do they determine gauge couplings, Higgs parameters, Yukawa
eigenvalues, or renormalization-group data.

A separate analytic problem is required to pass from the stabilized finite
gravity--matter structure to a coupled continuum theory.  The cofinal
commutant theorem controls the protected multiplicity sector under
refinement, but does not itself supply the nonlinear compactness or
measure-theoretic control required for a classical or quantum continuum
limit.  Renewal Geometry
\cite{PelissierSpacetime2026,PelissierSpectral2026} provides one predictive
upstream realization of the finite data considered here, while the
action--multiplicity duality remains independent of that realization.

The central conclusion is therefore not that gauge symmetry is spacetime in
disguise, nor that the Standard Model is the commutant of a bare
four-dimensional Clifford algebra.  Rather, once spacetime action, typed
internal action, and physical matter incidence are represented on a common
finite carrier, the complete visible action and the matter multiplicity that
it cannot resolve become mutually determining operator-algebraic structures.
The structural Standard-Model branch gives one distinguished realization of
this general dual pair.

\section*{Data and formalization availability}
Selected finite algebraic components of the construction have been formalized in
\texttt{Lean}.  Source code is available at
\url{https://github.com/AI-MathPhys/renewal-geometry}.


\appendix


\section{Finite operator algebras, support-polar forms, and incidence minimality}
\label{app:commutant}

This appendix supplies the coordinate, support, and deletion details for
Section~\ref{sec:finite-duality}.  The finite bicommutant input is the
reciprocal matrix decomposition in Corollary~\ref{cor:reciprocal-wedderburn}:
taking the commutant twice returns the original unital finite-dimensional
$C^*$-algebra.


\subsection{Coordinates, order reversal, and typing refinement}
\label{app:coefficient-consequences}

\begin{corollary}[Canonicality under Wedderburn coordinates]
\label{cor:typing-coordinate-canonicality}
Fix the represented pair $(\A_{\rm Cl},\A_{\rm type})$ and let
$z_\lambda$ be the minimal central projections of
$C^*(\A_{\rm Cl},\A_{\rm type})$.  Any two decompositions of the form
\[
 z_\lambda\Hh\cong K_\lambda\otimes M_\lambda,
 \qquad
 z_\lambda\Hh\cong\widetilde K_\lambda\otimes\widetilde M_\lambda
\]
which realize the same represented external--typing algebra differ, after
relabeling by the same central projections, by unitary multiplicity
identifications $W_\lambda:M_\lambda\to\widetilde M_\lambda$.  For every
represented incidence edge $e:\lambda\to\mu$,
\[
 \widetilde{\mathfrak B}_e
 =W_\mu\mathfrak B_eW_\lambda^*,
\]
and, with $W=\bigoplus_\lambda W_\lambda$,
\begin{equation}
 \widetilde{\Oalg}_Y=W\Oalg_YW^*,
 \qquad
 \widetilde{\M}_{\rm type}=W\M_{\rm type}W^*.
 \label{eq:typing-coordinate-canonicality}
\end{equation}
Thus the typed dual pair is independent of matrix units, carrier bases and
Wedderburn coordinates.  The physical typing algebra itself, however, remains
represented input.
\end{corollary}

\begin{proof}
The projections $z_\lambda$ are intrinsic minimal central projections of the
represented algebra.  On each central block the commutant is a full matrix
algebra, and finite-dimensional uniqueness of its defining representation
gives the multiplicity unitary $W_\lambda$.  A change of carrier coordinates
is absorbed by the full family of carrier linear functionals used in the
definition of the slice space, while the multiplicity coordinates transform by
left--right unitary transport.  This gives the slice-space identity and hence
the two generated-algebra identities.
\end{proof}

\begin{proposition}[Order reversal and represented typing refinement]
\label{prop:typing-refinement}
For two incidence banks on the same typed multiplicity carrier,
\begin{equation}
 \Oalg_Y\subseteq\Oalg_{\widetilde Y}
 \quad\Longleftrightarrow\quad
 \M_{\rm type}(\widetilde Y)\subseteq\M_{\rm type}(Y).
 \label{eq:incidence-closure-galois}
\end{equation}
Thus equality of the closed coefficient-action algebras is equivalent to
equality of the exact multiplicity algebras.
For a fixed external factor and incidence bank $\bm Y$, let
$\A_{\rm type}^{(1)}\subseteq\A_{\rm type}^{(2)}\subseteq\A_{\rm Cl}'$,
and put $\A_i=C^*(\A_{\rm Cl},\A_{\rm type}^{(i)},\bm Y,\bm Y^*)$,
$\M_i=\A_i'$.  Then
\begin{equation}
 \A_1\subseteq\A_2,\qquad\M_2\subseteq\M_1,
 \label{eq:typing-refinement-monotone}
\end{equation}
and
\[
 \M_1=\M_2\quad\Longleftrightarrow\quad\A_1=\A_2
 \quad\Longleftrightarrow\quad\A_{\rm type}^{(2)}\subseteq\A_1.
\]
A typing refinement is therefore multiplicity-neutral exactly when its
added operations already lie in the generated action.
\end{proposition}
\begin{proof}
Taking commutants reverses inclusions, and finite bicommutant closure
reverses them back.  The typing statement additionally uses
$\A_2=C^*(\A_1,\A_{\rm type}^{(2)})$.
\end{proof}


\subsection{Coefficient banks and support refinements}
\label{app:coefficient-spaces}

\begin{lemma}[Multiplicity-one incidence factorization]
\label{lem:multiplicity-factorization}
If the prescribed carrier coefficient space of an edge is a line
$\C D_e$, with $D_e\ne0$, then the edge has a unique factorization, relative
to this choice of $D_e$,
\begin{equation}
 Y_e=D_e\otimes F_e,\qquad F_e:M_\lambda\to M_\mu.
 \label{eq:incidence-factorization}
\end{equation}
In particular this holds when a declared equivariance condition has a
one-dimensional carrier intertwiner space and acts trivially on the
multiplicity factors.  On this branch the equations in
\eqref{eq:typed-multiplicity} reduce to
$R_\mu F_e=F_eR_\lambda$ and $R_\lambda F_e^*=F_e^*R_\mu$.
\end{lemma}

\begin{proof}
The admitted edge space is
$\C D_e\otimes\mathcal B(M_\lambda,M_\mu)$; contraction by a linear
functional taking value one on $D_e$ determines $F_e$.  The general
coefficient-space theorem then gives the two displayed equations.
\end{proof}

\noindent\textit{Support refinements and canonical data.} The intrinsic coefficient spaces, their adjoint-closed quiver, and
$\Oalg_Y$ are independent of coefficient bases.  The support algebra of a
chosen coefficient list can change when that list is linearly recombined.
Its isotypic decomposition consequently gives a presentation of the same
commutant, not an additional basis-independent physical decomposition.
All such support-polar refinements recover the same $\Oalg_Y$ because
they generate exactly the original coefficient algebra.  When a support
bank is part of the represented data, its refined quiver is intrinsic to
that bank up to the usual unitary multiplicity identifications.


\subsection{Singular polar edges, leakage, and support certification}
\label{app:polar-details}

Let $F:E\to H$ have polar decomposition
\[
F=UP,
\qquad
P=(F^*F)^{1/2},
\qquad
p=U^*U=\supp P,
\qquad
q=UU^*=\supp(FF^*).
\]

\begin{lemma}[Singular polar-edge equivalence]
\label{lem:polar-edge}
For $R_E\in\mathcal B(E)$ and $R_H\in\mathcal B(H)$, the equations
\[
R_HF=FR_E,
\qquad
R_EF^*=F^*R_H
\]
are equivalent to
\begin{equation}
[R_E,P^2]=0,
\qquad
R_HU=UR_E,
\qquad
R_EU^*=U^*R_H.
\label{eq:polar-edge-equivalence}
\end{equation}
No invertibility assumption is required.
\end{lemma}

\begin{proof}
The two incidence equations imply
$F^*R_HF=F^*FR_E=R_EF^*F$, hence $[R_E,P^2]=0$.
Functional calculus gives commutation with $P$ and its support.
Multiplication by the Moore--Penrose inverse of $P$ on its support gives the
partial-isometry transport equations.  The converse follows by multiplying
the transport equations by $P$ and using $[R_E,P]=0$.
\end{proof}

\begin{proposition}[Exact singular polar-edge energy and support leakage]
\label{prop:singular-polar-energy}
For a nonzero coefficient map $F:E\to H$ with polar decomposition
$F=UP$, support projections $p=U^*U$, $q=UU^*$, and
$\mu=\lambda_{\min}(P|_{pE})>0$, define
\[
\mathcal E_F(R_E,R_H)
=\norm{R_HF-FR_E}_{\HS}^2
 +\norm{R_EF^*-F^*R_H}_{\HS}^2.
\]
Put
\[
 A=pR_Ep,
 \qquad
 B=U^*qR_HqU,
 \qquad
 S=\frac{A+B}{2},
 \qquad
 \Delta=B-A.
\]
Then
\begin{align}
\mathcal E_F(R_E,R_H)={}&
2\norm{[S,P]}_{\HS}^2
+\frac12\norm{\{\Delta,P\}}_{\HS}^2
\nonumber\\
&+\norm{PpR_E(1-p)}_{\HS}^2
 +\norm{(1-p)R_EpP}_{\HS}^2
\nonumber\\
&+\norm{PU^*qR_H(1-q)}_{\HS}^2
 +\norm{(1-q)R_HqUP}_{\HS}^2.
\label{eq:singular-polar-energy}
\end{align}
Consequently
\begin{equation}
\mathcal E_F(R_E,R_H)
\ge
\mu^2\left(
\norm{[p,R_E]}_{\HS}^2
+\norm{[q,R_H]}_{\HS}^2
+2\norm{U^*qR_HqU-pR_Ep}_{\HS}^2
\right).
\label{eq:singular-polar-support-bound}
\end{equation}
The energy vanishes exactly when $R_E$ and $R_H$ reduce the source and target
supports, their supported blocks are transported by $U$, and the supported
source block commutes with $P$.  The blocks on $\ker F$ and $\ker F^*$ remain
independent.
\end{proposition}

\begin{proof}
Relative to $E=pE\oplus(1-p)E$ and $H=qH\oplus(1-q)H$, the map $F$ has
only the block $UP:pE\to qH$.  Direct block multiplication gives the four
support-mixing terms above and two supported terms which, after transport by
$U$, are $BP-PA$ and $AP-PB$.  Since
\[
BP-PA=[S,P]+\tfrac12\{\Delta,P\},
\qquad
AP-PB=[S,P]-\tfrac12\{\Delta,P\},
\]
the parallelogram identity gives \eqref{eq:singular-polar-energy}.  In an
eigenbasis of $P|_{pE}$,
$\{\Delta,P\}_{ij}=(\lambda_i+\lambda_j)\Delta_{ij}$, so its contribution is
at least $2\mu^2\norm{\Delta}_{\HS}^2$.  Each support-mixing term is bounded
below by $\mu^2$ times the corresponding bare block norm, yielding
\eqref{eq:singular-polar-support-bound}.  Equality to zero forces all four
mixing blocks to vanish; positivity of $P|_{pE}$ makes the anticommutator map
invertible, so $\Delta=0$, and the remaining term gives $[A,P]=0$.
\end{proof}

\begin{proposition}[Certified support stratum from a singular-value gap]
\label{thm:certified-support-stratum}
Let $F:E\to H$ be one represented coefficient map and let
$\widehat F$ satisfy
$\|\widehat F-F\|_{\mathrm{op}}\le\varepsilon$.  Suppose an independent
source or structural premise gives $\operatorname{rank}F\le r$ and the
$r$th observed singular value obeys
\[
 \widehat\sigma_r>2\varepsilon.
\]
Then $\operatorname{rank}F=r$.  If
\[
 p=\operatorname{supp}(F^*F),\qquad
 q=\operatorname{supp}(FF^*)
\]
and $\widehat p_r,\widehat q_r$ are the right and left singular projectors of
$\widehat F$ for its $r$ largest singular values, then
\begin{equation}
 \|p-\widehat p_r\|_{\mathrm{op}},\ 
 \|q-\widehat q_r\|_{\mathrm{op}}
 \le
 \frac{\varepsilon}{\widehat\sigma_r-2\varepsilon}.
 \label{eq:certified-support-stratum}
\end{equation}
The same statement holds under the stronger-looking but often more directly
available hypothesis
$\|\widehat F-F\|_{\HS}\le\varepsilon$.
\end{proposition}

\begin{proof}
Singular-value perturbation gives
$\sigma_r(F)\ge\widehat\sigma_r-\varepsilon>\varepsilon$, so the upper-rank
premise forces rank exactly $r$.  Put $s=\sigma_r(F)$.  If
$y\in\ker F$, then $\|\widehat Fy\|\le\varepsilon\|y\|$.  Orthogonality of
the left-singular images associated with $\widehat p_r$ and
$I-\widehat p_r$ gives
\[
 (s-\varepsilon)\|\widehat p_ry\|
 \le\|\widehat F\widehat p_ry\|
 \le\|\widehat Fy\|.
\]
Hence
$\|\widehat p_r(I-p)\|\le\varepsilon/(s-\varepsilon)$.
The projectors $p$ and $\widehat p_r$ have equal finite rank, so this is their
operator-norm projector distance.  Since
$s-\varepsilon\ge\widehat\sigma_r-2\varepsilon$, the source-support bound
follows; apply the same argument to $F^*$ for the target support.  Finally
$\|X\|_{\rm op}\le\|X\|_{\HS}$ gives the HS-error version.
\end{proof}

\noindent\textit{Support certification is one-sided without a rank premise.} An observed singular value above the error floor certifies a genuine support
direction, but a small observed tail cannot certify exact absence.  The upper
rank/source premise in Proposition~\ref{thm:certified-support-stratum} is therefore
essential; it is the support-level analogue of the one-sided coefficient-rank
firewall in Corollary~\ref{cor:incidence-one-sided-rank}.


\subsection{Noncommuting supports and the refined quiver}

Support projections themselves can fail to commute.  The correct residual
description is then a finite operator-space quiver, in the standard
representation-theoretic sense of a quiver with linear arrow spaces
\cite{Schiffler2014}.

\begin{definition}[Multiplicity quiver]
\label{def:multiplicity-quiver}
Let
\[
\mathcal S_F
=
C^*(Z_\lambda,\supp(F_e^*F_e),\supp(F_eF_e^*):\lambda,e)
\]
and write its canonical finite isotypic decomposition as
\[
\mathscr M
\cong
\bigoplus_{a\in\mathfrak A}V_a\otimes N_a,
\qquad
\mathcal S_F
=
\bigoplus_a\mathcal B(V_a)\otimes I_{N_a}.
\]
Let $\bm T$ be any finite $*$-closed family of further reconstructed support
metrics, polar arrows, cycle products, or action generators.  For a block
$T_{ba}:V_a\otimes N_a\to V_b\otimes N_b$, define
\[
\mathfrak B_{ba}(T)
=
\{(\varphi\otimes\id)(T_{ba}):
\varphi\in\mathcal B(V_a,V_b)^*\}
\subseteq\mathcal B(N_a,N_b).
\]
The multiplicity quiver $\mathfrak Q_F(\bm T)$ has vertex spaces $N_a$ and
arrow spaces
\[
\mathfrak B_{ba}
=
\Span_{T\in\bm T}\mathfrak B_{ba}(T).
\]
For this fixed support bank its endomorphism algebra is
\[
\End\mathfrak Q_F
=
\left\{(X_a)_a:
X_bB=BX_a
\ \text{for every }B\in\mathfrak B_{ba}\right\}.
\]
\end{definition}

\begin{proposition}[Noncommuting-support quiver theorem]
\label{thm:quiver}
With the preceding notation,
\begin{equation}
C^*(\mathcal S_F,\bm T)'
\cong
\End\mathfrak Q_F(\bm T).
\label{eq:quiver-commutant}
\end{equation}
If
\[
(T_j)_{ba}
=
\sum_\alpha E_\alpha^{ba}\otimes B_{j,\alpha}^{ba}
\]
in Hilbert--Schmidt orthonormal bases of
$\mathcal B(V_a,V_b)$, then the ambient commutator form restricts to
\begin{equation}
q_{\bm T}[(X_a)]
=
\sum_{j,b,a,\alpha}
\norm{B_{j,\alpha}^{ba}X_a-X_bB_{j,\alpha}^{ba}}_{\HS}^2,
\label{eq:quiver-form}
\end{equation}
whose kernel is exactly $\End\mathfrak Q_F$.
\end{proposition}

\begin{proof}
An operator commuting with $\mathcal S_F$ has the form
$\bigoplus_a I_{V_a}\otimes X_a$.  Expanding its commutator with every
$T_j$ in the Hilbert--Schmidt basis gives precisely the quiver intertwining
relations.  Orthogonality gives \eqref{eq:quiver-form}.
\end{proof}

For a fixed support algebra, changing Hilbert--Schmidt carrier bases
unitarily recombines the coefficient families and leaves each arrow space
unchanged.  The refined quiver is consequently independent of those bases,
up to multiplicity unitaries and block permutations.  Replacing the support
bank can change the refinement, but not the coefficient-action algebra or
its commutant.


\subsection{The connected invertible branch}

When a chosen coefficient bank has invertible arrows, the duality has a
particularly geometric form.  For a multiplicity-one physical edge bank these
are exactly its residue maps; general physical incidences may contribute
several coefficient arrows.

\begin{proposition}[Polar metric--holonomy commutant]
\label{thm:polar-holonomy}
Assume $\Gamma$ is connected, all multiplicity spaces have common dimension
$g$, and every edge map is invertible:
\[
F_e=U_eP_e,
\qquad
U_e\in\U(g),
\qquad
P_e\succ0.
\]
Choose a root $o$ and a spanning tree.  Let
$Q_\lambda:M_o\to M_\lambda$ be the ordered product of the $U_e$ along the
tree path.  For $e:\lambda\to\mu$, define on the root fibre
\[
K_e=Q_\lambda^*P_e^2Q_\lambda,
\qquad
W_e=Q_\mu^*U_eQ_\lambda,
\]
and put
\begin{equation}
\Oalg_F
=
C^*(K_e,W_e:e\in E)\subseteq M_g(\C).
\label{eq:root-holonomy-algebra}
\end{equation}
Then
\begin{equation}
\M_{\rm type}^{(o)}
=
\Oalg_F',
\qquad
\Oalg_F
=
(\M_{\rm type}^{(o)})'.
\label{eq:root-holonomy-duality}
\end{equation}
The multiplicity family associated with
$R\in\Oalg_F'$ is
\[
R_\lambda=Q_\lambda RQ_\lambda^*.
\]
\end{proposition}

\begin{proof}
The polar-edge equations give
$[R_\lambda,P_e^2]=0$ and
$R_\mu=U_eR_\lambda U_e^*$.
Along the spanning tree these force
$R_\lambda=Q_\lambda RQ_\lambda^*$ from the root value $R=R_o$.
The metric condition becomes $[R,K_e]=0$, while a non-tree edge gives the
cycle-consistency condition $[R,W_e]=0$.  Thus the residual family is
equivalent to one root operator commuting with the algebra
\eqref{eq:root-holonomy-algebra}.  Finite bicommutant closure gives the reverse
equality.
\end{proof}

\begin{corollary}[Flat and irreducible extremes]
\label{cor:polar-extremes}
Under Proposition~\ref{thm:polar-holonomy}:
\begin{enumerate}[label=\textup{(\roman*)},leftmargin=2.8em]
\item $\M_{\rm type}^{(o)}=M_g(\C)$ if and only if every $K_e$ and $W_e$ is
      scalar;
\item $\M_{\rm type}^{(o)}=\C I$ if and only if $\Oalg_F=M_g(\C)$;
\item if every $F_e$ is unitary, then $K_e=I$ and the obstruction reduces to
      relative cycle holonomy;
\item even on a tree, a non-scalar positive polar metric can reduce the
      residual multiplicity algebra.
\end{enumerate}
\end{corollary}


\subsection{Determinant supports and simultaneous deletion}
\label{app:incidence-details}

The rank theorem in Section~\ref{sec:incidence-skeleton} has a determinant
refinement useful for finite audits.  Keep the notation there and introduce
nonnegative incidence weights $w=(w_e)_{e\in E}$:
\begin{equation}
 G(w)=\sum_{e\in E}w_eG_e,
 \qquad P_E(w)=\det G(w).
 \label{eq:weighted-incidence-Gram}
\end{equation}

\begin{proposition}[Positive incidence determinant and sufficient supports]
\label{thm:incidence-determinant}
The polynomial $P_E$ is homogeneous of degree $n$ and has nonnegative real
coefficients.  For any $S\subseteq E$, if $w_e>0$ on $S$ and $w_e=0$ off
$S$, then
\begin{equation}
 P_E(w)>0\quad\Longleftrightarrow\quad f(S)=n
 \quad\Longleftrightarrow\quad \M(S)=\M(E).
 \label{eq:incidence-face-exactness}
\end{equation}
If
$\kappa_e=n-f(E\setminus\{e\})$, then the smallest power of $w_e$ occurring
with nonzero coefficient in $P_E$ is exactly $\kappa_e$.  Thus $e$ is
essential if and only if $w_e$ divides $P_E$.
\end{proposition}

\begin{proof}
Stack the row blocks $\sqrt{w_e}D_e$.  Cauchy--Binet expresses
$\det G(w)$ as a sum of squared absolute values of $n\times n$ row minors,
each multiplied by the monomial recording how many rows were selected from
each edge block.  All coefficients are therefore nonnegative.  A point on a
coordinate face has positive determinant exactly when the row spaces of the
retained blocks span $V$, which is $f(S)=n$.  For the last assertion, any row
basis uses at most $f(E\setminus\{e\})$ rows from the other blocks and hence
at least $\kappa_e$ rows from block $e$; extending a basis of the other-block
span using rows from $e$ attains this number.
\end{proof}

For quantitative deletion, put $G=G_E\succ0$ and
\[
 B_e=G^{-1/2}G_eG^{-1/2},\qquad \sum_eB_e=I,
 \qquad B_R=\sum_{e\in R}B_e.
\]

\begin{proposition}[Exact simultaneous-deletion margin]
\label{prop:deletion-margin}
For every $R\subseteq E$,
\begin{equation}
 G_{E\setminus R}=G^{1/2}(I-B_R)G^{1/2},
 \label{eq:deletion-factor}
\end{equation}
and therefore
\begin{equation}
 \M(E\setminus R)=\M(E)
 \quad\Longleftrightarrow\quad \|B_R\|_{\rm op}<1.
 \label{eq:deletion-margin}
\end{equation}
If $\beta_R=\|B_R\|_{\rm op}<1$, then
\[
 G_{E\setminus R}\succeq(1-\beta_R)G,
 \qquad
 \lambda_{\min}(G_{E\setminus R})
 \ge(1-\beta_R)\lambda_{\min}(G).
\]
\end{proposition}

\begin{proof}
The factorization follows from the definitions.  Since
$0\preceq B_R\preceq I$, $I-B_R$ is positive definite exactly when the
largest eigenvalue of $B_R$ is smaller than one.  Congruence by $G^{1/2}$ preserves
nullity and positivity, giving both the commutant criterion and the lower
bound.
\end{proof}

\begin{example}[Minimal incidence skeletons need not form matroid bases]
\label{ex:nonmatroid-skeleton}
On $\C^3$ with scalar base algebra let
\[
 Y_a=E_{12}+E_{23},\qquad Y_b=E_{12},\qquad Y_c=E_{23}.
\]
The singleton $\{a\}$ and the pair $\{b,c\}$ both generate $M_3(\C)$ after
adjoining adjoints, and hence both have scalar commutant.  Each is
inclusion-minimal: neither $Y_b$ nor $Y_c$ alone has scalar commutant.  Thus
minimal sufficient incidence sets can have unequal cardinality, so the
incidence labels do not carry a matroid-basis structure in general.
\end{example}


\section{Source identification and historical compilation}
\label{app:source-identifiability}

The finite duality presupposes represented operators on the acted-on
carrier.  This appendix distinguishes that input from coarser record data,
gives a coherent reconstruction route, and identifies historical operations
that can be evaluated through existing word moments.


\subsection{System-algebra identification}
\label{sec:source-identifiability}

The source-complete packet is a represented system object.  It cannot in
general be reconstructed from a classical outcome record or from the algebra
of an auxiliary environment alone.  The following finite statement separates
what such coarse data determine from what is needed for the duality.

\begin{proposition}[Classical record data do not determine the system algebra]
\label{thm:record-nonidentifiability}
There exist two finite normalized instruments on the same input and output
Hilbert spaces with the same resolved classical record channel at every word
length, the same one-step effects, the same one-letter Hilbert--Schmidt Gram,
and the same nonzero coarse Choi spectrum, but with nonisomorphic represented
system coefficient algebras and commutants.  Consequently no theorem using
only those record/environment invariants can identify the joint action algebra
of Definition~\ref{def:action-algebra}.
\end{proposition}

\begin{proof}
The explicit pair in Proposition~\ref{prop:record-witness} supplies the witness.
One coefficient family is simultaneously diagonal and generates
$\C^{20}$; the other generates $M_4(\C)\otimes I_5$.  Each resolved branch is
a common scalar multiple of a unitary, so every uninterrupted classical word
has the same effect in the two realizations, while the coefficient algebras
and their commutants are different.  The displayed Gram and Choi identities in
Proposition~\ref{prop:record-witness} complete the finite witness.
\end{proof}

Let $V:\Hh\to\Hh\otimes\mathcal E$ be an isometry and choose an orthonormal
environment basis $(|\nu\rangle)$.  Write
\[
 V=\sum_\nu C_\nu\otimes|\nu\rangle,
 \qquad
 C(z)=(I\otimes\langle z|)V,
\]
and let $\mathcal P_*$ be the complementary channel.  Its Heisenberg pullback
obeys
\begin{equation}
 \mathcal P^*(|z\rangle\langle w|)=C(z)^*C(w).
 \label{eq:complementary-pullback}
\end{equation}

\begin{proposition}[Identity-pivot reconstruction of the system action]
\label{thm:pullback-reconstruction}
Suppose the occurring coefficient algebra satisfies
$\mathcal A_{\rm sys}=C^*(C_\nu:\nu)=C^*(A_1,\ldots,A_s)$ and there are environment vectors
$z_0,z_1,\ldots,z_s$ with
\[
 C(z_0)=\alpha I,\qquad C(z_j)=\beta_jA_j,
 \qquad \alpha\beta_j\ne0.
\]
Then
\begin{equation}
 C^*\!\left(\mathcal P^*(\mathcal B(\mathcal E))\right)
 =\mathcal A_{\rm sys},
 \qquad
 [\mathcal P^*(\mathcal B(\mathcal E))]'=\mathcal A_{\rm sys}'.
 \label{eq:pullback-reconstruction}
\end{equation}
In particular, a complete operator-valued complementary pullback recovers the
represented system action whenever the physical source contains such an
identity pivot; no fitted remainder algebra is required.
\end{proposition}

\begin{proof}
Every pullback $C(z)^*C(w)$ belongs to $C^*(C_\nu)$ and hence to
$\mathcal A_{\rm sys}$.  Conversely,
\[
 \mathcal P^*(|z_0\rangle\langle z_j|)
 =\overline\alpha\,\beta_j A_j,
\]
so every generator $A_j$ lies in the pullback algebra.  Equality of the
commutants follows immediately.
\end{proof}

\begin{corollary}[Finite coherent system export]
\label{cor:finite-system-export}
For the normalized six-edge source class of
Appendix~\ref{app:source-identifiability}, ten fixed Hermitian environment
quadratures suffice to reconstruct all six edge coefficients on the acted-on
carrier.  With Hilbert--Schmidt coefficient error $e_K$ as defined in
\eqref{eq:appendix-reconstruction-error}, the associated star-closed
commutator map changes in operator norm by at most $2\sqrt2\,e_K$.  Hence an
independently verified protected commutant complement with measured least
singular value $\widehat s>2\sqrt2\,e_K$ has no additional commutant direction.
\end{corollary}

\noindent\textit{Logical role of source identification.} Proposition~\ref{thm:record-nonidentifiability} is a stopping theorem, not a
weakening of the double-centralizer result.  Once the represented system
operators are supplied or reconstructed, Theorem~\ref{thm:main-duality} is
exact.  What cannot be inferred from coarse record data is which system
operators occurred on the common carrier in the first place.


\subsection{A finite nonidentifiability pair and coherent reconstruction}

Let $H_E$ be a finite edge-system space, let
$Q_a$ ($1\le a\le4$) be nonzero orthogonal projections on a record factor
$H_R$ with $\sum_aQ_a=I$, and let $K_1,\ldots,K_6\in\mathcal B(H_E)$ satisfy
\begin{equation}
 \sum_{e=1}^6K_e^*K_e=I,
 \qquad
 \sum_{e=1}^6K_e=\sqrt2\,I,
 \qquad
 K_1,\ldots,K_6\ \text{linearly independent}.
 \label{eq:six-edge-normalization}
\end{equation}
For $0<\vartheta<1$ set
\[
 A_{e,a}=K_e\otimes Q_a,
 \qquad L_0=\sqrt{1-\vartheta}\,I,
 \qquad L_{e,a}=\sqrt\vartheta\,A_{e,a}.
\]
These operators define a normalized resolved instrument.

Choose a real $6\times5$ matrix $O$ with
\begin{equation}
 O^{\mathsf T}O=I_5,
 \qquad O^{\mathsf T}\mathbf1=0,
 \qquad OO^{\mathsf T}=I_6-\frac16\mathbf1\mathbf1^{\mathsf T}.
 \label{eq:helmert-frame}
\end{equation}

\begin{proposition}[Five-Kraus freedom behind the six-edge normalization]
\label{prop:five-kraus-freedom}
Equation~\eqref{eq:six-edge-normalization} holds if and only if there are
five matrices $B_1,\ldots,B_5$ with
$\sum_\mu B_\mu^*B_\mu=I$ and
$I,B_1,\ldots,B_5$ linearly independent such that
\begin{equation}
 K_e=\frac1{\sqrt{18}}I
       +\sqrt{\frac23}\sum_{\mu=1}^5O_{e\mu}B_\mu.
 \label{eq:six-edge-normal-form}
\end{equation}
Consequently the coarse edge channel has the form
\[
 \mathcal E_* =\frac13\id+\frac23\Psi_*,
 \qquad \Psi_*(\rho)=\sum_\mu B_\mu\rho B_\mu^*,
\]
and $C^*(K_e:e)=C^*(B_\mu:\mu)$.
\end{proposition}

\begin{proof}
Apply the orthogonal coefficient change whose first column is
$\mathbf1/\sqrt6$ and whose remaining columns are $O$.  The coherent-sum
constraint fixes the first transformed coefficient to $I/\sqrt3$; the other
five, $D_\mu=\sum_eO_{e\mu}K_e$, satisfy
$\sum_\mu D_\mu^*D_\mu=2I/3$.  Put
$B_\mu=\sqrt{3/2}\,D_\mu$ and invert the orthogonal change.
\end{proof}

\subsubsection{An exact nonidentifiability pair}

For this witness take $H_R=\C^4$ and $Q_a=|a\rangle\langle a|$.
Let $\mathscr S$ be the twenty three-element subsets of
$\{1,\ldots,6\}$.  On $\C^{\mathscr S}$ define commuting Hermitian
involutions
\[
 H_e^{\rm c}|S\rangle=(2\mathbf1_{e\in S}-1)|S\rangle.
\]
Let $\Gamma_1,\ldots,\Gamma_5$ be five anticommuting Hermitian Clifford
matrices on $\C^4$, and on $\C^4\otimes\C^5$ set
\[
 H_e^{\rm q}=\sqrt{\frac65}\sum_{\mu=1}^5O_{e\mu}\Gamma_\mu\otimes I_5.
\]
For $b\in\{{\rm c},{\rm q}\}$ put
\begin{equation}
 K_e^b=\frac1{\sqrt{18}}I_{20}+\frac{i}{3}H_e^b.
 \label{eq:two-system-completions}
\end{equation}

\begin{proposition}[Record-complete nonidentifiability witness]
\label{prop:record-witness}
Both lists in \eqref{eq:two-system-completions} satisfy
\eqref{eq:six-edge-normalization}, and
\begin{equation}
 (K_e^b)^*K_e^b=I_{20}/6,
 \qquad
 \Tr[(K_e^b)^*K_f^b]
 =\begin{cases}10/3,&e=f,\\2/3,&e\ne f.\end{cases}
 \label{eq:matching-record-data}
\end{equation}
Yet
\begin{equation}
 C^*(K_e^{\rm c})=\C^{20},
 \qquad
 C^*(K_e^{\rm q})=M_4(\C)\otimes I_5.
 \label{eq:inequivalent-system-algebras}
\end{equation}
For the resolved instrument above, every uninterrupted word in the original
classical alphabet has the same effect in the two realizations.  A word with
$m\ge1$ accepted letters, $n_0$ no-response letters and one common accepted
record type $a$ has effect
\begin{equation}
 (1-\vartheta)^{n_0}(\vartheta/6)^m I_{20}\otimes Q_a;
 \label{eq:record-word-effect}
\end{equation}
if accepted record types disagree the effect is zero.  Moreover the two one-step opportunity channels have the same nonzero Choi
spectrum,
\begin{equation}
 \spec_+(J_{\rm opp})=
 \left\{
 \left[\frac{20}{3}(12-11\vartheta)\right]^{(1)},
 \left(\frac{20\vartheta}{3}\right)^{(3)},
 \left(\frac{8\vartheta}{3}\right)^{(20)}
 \right\}.
 \label{eq:matching-Choi-spectrum}
\end{equation}
Thus the complete classical record channel at every word length is identical
although the system algebras differ.
\end{proposition}

\begin{proof}
For the commuting family, counting three-subsets gives zero mean and pair
correlation $-1/5$.  For the Clifford family the row vectors of
$\sqrt{6/5}\,O$ are unit vectors with the same pair correlations, and
Clifford anticommutation gives $(H_e^{\rm q})^2=I$.  Hence both families
satisfy the same normalization and Gram relations, and each $\sqrt6K_e^b$
is unitary.  The twenty sign patterns distinguish all basis vectors in the
commuting realization, giving $\C^{20}$.  The five Clifford generators span
and generate $M_4(\C)\otimes I_5$.  Finally, every classical word effect is
$L_w^*L_w$; because each accepted system factor is $6^{-1/2}$ times a
unitary and the record projectors are orthogonal, this gives
\eqref{eq:record-word-effect} in both realizations.

For the one-step Choi spectrum, the Gram of the twenty-four accepted Choi
vectors is $G_{24}=G_6\otimes I_4$ with
$G_6=(8/3)I_6+(2/3)\mathbf1\mathbf1^{\mathsf T}$.  It has eigenvalue
$20/3$ on the uniform edge direction with four record directions and
$8/3$ on the twenty edge-contrast directions.  The coefficient matrix
$M_\vartheta=\vartheta I+(1-\vartheta)cc^*$ acts by
$12-11\vartheta$ on the global uniform direction and by $\vartheta$ on its
orthogonal complement.  The nonzero Choi eigenvalues are therefore those of
$M_\vartheta^{1/2}G_{24}M_\vartheta^{1/2}$, giving
\eqref{eq:matching-Choi-spectrum}.
\end{proof}

\subsubsection{Ten coherent pullbacks reconstruct the six coefficients}

Enumerate the twenty-four accepted operators $A_{e,a}$ as $A_j$ and let
$\mathsf B$ synthesize their Choi vectors $|A_j\rangle\!\rangle$.  In
$\C^{24}$ put
\[
 c=\mathbf1/\sqrt2,
 \qquad
 M_\vartheta=\vartheta I+(1-\vartheta)cc^*.
\]
Let $C_\nu$ be the matrices synthesized by the columns of
$\mathsf B M_\vartheta^{1/2}$ and define the minimal Stinespring isometry
\[
 V_\vartheta=\sum_{\nu=1}^{24}C_\nu\otimes|\nu\rangle,
 \qquad
 \mathcal P_{\vartheta,*}(\rho)
 =\Tr_{\rm sys}(V_\vartheta\rho V_\vartheta^*).
\]
For $z\in\C^{24}$ write
$C(z)=(I\otimes\langle z|)V_\vartheta$.  Then
$\mathcal P_\vartheta^*(|z\rangle\langle w|)=C(z)^*C(w)$.

For the twenty-four coefficient labels $(e,a)$ define
\[
 u_{e,a}=1/\sqrt{24},
 \qquad (r_\mu)_{e,a}=O_{e\mu}/2,
 \qquad \langle u,r_\mu\rangle=0,
\]
and the fixed Hermitian environment quadratures
\begin{equation}
 X_\mu=|u\rangle\langle r_\mu|+|r_\mu\rangle\langle u|,
 \qquad
 Y_\mu=-i|u\rangle\langle r_\mu|+i|r_\mu\rangle\langle u|.
 \label{eq:ten-quadratures}
\end{equation}
Let $\mathcal P_\vartheta^*$ denote the Heisenberg complementary pullback of
the normalized coefficient frame and write
$\mathsf X_\mu=\mathcal P_\vartheta^*(X_\mu)$,
$\mathsf Y_\mu=\mathcal P_\vartheta^*(Y_\mu)$.  Set
$h=12-11\vartheta$ and $D_\mu=\sum_eO_{e\mu}K_e$.

\begin{proposition}[Explicit ten-panel reconstruction]
\label{prop:ten-panel-reconstruction}
The ten Hermitian operator-valued pullbacks determine all six system
coefficients through
\begin{equation}
 D_\mu\otimes I_{H_R}
 =\frac{2\sqrt3}{\sqrt{\vartheta h}}
   (\mathsf X_\mu+i\mathsf Y_\mu),
 \qquad
 K_e=\frac{I}{\sqrt{18}}+\sum_\mu O_{e\mu}D_\mu.
 \label{eq:ten-panel-reconstruction}
\end{equation}
Thus this finite operator-valued export determines the represented system
algebra and commutant.
\end{proposition}

\begin{proof}
In the normalized coefficient frame the uniform environment vector is an
eigenvector of the coefficient Gram with eigenvalue $h$, while each
$r_\mu$ has eigenvalue $\vartheta$.  Hence
$C(u)=\sqrt{h/12}\,I$ and
$C(r_\mu)=\sqrt\vartheta D_\mu\otimes I_{H_R}/2$.
Since $X_\mu+iY_\mu=2|u\rangle\langle r_\mu|$,
Equation~\eqref{eq:complementary-pullback} gives the first reconstruction
formula; orthogonality of $O$ gives the second.
\end{proof}

\begin{proposition}[Robust commutant stopping test]
\label{prop:reconstruction-error}
If the ten pullbacks have Hilbert--Schmidt errors
$\epsilon_{X,\mu}$ and $\epsilon_{Y,\mu}$, then
\begin{equation}
 \left(\sum_e\|\widehat K_e-K_e\|_{\HS}^2\right)^{1/2}
 \le e_K:=\frac{2\sqrt3}{\sqrt{\vartheta h}}
 \left[\sum_\mu(\epsilon_{X,\mu}+\epsilon_{Y,\mu})^2\right]^{1/2}.
 \label{eq:appendix-reconstruction-error}
\end{equation}
For the star-closed stacked commutator map,
\begin{equation}
 \|\widehat{\mathscr D}-\mathscr D\|_{2\to2}\le2\sqrt2\,e_K.
 \label{eq:appendix-commutator-error}
\end{equation}
Consequently a measured least singular value
$\widehat s>2\sqrt2\,e_K$ on an independently verified protected complement
certifies absence of any further commutant direction there.
\end{proposition}

\begin{proof}
Equation~\eqref{eq:ten-panel-reconstruction} and the triangle inequality give
the $D_\mu$ errors.  Orthogonality of $O$ preserves the sum of squared
coefficient errors.  For unit Hilbert--Schmidt $X$, each coefficient error
changes a commutator by at most twice its operator norm; retaining both a
coefficient and its adjoint gives the factor $2\sqrt2$.  The last statement
is the standard least-singular-value perturbation inequality.
\end{proof}

The ten-panel theorem is a sufficient finite tomography statement, not an
assertion that these coherent quadratures are automatically present in every
microscopic source.  If only the classical resolved record is supplied,
Proposition~\ref{prop:record-witness} shows that the acted-on system algebra
can remain underdetermined.  Conversely, once the operator-valued pullbacks
are physically available on a common frame, the reconstruction is exact and
feeds directly into the commutant audit of the main text.


\subsection{Historical operations on the regular word carrier}
\label{sec:historical-operation-recognition}

The preceding reconstruction problem has a second layer.  Even after a finite
history algebra has been reconstructed, an invariant endomorphism of a word
Hilbert space need not be an operation induced by the historical system
algebra.  The following module criterion separates these notions.

Let $\A\subseteq\mathcal B(\Hh)$ be a finite unital $C^*$-algebra generated by
the represented historical operations on one exposed carrier, and let a finite
group $G$ act by unitaries $u_g$ with $u_g\A u_g^*=\A$.  Write
$\alpha_g(a)=u_gau_g^*$ and
\[
 \A\cong\bigoplus_\beta M_{n_\beta}(\C),\qquad
 D=\dim_\C\A=\sum_\beta n_\beta^2,
\]
with normalized regular trace
\[
 \tau_{\rm reg}(a)=D^{-1}\sum_\beta n_\beta\Tr(a_\beta).
\]
On $\mathscr W=L^2(\A,\tau_{\rm reg})$ let
$L_ax=ax$, $R_bx=xb$ and $U_gx=\alpha_g(x)$.

\begin{proposition}[Word-module recognition of historical internal actions]
\label{thm:word-module-recognition}
The operations on the regular word carrier induced by invariant historical
system elements are exactly
\begin{equation}
 R(\A)'\cap U(G)'=L(\A^G),
 \qquad
 \A^G=\A\cap\{u_g:g\in G\}'.
 \label{eq:word-module-recognition}
\end{equation}
Equivalently, $T\in\End(\mathscr W)$ is left multiplication by a unique
invariant historical element if and only if it commutes with every right
multiplication and every $U_g$; the coefficient is $T1$.
\end{proposition}

\begin{proof}
Every $L_a$ commutes with every $R_b$.  Conversely, if $[T,R_b]=0$ for all
$b$, then $T(b)=T(R_b1)=R_bT1=(T1)b$, hence $T=L_{T1}$.  Moreover
$U_gL_aU_g^{-1}=L_{\alpha_g(a)}$, so faithfulness of left multiplication gives
$[U_g,L_a]=0$ exactly when $a\in\A^G$.
\end{proof}

\begin{proposition}[Exact module and symmetry defects]
\label{prop:word-module-defect}
Let $(e_\nu)_{\nu=1}^D$ be a $\tau_{\rm reg}$-orthonormal basis of $\A$ and
put $a=T1$.  Define
\[
 \delta_R(T)^2=\sum_{\nu=1}^D\norm{[T,R_{e_\nu}]1}_{\tau}^2,
 \qquad
 \delta_G(T)^2=\frac1{|G|}\sum_g\norm{[U_g,T]1}_{\tau}^2.
\]
If $E_G(a)=|G|^{-1}\sum_g\alpha_g(a)$, then
\begin{align}
 \norm{T-L_a}_{\HS(\mathscr W)}^2&=\delta_R(T)^2,
 \label{eq:word-module-defect}\\
 \norm{T-L_{E_G(a)}}_{\HS(\mathscr W)}
 &\le \delta_R(T)+\sqrt{D/2}\,\delta_G(T).
 \label{eq:word-module-invariant-defect}
\end{align}
Thus both defects vanish exactly on the algebra in
\eqref{eq:word-module-recognition}.
\end{proposition}

\begin{proof}
For each basis element,
$(T-L_a)e_\nu=T(R_{e_\nu}1)-R_{e_\nu}T1=[T,R_{e_\nu}]1$; summation gives
\eqref{eq:word-module-defect}.  Since $U_g1=1$,
$[U_g,T]1=\alpha_g(a)-a$.  Orthogonal group averaging gives
$|G|^{-1}\sum_g\norm{\alpha_g(a)-a}_\tau^2=2\norm{a-E_G(a)}_\tau^2$.
Finally $\norm{L_b}_{\HS(\mathscr W)}^2=D\norm b_\tau^2$, and the triangle
inequality proves \eqref{eq:word-module-invariant-defect}.
\end{proof}

\begin{corollary}[Word-space multiplicity firewall]
\label{cor:word-multiplicity-firewall}
Multiplicity of the conjugation representation on $L^2(\A,\tau_{\rm reg})$
need not equal multiplicity on the exposed physical carrier and cannot by
itself be interpreted as an executable internal factor.  Any such inference
must also pass Proposition~\ref{thm:word-module-recognition}.  In particular, the
multiplicity-birth analysis of Section~\ref{sec:multiplicity-birth} always
refers to an actual reachable--observable physical carrier, not merely to the
regular word representation.  Appendix~\ref{app:historical-source-selection}
contains an explicit twelve-root example with an invariant $M_7(\C)$ word-space
block but no physical internal $M_3(\C)$ factor.
\end{corollary}

We next identify a class of mixed source quantities that require no independent
measurement once the relevant typing operations are already historical.

\begin{definition}[Word-native linear shadow]
\label{def:word-native-shadow}
A linear map $\mathscr L:\A\to\A$ is \emph{word native} if
\begin{equation}
 \mathscr L(X)=\sum_{\alpha=1}^m c_\alpha A_\alpha X B_\alpha,
 \label{eq:word-native-shadow}
\end{equation}
where $A_\alpha,B_\alpha\in\A$ are reconstructed historical-algebra elements
and the coefficients $c_\alpha$ are fixed by the declared shadow before the
target conclusion is evaluated.  Finite compressions, conjugations, finite
group averages and finite functional-calculus operations are included whenever
their defining elements belong to $\A$.
\end{definition}

\begin{proposition}[Inserted word-moment compiler]
\label{thm:inserted-word-moments}
Let $x,y,z\in\A$ be represented words and let $\mathscr L,\mathscr M$ be
word-native maps, with
$\mathscr M(Z)=\sum_\beta d_\beta C_\beta ZD_\beta$.  Then
\begin{align}
 \langle y,\mathscr L(x)\rangle_{\HS}
 &=\sum_\alpha c_\alpha\,
   \tau_{\rm reg}(y^*A_\alpha xB_\alpha),
 \label{eq:inserted-one-moment}\\
 \langle\mathscr L(x),\mathscr M(z)\rangle_{\HS}
 &=\sum_{\alpha,\beta}\overline{c_\alpha}d_\beta\,
   \tau_{\rm reg}(B_\alpha^*x^*A_\alpha^*C_\beta zD_\beta).
 \label{eq:inserted-two-moments}
\end{align}
Hence every ordinary--shadow and shadow--shadow Gram block is a finite linear
combination of ordinary Grand-Tensor word moments.  If the historical word
algebra is already flat, the reconstructed multiplication table reduces all
inserted products and no deeper physical histories are required.
\end{proposition}

\begin{proof}
The formulas are the Hilbert--Schmidt inner product followed by substitution of
\eqref{eq:word-native-shadow}.  Cyclicity of $\tau_{\rm reg}$ rewrites each
summand as the inner product of two represented algebra elements.  At flatness
the finite word span is a $*$-algebra, so every inserted product reduces in its
known multiplication table.
\end{proof}

\begin{corollary}[Table-native ancestry short]
\label{cor:table-native-ancestry}
Let $W_r$ synthesize a flat historical word bank, let
$S=W_rB_S$, $T=\mathscr L W_rB_T$, and define
\[
 K^{00}=W_r^*W_r,\qquad
 K^{0L}=W_r^*\mathscr LW_r,\qquad
 K^{LL}=W_r^*\mathscr L^*\mathscr LW_r.
\]
Then
\begin{equation}
 G=B_S^*K^{00}B_S,
 \quad C=B_S^*K^{0L}B_T,
 \quad H=B_T^*K^{LL}B_T,
 \label{eq:table-native-GCH}
\end{equation}
and the ancestry short
\begin{equation}
 R_{T\mid S}=H-C^*G^\dagger C
 \label{eq:table-native-short}
\end{equation}
is a deterministic function of the historical word Gram, multiplication table
and the fixed word-native compiler.  No independently executed mixed
cross-Hankel panel is required on this branch.
\end{corollary}

\begin{proposition}[Word moments cannot manufacture an external shadow]
\label{thm:external-shadow-no-go}
There are conservative extensions with identical historical word algebra,
multiplication table, complete word-Gram hierarchy and resolved classical
cylinder probabilities but different values of a proposed grading-changing or
typed-incidence cross block whose defining shadow is not represented in the
historical algebra.  Therefore no function of the old word moments alone can
recover such an external shadow.
\end{proposition}

\begin{proof}
Tensor a faithful historical realization with an unread two-level spectator on
which every old operation acts trivially.  One extension contains no coherent
cross-support link; another admits a future-visible spectator link and its
typed compression.  All old word products and moments agree, whereas the
proposed external cross block is zero in the first extension and nonzero in the
second.
\end{proof}


\subsubsection{Large word multiplicities need not be physical internal factors}
\label{app:historical-source-selection}

Let $\Omega=\{(i,j):1\le i,j\le4,\ i\ne j\}$ and
$\Hh=\C^\Omega$, with the simultaneous relabeling action
$u_g|i,j\rangle=|g(i),g(j)\rangle$ of $S_4$.  Adjoin all diagonal point
projections and the two fibre averages
\[
 (Tf)(i,j)=\frac13\sum_{k\ne i}f(i,k),\qquad
 (Hf)(i,j)=\frac13\sum_{k\ne j}f(k,j).
\]
Their support graph is connected, so the generated system algebra is
$M_{12}(\C)$.

\begin{proposition}[Invariant word-space inflation without an internal $M_3$]
\label{prop:word-space-inflation}
For the preceding benchmark,
\[
 \A=M_{12}(\C),\qquad
 \A^{S_4}\cong\C\oplus M_2(\C)\oplus\C\oplus\C.
\]
In the irreducible order
$(\mathbf1,\operatorname{sgn},W,W\otimes\operatorname{sgn},V_2)$, the
multiplicities of the physical $S_4$ action on $\Hh$ are
$(1,0,2,1,1)$, whereas the multiplicities of the conjugation action on
$L^2(M_{12})$ are $(7,5,19,17,12)$.  Thus the latter commutant contains an
$M_7(\C)$ block although the word-induced invariant system algebra has no
$M_3(\C)$ corner.
\end{proposition}

\begin{proof}
The permutation character on ordered unequal pairs is
$(12,2,0,0,0)$ on the conjugacy classes
$1,(12),(12)(34),(123),(1234)$.  Character inner products with the five
irreducibles of $S_4$ give $(1,0,2,1,1)$ and hence the displayed physical
commutant.  Conjugation on $M_{12}\cong\Hh\otimes\overline\Hh$ has character
$|\chi_\Hh|^2=(144,4,0,0,0)$, whose character inner products are
$(7,5,19,17,12)$.  Proposition~\ref{thm:word-module-recognition} retains only
left multiplication by $\A^{S_4}$ as word-induced invariant system actions.
\end{proof}


\subsection{Typed compilation from historical tables}

\begin{corollary}[Historical typed-compilation criterion]
\label{cor:historical-typed-compilation}
Let the relevant historical word hierarchy be flat.  Suppose the grading,
type, weak-leg and incidence source/target projections used to form the
Store--Clifford and matter-incidence shadows are represented elements, or
finite functional-calculus elements, of that same historical $*$-algebra, and
that every router or finite symmetry average entering the compression is
word native in the sense of Definition~\ref{def:word-native-shadow}.  Then the
complete mixed ancestry Gram of those typed shadows, its Schur ancestry
residual, and the alternating determinant shadow are deterministic functions
of the existing Grand-Tensor word Gram and multiplication table.  No
independent mixed cross-Hankel source row is required for these word-native
quantities.  The same table computes the six naturality defects when their
amplitudes and tensor-leg actions are historical; for a mixed branch,
Proposition~\ref{prop:naturality-compilation} instead gives the contextual
Choi-functional criterion.

Conversely, if a required typing projector or router is external to the
historical algebra, Proposition~\ref{thm:external-shadow-no-go} prevents its
historical occurrence from being inferred from the old word moments alone.
\end{corollary}

\begin{proof}
The typed compressions are word-native linear shadows, so
Proposition~\ref{thm:inserted-word-moments} and
Corollary~\ref{cor:table-native-ancestry} reconstruct all ordinary--shadow and
shadow--shadow blocks.  Complete alternation on the already reconstructed typed
Hom space is fixed linear postprocessing, so it introduces no additional
physical source datum.  The converse is exactly
Proposition~\ref{thm:external-shadow-no-go}.
\end{proof}


\section{Tetrahedral representations and internal-algebra reconstruction}
\label{app:internal-landing}

The representation decomposition, general router criterion, reciprocal
return test, and scalar-separation argument support the finite branch
classification in Section~\ref{sec:internal-algebra}.


\subsection{The oriented tetrahedral carrier}
\label{app:S4}

This appendix records the representation-theoretic decomposition used in
Proposition~\ref{thm:geometry-short} for the tetrahedral relational carrier
supplied by \cite{PelissierSpacetime2026}.  The decomposition itself is used
here only as finite representation-theoretic input and is independent of the
probabilistic origin of the source Gram.

Let $X=\{1,2,3,4\}$ and let $\Phi=\{(i,j):i\ne j\}$ be the twelve oriented
edges.  Root reversal is
\[
\iota(i,j)=(j,i).
\]
Let $P_\C=\C^\Phi$.  Its even and odd reversal subspaces are
\[
P^\pm=\{f:f(j,i)=\pm f(i,j)\},
\qquad
\dim P^\pm=6.
\]

\subsubsection{The even sector}

The even sector is the permutation representation on the six unordered edges
of a tetrahedron.  Let $\chi_+$ be its character.  For the five conjugacy
classes of $S_4$ represented by
\[
1,\quad (12),\quad (12)(34),\quad (123),\quad (1234),
\]
the number of fixed unordered edges is
\[
6,\quad 2,\quad 2,\quad 0,\quad 0.
\]
Comparing with the $S_4$ character table gives
\begin{equation}
P^+
\cong
\mathbf1\oplus W\oplus V_2,
\label{eq:appendix-even}
\end{equation}
where $W$ is the standard three-dimensional representation and $V_2$ is the
two-dimensional irreducible representation.

\subsubsection{The odd sector}

The odd sector is the oriented edge space
$\Lambda^2\C^4$ after identifying
$e_i\wedge e_j$ with the signed edge $i\to j$.  Write
\[
\C^4=\C u_0\oplus W,
\qquad
u_0=\frac12(1,1,1,1).
\]
Then
\[
\Lambda^2\C^4
=
u_0\wedge W\oplus\Lambda^2W.
\]
The first summand is isomorphic to $W$.  Since $\dim W=3$,
\[
\Lambda^2W\cong W\otimes{\rm sgn}.
\]
Therefore
\begin{equation}
P^-
\cong
W\oplus(W\otimes{\rm sgn}).
\label{eq:appendix-odd}
\end{equation}
The boundary map sends the first summand isomorphically onto the centred
endpoint space and annihilates the second.  Hence the two odd triplets are
canonically distinguished as endpoint and cycle sectors.

Combining \eqref{eq:appendix-even} and \eqref{eq:appendix-odd}, removing the
even root mean, and adjoining the scalar type contrast gives
\[
\C_{\rm type}^{+}
\oplus W_+\oplus(V_2)_+
\oplus W_-\oplus(W\otimes{\rm sgn})_-,
\]
which is \eqref{eq:residual-S4}.


\subsection{General routers and minimal support residuals}

\begin{proposition}[Connected-router full-matrix theorem]
\label{thm:connected-router}
Let
\[
\mathcal D_n=\Span\{p_1,\ldots,p_n\}\cong\C^n
\]
be the diagonal algebra of mutually orthogonal rank-one projections on
$\C^n$, and let $u\in M_n(\C)$.  Form the undirected support graph with an
edge $i-j$ whenever $p_iup_j\ne0$ or $p_jup_i\ne0$.  Then
\begin{equation}
C^*(\mathcal D_n,u)=M_n(\C)
\quad\Longleftrightarrow\quad
\text{the support graph of $u$ is connected}.
\label{eq:router-full-matrix}
\end{equation}
If the graph is disconnected, every connected component is a reducing
subspace and the generated algebra is a proper block-diagonal subalgebra.
\end{proposition}

\begin{proof}
For every support edge, $p_iup_j$ is a nonzero scalar multiple of the matrix
unit $E_{ij}$; adjoints give $E_{ji}$.  Products along a graph path generate
all matrix units between vertices in the same connected component.  Thus a
connected graph generates all of $M_n(\C)$.  Conversely the sum of the
$p_i$ over any connected component is a nontrivial central reducing
projection when the graph is disconnected.
\end{proof}

\begin{corollary}[One-router source minimality]
\label{cor:one-router-minimality}
Once a coherent $n$-dimensional multiplicity carrier and its diagonal minimal
projections have occurred, one cyclic nonnormal router is sufficient to
generate $M_n(\C)$.  Without any off-diagonal word the generated algebra
remains commutative.  Thus one router is source-minimal at fixed multiplicity
carrier.
\end{corollary}

\begin{proposition}[Minimal bridge and weak-router residuals]
\label{prop:bridge-router-residuals}
On $T\oplus H_{\rm priv}$ define
\[
\delta_{\rm col}(u)
:=
\norm{p_Hup_T}_{\HS}^2+\norm{p_Tup_H}_{\HS}^2.
\]
On a weak multiplicity plane define
\[
\delta_{\rm wk}(p,q):=\norm{[p,q]}_{\HS}^2.
\]
Then $\delta_{\rm col}=0$ exactly on the reducing colour branch
$M_2\oplus\C$, while $\delta_{\rm col}>0$ gives $M_3$.  For rank-one
projections $p,q$ on $\C^2$, one has
\begin{equation}
\delta_{\rm wk}(p,q)
=2c^2(1-c^2),
\qquad c=\abs{\langle t,h\rangle},
\label{eq:weak-router-residual}
\end{equation}
so $\delta_{\rm wk}>0$ exactly when the projections are distinct and
nonorthogonal, in which case they generate $M_2$.
\end{proposition}

\begin{proof}
The colour statement is Proposition~\ref{thm:colour-bridge}.  For the weak
statement, choose a basis with $p=|e_1\rangle\langle e_1|$ and
$h=ce_1+\sqrt{1-c^2}e_2$ after absorbing a phase.  Direct multiplication gives
\eqref{eq:weak-router-residual}.  Its strict positivity is equivalent to
$0<c<1$, which is the Burnside condition of Proposition~\ref{thm:colour-bridge}.
\end{proof}


\subsection{Subdirect matrix algebras and scalar separation}

\begin{proposition}[Finite subdirect matrix-algebra classification]
\label{thm:subdirect-classification}
Let $\mathcal B\subseteq\bigoplus_{i=1}^rM_{n_i}(\C)$ be a unital
$C^*$-subalgebra whose coordinate projection onto every displayed factor is
surjective.  Then the coordinates partition into classes of equal matrix size,
and $\mathcal B$ contains one independent full matrix algebra per class,
embedded diagonally across that class up to inner automorphisms.  In
particular, distinct matrix sizes cannot be diagonally identified.
\end{proposition}

\begin{proof}
Write the intrinsic decomposition
$\mathcal B\cong\bigoplus_\beta M_{m_\beta}(\C)$.  A surjective
$*$-homomorphism to a simple factor $M_{n_i}(\C)$ annihilates every intrinsic
summand except one: the images of the intrinsic central units are orthogonal
central projections summing to the identity, so exactly one is nonzero.  On
that summand the map is a nonzero homomorphism between simple full matrix
algebras and hence an isomorphism, forcing $m_\beta=n_i$.  Coordinates seeing
the same intrinsic summand are related by automorphisms of a full complex
matrix algebra, hence by inner automorphisms.  Different intrinsic summands
are independent.
\end{proof}

Applying this proposition to matrix sizes $(3,2,1,1)$ leaves only the two
scalar factors available for diagonal identification.  Their independent
central supports, or the positive character-separation residual, exclude it.


\subsection{Reciprocal returns through a represented mediator}
\label{app:reciprocal-return}

The bridge need not be a primitive operation when the historical route table
contains reciprocal access through one common mediator.  Let $G$ be the
external finite symmetry, let $V$ be an irreducible unitary $G$-module of
dimension $d$, and let the mediator carrier be
$\mathcal E=V\otimes\mathcal N$ with $n=\dim\mathcal N$.  Suppose actual entry
and reverse-exit operations have the covariant form
\[
 A_g=(v_g\otimes I)A:T\to\mathcal E,
 \qquad
 B_h^*=B^*(v_h^*\otimes I):\mathcal E\to H_{\rm priv},
\]
and that the actual mediator dynamics contains
$D=I_V\otimes h$ with $h=h^*$.  Put
\[
 r_{g,k}=B^*(v_g\otimes h^k)A,
 \qquad 0\le k<n,
\]
\[
 \rho_A=\Tr_V(AA^*),\qquad
 \rho_B=\Tr_V(BB^*),\qquad
 W_A=\sum_{k=0}^{n-1}h^k\rho_Ah^k.
\]

\begin{proposition}[Finite reciprocal-return criterion]
\label{thm:reciprocal-return}
For the declared route bank,
\begin{equation}
 \mathfrak m_{\rm ret}
 :=\frac1{|G|}\sum_{g\in G}\sum_{k=0}^{n-1}
       \norm{r_{g,k}}_{\HS}^2
 =\frac1d\Tr(\rho_BW_A)\ge0.
 \label{eq:reciprocal-return-mass}
\end{equation}
The following are equivalent:
\begin{enumerate}[label=\textup{(\roman*)},leftmargin=2.8em]
\item an admitted word in the route bank has a nonzero
      $H_{\rm priv}\leftarrow T$ corner;
\item some $r_{g,k}$ is nonzero;
\item $\mathfrak m_{\rm ret}>0$;
\item the $h$-cyclic multiplicity spaces generated by $\Ran\rho_A$ and
      $\Ran\rho_B$ are not orthogonal.
\end{enumerate}
If these conditions fail, an explicit reducing projection separates $T$ from
$H_{\rm priv}$ for the entire declared bank, so arbitrary repeated excursions
cannot create a colour bridge within that bank.  If they hold, a witness has
at most $n+1$ route letters.
\end{proposition}

\begin{proof}
Schur averaging on the irreducible $V$ factor gives
\[
 \frac1{|G|}\sum_g
 (v_g\otimes I)AA^*(v_g^*\otimes I)
 =\frac{I_V}{d}\otimes\rho_A.
\]
Taking the Hilbert--Schmidt norm of $B^*(v_g\otimes h^k)A$ and summing yields
\eqref{eq:reciprocal-return-mass}.  Since it is a finite sum of squared norms,
positivity is equivalent to a nonzero returned word.  The range of $W_A$ is
$\Span\{h^k\Ran\rho_A:0\le k<n\}$ by Cayley--Hamilton.  Thus
$\Tr(\rho_BW_A)=0$ exactly when the two cyclic multiplicity spaces are
orthogonal.  In that case the projection onto
$T\oplus(V\otimes\Ran W_A)$ commutes with the route generators and their
adjoints while annihilating $H_{\rm priv}$, excluding all longer returns.  The
finite horizon $k<n$ gives the stated word-length bound.
\end{proof}

\begin{corollary}[Port-relative reachability forces colour]
\label{cor:reciprocal-return-force}
If $W_A\succeq\kappa I_{\mathcal N}$ with $\kappa>0$ and $B\ne0$, then
\begin{equation}
 \mathfrak m_{\rm ret}\ge\frac{\kappa}{d}\norm B_{\HS}^2,
 \qquad
 \max_{g,k<n}\norm{r_{g,k}}_{\HS}^2
 \ge\frac{\kappa}{dn}\norm B_{\HS}^2.
 \label{eq:reciprocal-return-margin}
\end{equation}
Hence the colour bridge follows from the upstream route structure and
Proposition~\ref{thm:colour-bridge}.  In particular, for a multiplicity-one
mediator, two nonzero reciprocal ports force a two-leg colour return and no
independent direct bridge is required.
\end{corollary}

For the mediator notation of Proposition~\ref{thm:reciprocal-return}, set
\[
 \mathcal N_A=\Span\{h^k\Ran\rho_A:0\le k<n\},\qquad
 \mathcal N_B=\Span\{h^k\Ran\rho_B:0\le k<n\}.
\]
Cayley--Hamilton makes both spaces $h$-invariant, and self-adjointness of $h$
makes them reducing.

\begin{proposition}[All-excursion return alternative]
\label{prop:all-excursion-return}
If $\mathfrak m_{\rm ret}=0$, then
$\mathcal N_A\perp\mathcal N_B$ and
\[
 Q=I_T\oplus0_{H_{\rm priv}}\oplus(I_V\otimes P_{\mathcal N_A})
\]
commutes with every generator of the declared reciprocal route bank and its
adjoint.  Consequently every word $w$ in that bank has
$p_Hwp_T=0$.  Thus failure of the finite return panel is an exact obstruction
to arbitrarily long repeated excursions within the same bank.
\end{proposition}

\begin{proof}
The equality $\mathfrak m_{\rm ret}=d^{-1}\Tr(\rho_BW_A)=0$ and positivity
imply $\Ran\rho_B\perp\mathcal N_A$, hence
$\mathcal N_B\perp\mathcal N_A$.  Partial-trace positivity gives
$\Ran A\subseteq V\otimes\Ran\rho_A$ and
$\Ran B\subseteq V\otimes\Ran\rho_B$.  The stated projection therefore
reduces the entry, exit, mediator and adjoint generators, contains the type
line and annihilates the private plane.  Every generated word has zero
$H_{\rm priv}\leftarrow T$ corner.
\end{proof}


\subsection{Router failure loci and the finite alternatives}

For the diagonal algebra $\mathcal D_n$ of
Proposition~\ref{thm:connected-router}, define the router failure set
\[
\mathfrak F_n
=
\{u\in M_n(\C):C^*(\mathcal D_n,u)\ne M_n(\C)\}.
\]
For every nonempty proper subset $S\subset\{1,\ldots,n\}$ let
\[
L_S
=
\{u:p_iup_j=p_jup_i=0\ \text{for all }i\in S,
\ j\notin S\}.
\]

\begin{proposition}[Algebraic structure of the router failure set]
\label{prop:router-failure-locus}
One has
\begin{equation}
\mathfrak F_n
=
\bigcup_{\varnothing\ne S\subsetneq\{1,\ldots,n\}}L_S.
\label{eq:router-failure-locus}
\end{equation}
Each $L_S$ is a proper complex linear subspace of $M_n(\C)$.  Hence
$\mathfrak F_n$ is a proper Zariski-closed set and its complement is Zariski
open, Euclidean open, and dense.
\end{proposition}

\begin{proof}
The support graph is disconnected exactly when some nontrivial union $S$ of
connected components has no support edge to its complement, which is precisely
membership in $L_S$.  Each $L_S$ is defined by homogeneous linear equations
on matrix entries and is proper because, for example, a cyclic shift matrix
has connected support.  A finite union of closed algebraic subsets is closed,
and a finite union of proper linear subspaces cannot fill $M_n(\C)$.
\end{proof}

In particular, fixing diagonal minimal projections on the admitted
three- and two-dimensional multiplicity carriers gives a Zariski-open dense
set of routers generating $M_3(\C)\oplus M_2(\C)$.  This statement is
conditional on those carriers having occurred.

\begin{proposition}[Finite internal landing alternative]
\label{thm:finite-internal-landing-alternative}
The following finite implications apply to the colour/weak landing problem.
Failure of a finite-depth test without flatness is not a terminal obstruction;
the implications are not asserted to be mutually exclusive.
\begin{enumerate}[label=\textup{(A\arabic*)},leftmargin=3em]
\item If the word hierarchy is flat and $m_{0,r}<3$, an external-trivial
      internal colour $M_3(\C)$ is impossible at every later depth.
\item If $m_{0,r}\ge3$ but every represented type--private bridge has
      $\delta_{\rm col}=0$, the colour block is reducible and remains
      $M_2\oplus\C$ on the displayed three-space.
\item If the weak multiplicity plane is present but every available pair has
      $\delta_{\rm wk}=0$, the represented weak multiplicity algebra is
      commutative.
\item If the complete census is $\bm m_r=(3,0,2,1,1)$, both a positive colour
      residual and a positive weak-router residual occur, and the two scalar
      characters are centrally separated, the represented
      internal word algebra equals the full relative commutant
      $M_3\oplus M_2\oplus\C\oplus\C$.
\end{enumerate}
On the fixed complete census $\bm m_r=(3,0,2,1,1)$ with independent displayed central supports,
the intermediate algebra dimensions are exactly those of
Proposition~\ref{thm:internal-deficit-alternatives}.  Without scalar central
separation the blockwise-surjective terminal branch may instead have
dimension fourteen.  Therefore
every failure returns a finite algebraic witness and no failed condition is
silently completed to the terminal branch.
\end{proposition}

\begin{proof}
Item (A1) is Proposition~\ref{thm:active-isotypic}.  Item (A2) is
Proposition~\ref{thm:colour-bridge}; item (A3) is
Proposition~\ref{thm:colour-bridge}; and item (A4) is
Theorem~\ref{thm:internal-seed-saturation}.  The dimension statement is
Proposition~\ref{thm:internal-deficit-alternatives}.
\end{proof}


\section{Determinant naturality, incidence normalization, and anomalies}
\label{app:naturality}

The principal determinant and gauge-group proofs are in
Section~\ref{sec:gauge-matter}.  Here we retain the mixed-occurrence alignment
criterion, the unrestricted incidence envelopes, the exact naturality
spectrum, and the source and anomaly bookkeeping.


\subsection{Tensor alignment from mixed source occurrences}
\label{app:tensor-alignment}

\begin{proposition}[Fusion-native sufficient criterion for same-carrier alignment]
\label{thm:fusion-native-alignment}
Let $H_C,H_1,H_2$ be finite source-minimal Hilbert spaces with source
syntheses $S_A:E_A\to H_A$ and Grams $G_A=S_A^*S_A$.  Let
$H_{C1},H_{12},H_{C12}$ be the source-minimal carriers of the represented
pair and triple occurrences, with syntheses $S_{C1},S_{12},S_{C12}$ and
coefficient occurrence maps
\[
 m_{C,1}:E_C\otimes E_1\to E_{C1},\qquad
 m_{1,2}:E_1\otimes E_2\to E_{12},
\]
\[
 m_{C1,2}:E_{C1}\otimes E_2\to E_{C12},\qquad
 m_{C,12}:E_C\otimes E_{12}\to E_{C12}.
\]
Assume the coefficient occurrence law is associative,
\begin{equation}
 m_{C1,2}(m_{C,1}\otimes I)=m_{C,12}(I\otimes m_{1,2}).
 \label{eq:fusion-coefficient-associativity}
\end{equation}
For each composable pair $(A,B)$ suppose
\begin{equation}
 m_{A,B}^*G_{AB}m_{A,B}=G_A\otimes G_B.
 \label{eq:fusion-metric-identity}
\end{equation}
Then the rule on represented source vectors
\begin{equation}
 \Gamma_{A,B}(S_Ax\otimes S_By)
 :=S_{AB}m_{A,B}(x\otimes y)
 \label{eq:fusion-isometry-definition}
\end{equation}
extends uniquely to an isometry $\Gamma_{A,B}:H_A\otimes H_B\to H_{AB}$.
Assume in addition that the four fusion innovations vanish,
\begin{equation}
 S_{AB}^*(I-\Gamma_{A,B}\Gamma_{A,B}^*)S_{AB}=0
 \quad\text{for }(A,B)=(C,1),(1,2),(C1,2),(C,12),
 \label{eq:fusion-zero-innovation}
\end{equation}
and that the associator defect vanishes,
\begin{equation}
 \left\|
 \Gamma_{C1,2}(\Gamma_{C,1}\otimes I)
 -\Gamma_{C,12}(I\otimes\Gamma_{1,2})
 \right\|_{\HS}^2=0.
 \label{eq:fusion-associator-zero}
\end{equation}
Then the two parenthesizations define the same unitary
\begin{equation}
 \Gamma:H_C\otimes H_1\otimes H_2\xrightarrow{\sim}H_{C12}.
 \label{eq:fusion-triple-unitary}
\end{equation}
Consequently the three transported factor actions
\begin{align}
 \widehat\pi_C(a)&=\Gamma(a\otimes I\otimes I)\Gamma^*,\nonumber\\
 \widehat\pi_1(b)&=\Gamma(I\otimes b\otimes I)\Gamma^*,\nonumber\\
 \widehat\pi_2(c)&=\Gamma(I\otimes I\otimes c)\Gamma^*
 \label{eq:fusion-native-actions}
\end{align}
commute pairwise.  If $\dim H_C=3$ and
$\dim H_1=\dim H_2=2$, they generate $\mathcal B(H_{C12})\cong M_{12}(\C)$,
and the same-carrier defect $\Delta_{\otimes}$ of
Proposition~\ref{thm:tensor-factor-alignment} vanishes for these factor actions.
With an additional reducing multiplicity space $M$, the generated algebra is
$M_{12}(\C)\otimes I_M$.
\end{proposition}

\begin{proof}
Equation~\eqref{eq:fusion-metric-identity} says exactly that
\eqref{eq:fusion-isometry-definition} preserves inner products on the source
span, hence defines an isometry.  Source minimality gives
$\Ran S_{AB}=H_{AB}$; therefore
\eqref{eq:fusion-zero-innovation} implies that every $\Gamma_{A,B}$ is onto
and hence unitary.  The associator defect identifies the two triple unitaries.
The three tensor-leg matrix algebras commute before conjugation by $\Gamma$ and
therefore after conjugation.  Their tensor product is
$\mathcal B(H_C\otimes H_1\otimes H_2)$, so in dimensions $(3,2,2)$ the
transported factors generate $M_{12}(\C)$.  The commutator defect is a sum of
squared norms and consequently vanishes.  Tensoring all factors with $I_M$
gives the final statement.
\end{proof}

The direct commutation residual applies when factor actions are supplied
independently.  The preceding sufficient criterion instead reconstructs
commuting actions from the actual mixed occurrence packet.  Neither follows
from separate factor marginals, as the next witness shows.

\begin{proposition}[Factor marginals do not fix their relative tensor position]
\label{cth:factor-marginals-no-alignment}
There are two pairs of unital factor embeddings on the same six-dimensional
Hilbert space such that each colour marginal is $M_3(\C)$, each weak marginal
is $M_2(\C)$, and all intrinsic single-factor traces and Hilbert--Schmidt Grams
agree, while one pair commutes and the other does not.
\end{proposition}

\begin{proof}
On $\C^3\otimes\C^2$ put
\[
\mathcal A=M_3(\C)\otimes I_2,
\qquad
\mathcal B=I_3\otimes M_2(\C).
\]
Choose a non-scalar Hermitian $D\in M_3(\C)$, a Pauli matrix
$X\in M_2(\C)$, and for generic nonzero $\theta$ define
\[
U=\exp(i\theta D\otimes X),
\qquad
\mathcal B'=U\mathcal B U^*.
\]
The algebras $\mathcal B$ and $\mathcal B'$ are unitarily conjugate, so every
intrinsic algebraic invariant, trace, and Hilbert--Schmidt Gram computed within
the weak marginal agrees.  The colour marginal is unchanged.  The pair
$(\mathcal A,\mathcal B)$ commutes.  For generic $\theta$, however, choose a
colour matrix unit not commuting with $D$ and a weak generator not commuting
with $X$; conjugation by $U$ produces an element of $\mathcal B'$ with a
nonzero commutator with that colour element.  Hence the second pair is not a
commuting tensor-factor realization despite identical separate marginals.
\end{proof}


\subsection{Rectangular and semi-intertwiner envelopes}

Let an input carrier split as $A=\bigoplus_{s\in\mathcal S}A_s$,
with $p_s=\dim A_s$.  A recorded branch $a$ has source $A_{s(a)}$
and target $T_a$ of dimension $q_a$.  Suppose its only amplitude restriction
is the full rectangular envelope $\Hom(A_{s(a)},T_a)$.
Use the coefficient convention in which the normalization synthesis of
matrix-unit amplitudes $B_{a,i\alpha}=|f_i\rangle\langle e_\alpha|$
sends a coefficient matrix unit to $B_{a,i\alpha}^*B_{a,j\beta}$.
This is the usual Choi normalization map after the fixed input-transpose
identification; its rank and kernel do not depend on that convention.

\begin{proposition}[Exact rectangular normalization kernel]
\label{thm:rectangular-normalization-kernel}
If every source sector occurs in at least one branch, the normalization map is
\begin{equation}
 \Gamma_{\rm rect}((X_a)_a)
 =\bigoplus_{s\in\mathcal S}
       \sum_{a:s(a)=s}\Tr_{T_a}X_a,
 \qquad X_a\in M_{q_ap_s}(\C).
 \label{eq:rectangular-normalization-map}
\end{equation}
It is surjective onto $\bigoplus_s M_{p_s}(\C)$ and has complex nullity
\begin{equation}
 \dim_\C\Ker\Gamma_{\rm rect}
 =\sum_s p_s^2\left(\sum_{a:s(a)=s}q_a^2-1\right).
 \label{eq:rectangular-normalization-nullity}
\end{equation}
It is injective exactly when every source sector has one outgoing recorded
branch and that branch has one-dimensional target.
\end{proposition}

\begin{proof}
The products of amplitude matrix units are
$B_{a,i\alpha}^*B_{a,j\beta}=\delta_{ij}|e_\alpha\rangle\langle e_\beta|$.
This proves \eqref{eq:rectangular-normalization-map}.  Any one branch on a
source sector makes its partial trace surjective onto $M_{p_s}(\C)$.
The domain dimension on that sector is $p_s^2\sum_{a:s(a)=s}q_a^2$ and
the image dimension is $p_s^2$, giving the nullity.  All $q_a$ are positive
integers, so vanishing of each summand is equivalent to the asserted
one-branch, one-dimensional-target condition.
\end{proof}

For the incidence port of \eqref{eq:naturality-port}, $p=12$, $q=3$;
therefore the normalization kernel has dimension $1152$.  This is a
statement about the unrestricted envelope, not the nullity at an already
rank-one-supported Choi packet.  Finite naturality removes the unwanted
amplitude directions before normalization is used.


\subsubsection{Schur compression before normalization}

\begin{proposition}[Typed semi-intertwiner envelope]
\label{prop:schur-incidence-envelope}
Let a compact group act on a source $H_s$ and target $H_t$, and let $\chi$
be a unitary character.  Decompose
\[
 H_s=\bigoplus_\lambda V_\lambda\otimes\C^{m_\lambda},\qquad
 \chi^{-1}H_t=\bigoplus_\lambda V_\lambda\otimes\C^{n_\lambda}.
\]
Then the amplitude envelope
\[
 \mathsf S_\chi=\{K:R_t(g)KR_s(g)^*=\chi(g)K\}
\]
is isomorphic to
$\bigoplus_\lambda I_{V_\lambda}\otimes
M_{n_\lambda\times m_\lambda}(\C)$ and has dimension
$\sum_\lambda m_\lambda n_\lambda$.
Every $K^*L$ with $K,L\in\mathsf S_\chi$ lies in $R_s(G)'$.
For several recorded semi-intertwiner envelopes on this source, injectivity
of their joint normalization map therefore requires
\begin{equation}
 \sum_a\left(\sum_\lambda m_\lambda n_{a,\lambda}\right)^2
 \le\sum_\lambda m_\lambda^2.
 \label{eq:schur-normalization-screen}
\end{equation}
\end{proposition}

\begin{proof}
Twist the target by $\chi^{-1}$ and apply Schur's lemma.  The character
cancels in $K^*L$, so the normalization image lies in the source commutant,
of the dimension on the right of \eqref{eq:schur-normalization-screen}.
Injectivity requires the coefficient-domain dimension not to exceed it.
\end{proof}


\subsection{Exact naturality spectrum}
\label{app:naturality-spectrum}

Fix an oriented orthonormal basis $(c_1,c_2,c_3)$ of $C$ and its dual
$(c^1,c^2,c^3)$.  The unitary volume identification
$J_C:C\to\Lambda^2C^*$ is
\[
 J_Cc_1=c^2\wedge c^3,\qquad
 J_Cc_2=-c^1\wedge c^3,\qquad
 J_Cc_3=c^1\wedge c^2.
\]
It obeys
\begin{equation}
 J_C^{-1}d\rho_C^{\rm out}(X)J_C=X-(\Tr X)I_C.
 \label{eq:volume-infinitesimal-action}
\end{equation}
This identity follows by evaluating on the three exterior basis vectors,
or by differentiating the natural action on
$C\otimes(\det C)^*\cong\Lambda^2C^*$.
For every off-diagonal matrix unit the trace term vanishes.
Thus on
\[
 \mathscr U\cong M_3(\C)\otimes(W_2\otimes W_2)^*
\]
the positive defect splits into commuting tensor summands
\begin{equation}
 H_{\rm nat}\cong H_C\otimes I_4+I_9\otimes H_W,
 \qquad H_C=\sum_{X\in\mathcal G_C}\operatorname{ad}_X^*
                                       \operatorname{ad}_X.
 \label{eq:naturality-tensor-sum}
\end{equation}
Here $H_W$ is the dual-space version of $L_+L_-+L_-L_+$, where
$L_\pm=F_{12/21}\otimes I+I\otimes F_{12/21}$.  Transposition onto
the dual space preserves the following spectral calculation.

\begin{lemma}[Colour defect spectrum]
\label{lem:colour-naturality-spectrum}
The spectrum of $H_C$ is
\begin{equation}
 \spec H_C=\{0^{(1)},2^{(3)},3^{(4)},6^{(1)}\}.
 \label{eq:colour-naturality-spectrum}
\end{equation}
\end{lemma}

\begin{proof}
For the undirected edges $1-2$ and $2-3$, the two opposite matrix-unit
terms contribute
\[
 (E_{ii}+E_{jj})A+A(E_{ii}+E_{jj})
       -2(E_{ij}AE_{ji}+E_{ji}AE_{ij}).
\]
On an off-diagonal matrix unit $E_{ab}$, this is diagonal with eigenvalue
$d_a+d_b$, where $(d_1,d_2,d_3)=(1,2,1)$.  Thus $E_{13},E_{31}$ have
eigenvalue $2$ and $E_{12},E_{21},E_{23},E_{32}$ have eigenvalue $3$.
On diagonal matrices the action is twice the three-vertex path Laplacian:
\[
 2\begin{pmatrix}1&-1&0\\-1&2&-1\\0&-1&1\end{pmatrix}.
\]
Its eigenvectors $(1,1,1)$, $(1,0,-1)$, and $(1,-2,1)$ have eigenvalues
$0$, $2$, and $6$, respectively.  These nine eigenvectors give the spectrum.
\end{proof}

\begin{lemma}[Weak defect spectrum]
\label{lem:weak-naturality-spectrum}
The spectrum of $H_W$ is
\begin{equation}
 \spec H_W=\{0^{(1)},2^{(2)},4^{(1)}\}.
 \label{eq:weak-naturality-spectrum}
\end{equation}
Its zero eigenspace is the alternating weak covector line.
\end{lemma}

\begin{proof}
With an orthonormal basis $(w_1,w_2)$, both $L_+$ and $L_-$ annihilate
$(w_1\otimes w_2-w_2\otimes w_1)/\sqrt2$.
On the symmetric orthonormal vectors
\[
 w_1\otimes w_1,\quad
 (w_1\otimes w_2+w_2\otimes w_1)/\sqrt2,\quad
 w_2\otimes w_2,
\]
direct multiplication of $L_+L_-+L_-L_+$ gives eigenvalues $2$, $4$, and
$2$.  Passing to the dual representation gives the same spectrum.
\end{proof}

\begin{proposition}[Exact naturality spectrum]
\label{prop:exact-naturality-spectrum}
The complete spectrum is
\begin{equation}
 \spec H_{\rm nat}
 =\{0^{(1)},2^{(5)},3^{(4)},4^{(7)},5^{(8)},
       6^{(4)},7^{(4)},8^{(2)},10^{(1)}\}.
 \label{eq:exact-naturality-spectrum}
\end{equation}
In particular, the one-dimensional kernel is the tensor product of the
scalar colour matrix and alternating weak covector, and the first positive
eigenvalue is exactly two.
\end{proposition}

\begin{proof}
Add the eigenvalues of \eqref{eq:colour-naturality-spectrum} and
\eqref{eq:weak-naturality-spectrum}, multiplying their multiplicities, in
the tensor sum \eqref{eq:naturality-tensor-sum}.  The resulting multiplicities
sum to $36$.  The smallest positive sum is two, which proves
\eqref{eq:naturality-gap} without a numerical spectral assumption.
\end{proof}


\subsection{Source-independent selection and positive branch refinements}

\begin{proposition}[Determinant selector, population, and branchwise purity]
\label{thm:determinant-skeleton-population}
On the aligned tensor-factor branch, $(H_{\rm nat},P_{\det})$ is defined
entirely by the factor actions and typed slots, before any source population.
Let $J=\sum_bJ_b$, $J_b\succeq0$, and put
\begin{equation}
 m_b=\Tr J_b,\qquad d_b=\Tr(H_{\rm nat}J_b),\qquad
 m_{D,b}=\Tr(P_{\det}J_b).
 \label{eq:determinant-branch-masses}
\end{equation}
For the corresponding totals, $d=\sum_bd_b$, $m_D=\sum_bm_{D,b}$, and
$d_b\ge0$.  The exact spectrum gives
\begin{equation}
 2(I-P_{\det})\preceq H_{\rm nat}\preceq10(I-P_{\det}),
 \label{eq:naturality-two-sided-gap}
\end{equation}
so, branchwise and in total,
\begin{equation}
 \Tr[(I-P_{\det})J_b]\le d_b/2,\qquad
 \max\{0,m_b-d_b/2\}\le m_{D,b}\le m_b-d_b/10.
 \label{eq:determinant-branch-bounds}
\end{equation}
If $d=0$, every positive branch satisfies
\begin{equation}
 J_b=m_bP_{\det}.\label{eq:determinant-branch-purity}
\end{equation}
Hence $m>0$ gives $J=mP_{\det}$ and forgetting positive labels cannot hide
non-determinant contamination.  If merely $d<2m$, then $m_D>0$, certifying
occurrence without exact purity.  Simultaneous unitary changes of factor
coordinates preserve all these scalar statements.
\end{proposition}
\begin{proof}
The defect maps depend on the factor actions, not on $J$ or its Kraus
coordinates.  Linearity and positivity give additivity and $d_b\ge0$.
Proposition~\ref{prop:exact-naturality-spectrum} places the positive spectrum
in $[2,10]$; tracing the resulting operator bounds against each $J_b$
gives the inequalities.  If $d=0$, every $d_b=0$, so positivity forces each
$J_b$ onto the one-dimensional kernel of $H_{\rm nat}$.  Its trace then
gives $J_b=m_bP_{\det}$.  Finally $m_D\ge m-d/2$, and unitary covariance
of the factor actions conjugates the selector and source simultaneously.
\end{proof}


\subsection{Determinant realization, normalization, and tensor arity}
\label{app:anomaly}

For $C\cong\C^3$ and $W_2\cong\C^2$, define
\[
\mathscr D=(\det C\otimes\det W_2)^*,
\qquad
\mathscr U=
\Hom(C\otimes W_2\otimes W_2,\Lambda^2C^*).
\]
For $\tau\in\mathscr D$, define
\[
\bigl(\Theta_\tau(c\otimes w_1\otimes w_2)\bigr)(c_1,c_2)
=
\tau(c\wedge c_1\wedge c_2\otimes w_1\wedge w_2).
\]
Complete alternation in the three colour and two weak arguments gives a map
$\mathfrak A_{\det}:\mathscr U\to\mathscr D$.

\begin{proposition}[Determinant-incidence line]
\label{prop:det-incidence}
The maps $\Theta$ and $\mathfrak A_{\det}$ restrict to inverse isomorphisms
between $\mathscr D$ and the fully alternating incidence line.  In oriented
orthonormal bases,
\[
\mathfrak A_{\det}(\Theta_\tau)=\tau,
\qquad
\norm{\Theta_\tau}_{\HS}^2=6\norm{\tau}^2.
\]
Thus a nonzero fully alternating incidence shadow is equivalent to a nonzero
determinant seed.
\end{proposition}

\begin{proof}
Complete alternation is the orthogonal projection onto the unique
$(\det C\,\det W_2)^{-1}$ symmetry type.  Evaluation on oriented orthonormal
bases gives the factor $3!\,2!/2=6$ in the Hilbert--Schmidt normalization.
\end{proof}


\subsubsection{Partial-isometry normalization and competing branches}

Let $a=(w_1\otimes w_2-w_2\otimes w_1)/\sqrt2$ and define
$T_0(c\otimes\zeta)=J_Cc\,\langle a,\zeta\rangle$.
Then
\begin{equation}
 T_0^*T_0=I_C\otimes|a\rangle\langle a|=P_{\rm alt},
 \qquad T_0T_0^*=I_{\Lambda^2C^*},\qquad\norm{T_0}_{\HS}^2=3.
 \label{eq:normalized-determinant-effect}
\end{equation}
For unit $\tau$ in Proposition~\ref{prop:det-incidence},
$T_0=\Theta_\tau/\sqrt2$, up to a common phase.  This distinguishes
partial-isometry normalization from unit Hilbert--Schmidt normalization.

\begin{proposition}[Alternating-sector isolation fixes the branch weight]
\label{prop:determinant-normalization-isolation}
Let a normalized instrument contain a branch satisfying
\eqref{eq:naturality-ray}, with $T_0$ normalized by
\eqref{eq:normalized-determinant-effect}.  Suppose every Kraus amplitude $K$
of every other recorded branch satisfies $KP_{\rm alt}=0$.  Then $x=1$.
If instead $k$ recorded branches on this sector have the same normalized
one-dimensional envelope $\C T_0$, their effects constrain only
$\sum_{a=1}^k x_a=1$; the linear scalar normalization kernel has dimension
$k-1$.
\end{proposition}

\begin{proof}
Compress the trace-preservation identity on both sides by $P_{\rm alt}$.
The singled-out branch contributes $xP_{\rm alt}$ and the other branches
contribute zero, proving $x=1$.  For $k$ competing ray branches, compression
gives $(\sum_a x_a)P_{\rm alt}=P_{\rm alt}$, a single independent scalar
equation.  Positivity gives $x_a\ge0$ but does not remove the scalar
ambiguity in the interior of that simplex.
\end{proof}


\subsubsection{Sharp arity of the determinant character}

Let a native finite carrier have defining central weights $(1,0)$ on
$C\cong\C^3$, $(0,1)$ on $W_2\cong\C^2$, and $(0,0)$ on any neutral sector.
The centre of $\U(3)\times\U(2)$ acts on a weight $(a,b)$ by
$e^{i(a\alpha+b\beta)}$.

\begin{proposition}[Sharp five-slot determinant-mode floor]
\label{thm:determinant-arity-floor}
The central-weight support of the single-system endomorphism algebra contains
only
\[
 (0,0),\qquad \pm(1,0),\qquad \pm(0,1),\qquad \pm(1,-1),
\]
and therefore does not contain the determinant-incidence weight $(-3,-2)$.
On a diagonal $k$-fold tensor carrier this character is absent for $k<5$.
If the native carrier contains a nonzero neutral line, the bound is attained
at $k=5$.
\end{proposition}

\begin{proof}
Single-system endomorphism weights are target-minus-source differences among
$(1,0),(0,1),(0,0)$, giving the displayed list.  In $k$ tensor slots every
vector weight $(a,b)$ has $a,b\ge0$ and $a+b\le k$, so a target-minus-source
weight with total negative charge $-5$ cannot occur for $k<5$.  At $k=5$ the
alternating line in $C^{\otimes3}\otimes W_2^{\otimes2}$ transforms by
$\det C\,\det W_2$; a rank-one map from that line to a neutral line therefore
has weight $(-3,-2)$.
\end{proof}

This result concerns tensor arity, not temporal word length.  It explains why
longer evolution inside a fixed low-arity memory cannot substitute for the
typed five-leg incidence of Theorem~\ref{thm:six-naturality}.


\subsection{Finite functional and word-moment evaluation}

\begin{proposition}[Finite functional and word-moment realization]
\label{prop:naturality-compilation}
Fix the represented tensor-leg actions.  Suppose an actual contextual test
bank determines numbers $\Tr(F_\ell J_D)$ and its real linear span contains
$H_{\rm nat}$ and $I_{\mathscr U}$.  Then this bank determines
$d_{\rm nat}$ and $\Tr J_D$ without determining the full Choi operator.
More generally it suffices that those two functionals belong to the span
modulo linear functionals constant on the declared feasible Choi class.

Alternatively, if a historical amplitude bank $(T_a)$ and all inserted
matrix-factor actions are represented historical operations, then
\begin{equation}
 d_{\rm nat}
 =\sum_a\left(\sum_{X\in\mathcal G_C}\norm{\mathcal R_X(T_a)}_{\HS}^2
              +\sum_{Y\in\mathcal G_W}\norm{\mathcal S_Y(T_a)}_{\HS}^2\right)
 \label{eq:word-naturality-defect}
\end{equation}
is a finite sum of inserted word moments.  At flatness their multiplication
table computes it exactly.  This second route is conditional on amplitude
access; the contextual route does not require a Kraus refinement.
\end{proposition}

\begin{proof}
Express each required functional as a real linear combination of the
available tester functionals, adding its known constant when necessary.
For the second route, use
$J_D=\sum_a|T_a\rangle\!\rangle\langle\!\langle T_a|$ and
\eqref{eq:naturality-operator}.  Each squared norm expands into traces of
products of the represented amplitudes and inserted actions, to which
Proposition~\ref{thm:inserted-word-moments} applies on the corresponding
represented source algebra.  A known faithful trace convention must be
used consistently; converting between a physical trace and the regular
trace requires their reconstructed block weights.  Finite table reduction
then evaluates the products.  Identification of a scalar functional does
not itself implement that functional as a physical measurement operation.
\end{proof}


\subsection{Chiral conventions and the full gauge-charged audit}

With $(a,b)=(-2,3)$, the physical right-handed convention gives
\[
Q:(\mathbf3,\mathbf2)_1,\qquad
u:(\mathbf3,\mathbf1)_4,\qquad
d:(\mathbf3,\mathbf1)_{-2},
\]
\[
L:(\mathbf1,\mathbf2)_{-3},\qquad
e:(\mathbf1,\mathbf1)_{-6},\qquad
H:(\mathbf1,\mathbf2)_3.
\]
Writing all fermions as left-handed Weyl fields replaces
\[
u\mapsto u^c:(\overline{\mathbf3},\mathbf1)_{-4},
\qquad
d\mapsto d^c:(\overline{\mathbf3},\mathbf1)_2,
\qquad
e\mapsto e^c:(\mathbf1,\mathbf1)_6.
\]
The common kernel of these representations of
$\SU(3)\times\SU(2)\times\U(1)_y$ is generated by
\[
\left(e^{2\pi i/3}I_3,-I_2,e^{i\pi/3}\right),
\]
hence is $\Z_6$.  This agrees with the quotient in
\eqref{eq:SM-gauge-group}.

\begin{proposition}[Hypercharge support can miss a gauge-charged exotic]
\label{cth:zero-hypercharge-exotic}
There are finite fermion carriers for which every target-sector anomaly and
every nonzero-hypercharge inventory test is unchanged, while the full
gauge-charged inventory residual of
Proposition~\ref{thm:canonical-chiral-projectors} is nonzero.  Adjoin to a passing
target packet a pair of weak adjoint triplets $(\mathbf1,\mathbf3)_0$,
one in each physical chirality.  These are representations of the same
$\Z_6$ quotient: the weak centre acts trivially in the adjoint representation,
and the colour and abelian factors act trivially.  The pair is invisible to
$\operatorname{supp}Y$ but spans a six-dimensional nontrivial weak sector.
Consequently, if these directions are included in the physical source,
\begin{equation}
 \rank P_{0Y}^{\rm nt}=6,\qquad \Delta_{\rm inv}^F=6,
 \label{eq:zero-hypercharge-exotic-gap}
\end{equation}
where $P_{0Y}^{\rm nt}$ is the zero-hypercharge subspace with the fully
gauge-trivial part removed.
\end{proposition}

\begin{proof}
The added triplets lie in $\Ker Y$ but not in the gauge-trivial sector, so
they are absent from $\operatorname{supp}Y$ and present in $P_{\rm gch}^F$.
Their total dimension is six.  Their zero hypercharge makes every local
anomaly term containing the abelian generator vanish; opposite physical
chiralities cancel the remaining local chiral contributions.  They add no
weak defining doublet.  Nevertheless, their irreducible gauge type differs
from every target type, so the full source-minimal inventory detects all six
directions as exotic.  The trivial action of the quotient kernel ensures
that this example does not change the represented global gauge group.
\end{proof}


\section{Fermion sources, occurrence, and allocation certificates}
\label{app:fermion-source}

These results make explicit the completeness and occurrence conditions
used in the charged-copy theorem.  They also separate physical source
support from ambient multiplicity, and process-level grading change from
the structure of a chosen generator.


\subsection{Complete fermion and typed sources}

\begin{proposition}[Finite certificate for completeness of the retained fermion source]
\label{prop:fermion-source-completeness}
Let $U_F:\C^M\to\Hh_F^{\rm amb}$ synthesize an independently complete
physical fermion-source atlas.  For an $\mathcal A_F$-reducing candidate
carrier $\mathcal K$ containing the retained source bank, define
\begin{equation}
 \mathbb C_{F\mid\mathcal K}
 :=U_F^*(I-P_{\mathcal K})U_F\succeq0.
 \label{eq:fermion-source-completeness}
\end{equation}
Then
\begin{equation}
 \mathbb C_{F\mid\mathcal K}=0
 \iff \Ran U_F\subseteq\mathcal K
 \iff \mathcal A_F\Ran U_F\subseteq\mathcal K.
 \label{eq:fermion-source-completeness-zero}
\end{equation}
Moreover,
\begin{align}
 \Tr\mathbb C_{F\mid\mathcal K}
 &=\|(I-P_{\mathcal K})U_F\|_{\HS}^2,\nonumber\\
 \rank\mathbb C_{F\mid\mathcal K}
 &=\dim\Ran U_F-\dim(\Ran U_F\cap\mathcal K).
 \label{eq:fermion-source-completeness-rank}
\end{align}
The zero test is exact; the trace measures source leakage and the rank counts
independent source directions modulo $\mathcal K$.  This rank need not equal
the dimension of their subsequent cyclic closure.
\end{proposition}

\begin{proof}
Writing $V=(I-P_{\mathcal K})U_F$ gives
$\mathbb C_{F\mid\mathcal K}=V^*V$, hence positivity, the trace identity, and
$\mathbb C_{F\mid\mathcal K}=0$ exactly when $\Ran U_F\subseteq\mathcal K$.
The latter inclusion is equivalent to
$\mathcal A_F\Ran U_F\subseteq\mathcal K$ because $\mathcal K$ is invariant
under $\mathcal A_F$ and $I\in\mathcal A_F$.  Finally, apply rank--nullity to
$(I-P_{\mathcal K})|_{\Ran U_F}$ and use $\rank(V^*V)=\rank V$.
\end{proof}

\noindent\textit{Exact completeness and quantitative separation.} A lower frame bound for $U_F$ on $\Ran U_F$ alone does not separate every
nonzero completeness residual from zero: the source range can meet
$\mathcal K$ at an arbitrarily small angle.  A uniform positive separation
requires additional geometry, for example $\mathcal K\subseteq\Ran U_F$.
Appendix~\ref{app:source-allocation-certificates} proves the contained-range
bound and gives an explicit counterexample without that inclusion.


\begin{lemma}[Complete multiplicity synthesis from the saturated fermion source]
\label{lem:complete-typed-source}
For a populated charged type $f$, set
$N_f=\mathcal I_{f,\epsilon_f}$ with the multiplicity inner product of
Proposition~\ref{thm:canonical-chiral-projectors}.  Its evaluation map
\[
 \operatorname{ev}_f:V_f\otimes N_f
 \longrightarrow P_{f,\epsilon_f}\mathcal K_F^{\min},\qquad
 \operatorname{ev}_f(v\otimes T)=Tv
\]
is unitary.  Let $Z_F:E_{\rm sat}\to\mathcal K_F^{\min}$ synthesize a
terminal source-word bank from
Proposition~\ref{thm:fermion-source-saturation}, so that
$\Ran Z_F=\mathcal K_F^{\min}$.  Define
\begin{equation}
 \begin{split}
 \widehat Z_f&=\operatorname{ev}_f^*P_{f,\epsilon_f}Z_F,\\
 S_f:\overline{V_f}\otimes E_{\rm sat}&\longrightarrow N_f,\qquad
 S_f(\overline v\otimes e)
       =(\langle v|\otimes I_{N_f})\widehat Z_fe.
 \end{split}
 \label{eq:complete-typed-source-synthesis}
\end{equation}
Then $\Ran S_f=N_f$.  Thus the faithful quotient of this typed source Gram
is precisely the canonical charged multiplicity, even if the displayed bank
has redundant columns.
\end{lemma}

\begin{proof}
For an orthonormal intertwiner basis, the identities
$T_a^*T_b=\delta_{ab}I_{V_f}$ make evaluation an isometry onto the complete
type block.  Since $Z_F$ is onto the saturated carrier, $\widehat Z_f$ is onto
$V_f\otimes N_f$.  Given a unit $v\in V_f$ and $n\in N_f$, choose $e$ with
$\widehat Z_fe=v\otimes n$.  Then
$S_f(\overline v\otimes e)=n$, proving surjectivity.  The conjugate space
$\overline{V_f}$ makes this contraction linear in its source coordinates.
\end{proof}


\subsection{Source-native allocation formulas}

\begin{corollary}[Source-native endpoint allocation]
\label{cor:source-native-endpoint-allocation}
Under Proposition~\ref{thm:endpoint-ancestry-compiler}, assume
$\Delta_f^{\rm anc}(\Xi_f)=0$ and $T_f^*T_f\succ0$.  Let
$S_f:E_f\to N_f$ be the complete typed synthesis of
Lemma~\ref{lem:complete-typed-source}.  Then every quantity in the endpoint
allocation short of Theorem~\ref{thm:endpoint-allocation} is determined by
$S_f$ and the represented ancestry operator $\Xi_f$ through the canonical
coefficient \eqref{eq:endpoint-ancestry-partial-trace}.  In particular,
\begin{equation}
a_{f\mid G}=0
 \quad\Longleftrightarrow\quad
 \Ran\Xi_f=P_{f,\epsilon_f}\mathcal K_F^{\min}
 \quad\Longleftrightarrow\quad n_f=3,
\label{eq:source-native-three-count}
\end{equation}
where the middle equality is understood after the canonical identification
$P_{f,\epsilon_f}\mathcal K_F^{\min}\cong V_f\otimes N_f$: namely
$\Ran\Xi_f=V_f\otimes N_f$.
\end{corollary}

\begin{proof}
The compiler supplies the unique faithful $T_f:G_{\rm gen}\to N_f$ required
by Theorem~\ref{thm:endpoint-allocation}.  Since
$\widehat\Xi_f=I_{V_f}\otimes T_f$, one has
$\Ran\Xi_f=V_f\otimes\Ran T_f$ under evaluation coordinates.  The zero
allocation residual is equivalent to $\Ran T_f=N_f$, and completeness of
$S_f$ then gives $n_f=3$ by
\eqref{eq:endpoint-allocation-zero}.  The converse implications are the same
identities in reverse.
\end{proof}

\begin{corollary}[Same-source coefficient formula]
\label{cor:endpoint-allocation-compiler}
Suppose the actual endpoint subsource is specified by a represented
coefficient map $J_{fG}:G_{\rm gen}\to E_f$ with
$T_f=S_fJ_{fG}$ and $J_{fG}^*G_fJ_{fG}\succ0$.  Then
\begin{equation}
 \mathbb A_{f\mid G}
 =G_f-G_fJ_{fG}(J_{fG}^*G_fJ_{fG})^{-1}J_{fG}^*G_f.
 \label{eq:endpoint-allocation-compiler}
\end{equation}
Thus the allocation residual, its trace, and its excess rank are determined
by the same typed source Gram and the specified endpoint compiler, with no
independent mixed Gram required.  If these data are word native at a flat
historical depth, the existing multiplication and Gram tables compute the
residual.
\end{corollary}

\begin{proof}
Use $A_f=J_{fG}^*G_fJ_{fG}$ and $C_f=J_{fG}^*G_f$ in
\eqref{eq:endpoint-allocation-short}.  The historical assertion follows from
Proposition~\ref{thm:inserted-word-moments} and finite matrix operations on the
reconstructed Gram.  The map $J_{fG}$ must specify the actual endpoint subsource.  On the
covariant same-carrier branch of
Proposition~\ref{thm:endpoint-ancestry-compiler}, it is replaced by the canonical
coefficient extracted from $\Xi_f$; choosing an arbitrary rank-three
coefficient subspace is not a source-occurrence argument.
\end{proof}


\subsection{Common-carrier source quotients}
\label{app:coupled-source}

The source-short argument also applies when the two banks are interpreted
as gravitational and internal variations.  It is a Gram-space statement;
it does not specify a variational action or its continuum limit.

Let
\[
S_{\rm g}:E_{\rm g}\to\Hh_{\rm cyl},
\qquad
S_{\rm i}:E_{\rm i}\to\Hh_{\rm cyl}
\]
be the source syntheses of the complete gravitational and internal physical
variation banks on one same-history cylinder.  Assume the gravitational bank
has already been reduced to its faithful quotient and put
\[
G_{\rm g}=S_{\rm g}^*S_{\rm g},
\qquad
G_{\rm i}=S_{\rm i}^*S_{\rm i},
\qquad
C=S_{\rm i}^*S_{\rm g}.
\]

\begin{proposition}[Universal coupled source quotient]
\label{thm:coupled-schur}
Define
\begin{equation}
G_{{\rm i}\mid{\rm g}}
=
G_{\rm i}-CG_{\rm g}^{-1}C^*
=
S_{\rm i}^*(I-P_{\rm g})S_{\rm i}.
\label{eq:coupled-schur}
\end{equation}
Then
\begin{equation}
\begin{pmatrix}
G_{\rm g}&C^*\\
C&G_{\rm i}
\end{pmatrix}
=
\begin{pmatrix}
I&0\\
CG_{\rm g}^{-1}&I
\end{pmatrix}
\begin{pmatrix}
G_{\rm g}&0\\
0&G_{{\rm i}\mid{\rm g}}
\end{pmatrix}
\begin{pmatrix}
I&G_{\rm g}^{-1}C^*\\
0&I
\end{pmatrix}.
\label{eq:coupled-factorization}
\end{equation}
Consequently:
\begin{enumerate}[label=\textup{(\roman*)},leftmargin=2.8em]
\item the joint variation carrier is faithful exactly when
      $G_{{\rm i}\mid{\rm g}}\succ0$;
\item the number of genuinely new internal source directions is
      $\rank G_{{\rm i}\mid{\rm g}}$;
\item the source-minimal internal variation space is
      \[
      E_{{\rm i}\mid{\rm g}}^{\min}
      =
      E_{\rm i}/\Ker G_{{\rm i}\mid{\rm g}}.
      \]
\end{enumerate}
\end{proposition}

\begin{proof}
The first equality in \eqref{eq:coupled-schur} is the Gram of
$(I-P_{\rm g})S_{\rm i}$.  Direct block multiplication gives
\eqref{eq:coupled-factorization}; the remaining statements are the standard
Schur-complement rank and positivity consequences.
\end{proof}

Thus ``adding the Standard Model'' does not mean appending every internal
variation as a new primitive.  Only the component orthogonal to the already
represented gravitational source counts as genuinely new.


\subsection{Completeness residuals and quantitative allocation thresholds}
\label{app:source-allocation-certificates}

\begin{example}[A source frame bound does not control an arbitrary transverse candidate]
\label{ex:source-angle-no-gap}
On $\C^3$ take $\mathcal A_F=\C I$, let the retained source be
$B_Fz=ze_1$, and let $U_F:\C^2\to\C^3$ be
$U_F(a,b)=ae_1+be_2$.  It has frame bound one on
$H_U=\Span\{e_1,e_2\}$.  For $0<\theta<\pi/2$, the carrier
\[
 \mathcal K_\theta
 =\Span\{e_1,\cos\theta\,e_2+\sin\theta\,e_3\}
\]
reduces $\mathcal A_F$ and contains the retained source, but
\begin{equation}
 U_F^*(I-P_{\mathcal K_\theta})U_F
 =\begin{pmatrix}0&0\\0&\sin^2\theta\end{pmatrix}.
 \label{eq:source-angle-counterexample}
\end{equation}
Thus a nonzero completeness residual can have arbitrarily small norm while
the frame bound remains one.  Notice that
$\mathcal K_\theta\not\subseteq H_U$.
\end{example}

\begin{proposition}[Contained-source frame separation]
\label{prop:contained-source-frame-gap}
Let $U:E\to\mathcal H$ have range $H_U$ and satisfy
$UU^*\succeq mP_{H_U}$ with $m>0$.  If $K\subseteq H_U$, put
\[
 C_{U\mid K}=U^*(I-P_K)U,
 \qquad r=\dim H_U-\dim K.
\]
Then
\begin{equation}
 \rank C_{U\mid K}=r,\qquad
 \Tr C_{U\mid K}\ge mr.
 \label{eq:contained-source-frame-trace}
\end{equation}
If $r>0$, every positive eigenvalue of $C_{U\mid K}$ is at least $m$.
Consequently $C_{U\mid K}\ne0$ implies
$\|C_{U\mid K}\|_{\rm op}\ge m$; equivalently,
$\Tr C_{U\mid K}<m$ forces $r=0$ and $C_{U\mid K}=0$.
\end{proposition}

\begin{proof}
Since $K\subseteq H_U$, the operator $Q=P_{H_U}-P_K$ is the orthogonal
projection onto $H_U\cap K^\perp$, and $(I-P_K)U=QU$.  Thus
\[
 C_{U\mid K}=(QU)^*(QU),\qquad
 (QU)(QU)^*=QUU^*Q\succeq mQ.
\]
On $\Ran Q$ the latter operator is strictly positive; outside $\Ran Q$ it
vanishes.  It therefore has exactly $r$ positive eigenvalues, all at least
$m$.  The two finite Gram operators have the same nonzero spectrum, proving
the assertions.  The conclusion includes the zero-dimensional case.
\end{proof}

For the complete typed source of
Theorem~\ref{thm:endpoint-allocation}, one has $H_U=N_f$ and
$K=\Ran T_f\subseteq N_f$.  Hence an independently known bound
$S_fS_f^*\succeq m_fI_{N_f}$ yields
\begin{equation}
 a_{f\mid G}\ge m_f(n_f-3),\qquad
 a_{f\mid G}<m_f\Longrightarrow n_f=3.
 \label{eq:endpoint-allocation-frame-threshold}
\end{equation}
In particular, if a measured scalar $\widehat a$ satisfies
$|\widehat a-a_{f\mid G}|\le\varepsilon_a$ and
$\widehat a+\varepsilon_a<m_f$, the exact three-dimensional allocation
follows.  This implication requires the frame bound and completeness to be
verified independently of the allocation measurement.

\begin{example}[Small allocation mass without a uniform source floor]
\label{ex:allocation-small-extra-source}
Let $N_f=\C^4$, take
$T_f:\C^3\to\C^4$ to be the inclusion of the first three coordinate axes,
and put
\[
 S_{f,t}=\diag(1,1,1,t),\qquad 0<t\le1.
\]
Every $S_{f,t}$ is a complete source for $N_f$, yet
\begin{equation}
 \mathbb A_{f\mid G,t}=\diag(0,0,0,t^2),\qquad
 a_{f\mid G,t}=t^2,\qquad n_f=4.
 \label{eq:allocation-small-extra-source}
\end{equation}
Thus small positive allocation mass cannot alone certify exactly three
multiplicity directions.  The least source frame bound is $m_f=t^2$, so the
quantitative estimate \eqref{eq:endpoint-allocation-frame-threshold} is sharp.
\end{example}

\noindent\textit{Source range and physical multiplicity.} For an arbitrary synthesis $S:E\to\mathcal H$ and faithful
$T:G_{\rm gen}\to\mathcal H$ with $\Ran T\subseteq\Ran S$, the same proof
as Theorem~\ref{thm:endpoint-allocation} gives
\[
 \rank(S^*S)
 =3+\rank\bigl(S^*(I-P_{\Ran T})S\bigr).
\]
This identity counts the source range.  It counts the canonical charged
multiplicity only when that range is its complete typed realization.
Lemma~\ref{lem:complete-typed-source} provides precisely this identification
from a saturated physical fermion bank.  An incomplete coefficient bank, an
unread Kraus index, or a source on an unrelated carrier cannot replace that
identification.


\subsection{Kraus-independent grading-changing occurrence}
\label{app:odd-occurrence}

The Clifford and internal packets must occur on one carrier.  Let
$J=P_+-P_-$ be the protected pre-central grading and let
$\sigma_0,\ldots,\sigma_3$ be Hermitian Clifford axes generating the external
matrix factor.  Before selecting a coherent grading-odd generator, primitive
matter occurrence can be tested directly at the completely positive process
level.

For a completely positive map
\[
\Phi(X)=\sum_\alpha K_\alpha X K_\alpha^*
\]
on the active carrier, let $J_\Phi$ be its Choi matrix.  On Hilbert--Schmidt
operator space define
\[
\Theta_J(X)=JXJ,
\qquad
\mathbf P_{\rm odd}=\frac12(I-\Theta_J).
\]

\begin{proposition}[Kraus-independent odd-Choi occurrence]
\label{thm:odd-choi-occurrence}
The quantity
\begin{equation}
\mu_J(\Phi):=\Tr(\mathbf P_{\rm odd}J_\Phi)
\label{eq:odd-choi-mass}
\end{equation}
is independent of the Kraus representation and satisfies
\begin{equation}
\begin{aligned}
\mu_J(\Phi)
&=\sum_\alpha\norm{(K_\alpha)_{\rm odd}}_{\HS}^2\\
&=\Tr[P_-\Phi(P_+)]+\Tr[P_+\Phi(P_-)]\\
&=\sum_\alpha\left(
\norm{P_-K_\alpha P_+}_{\HS}^2
+\norm{P_+K_\alpha P_-}_{\HS}^2\right),
\end{aligned}
\label{eq:odd-choi-identities}
\end{equation}
where $(K_\alpha)_{\rm odd}=(K_\alpha-JK_\alpha J)/2$.
Moreover,
\begin{equation}
\mu_J(\Phi)=0
\quad\Longleftrightarrow\quad
\Ran J_\Phi\ \text{is supported on grading-even operators}.
\label{eq:odd-choi-zero}
\end{equation}
Equivalently, every Kraus direction commutes with $J$.  If
$d_\pm=\Tr P_\pm>0$ and $\rho_\pm=P_\pm/d_\pm$, then
\begin{equation}
\mu_J(\Phi)
=d_+\Tr[P_-\Phi(\rho_+)]
+d_-\Tr[P_+\Phi(\rho_-)].
\label{eq:odd-choi-two-sign}
\end{equation}
Thus primitive grading change is decidable from two positive cross-sign
population experiments without choosing Kraus phases.
\end{proposition}

\begin{proof}
Vectorization gives
$J_\Phi=\sum_\alpha|K_\alpha\rangle\!\rangle
\langle\!\langle K_\alpha|$.  Taking the trace against the orthogonal
projector $\mathbf P_{\rm odd}$ gives the first line of
\eqref{eq:odd-choi-identities} and proves Kraus-coordinate independence.
For one Kraus operator,
\[
\Tr(P_-KP_+K^*)=\norm{P_-KP_+}_{\HS}^2,
\qquad
\Tr(P_+KP_-K^*)=\norm{P_+KP_-}_{\HS}^2.
\]
The two off-diagonal corners are Hilbert--Schmidt orthogonal and sum to the
odd part of $K$, proving the remaining identities.  Since every summand is
nonnegative, the mass vanishes exactly when the Choi support lies in the even
operator subspace.  Substituting $P_\pm=d_\pm\rho_\pm$ gives
\eqref{eq:odd-choi-two-sign}.
\end{proof}

\begin{corollary}[Primitive active no-matter or occurrence alternative]
\label{cor:primitive-odd-alternative}
Let $\Sigma$ be a finite alphabet of actual autonomous CP letters and put
\begin{equation}
\mathfrak m_{\rm odd}(\Sigma)
:=\sum_{a\in\Sigma}\mu_J(\Phi_a).
\label{eq:primitive-odd-profile}
\end{equation}
Assume the autonomous word grammar is generated by these letters together with
grading-even passive writers and initializations.  Then
\begin{equation}
\mathfrak m_{\rm odd}(\Sigma)=0
\quad\Longleftrightarrow\quad
\text{every finite autonomous resolved word is grading even}.
\label{eq:primitive-odd-stop}
\end{equation}
If $\mathfrak m_{\rm odd}(\Sigma)>0$, at least one represented primitive
letter has nonzero cross-sign support.  This is a process-level occurrence
statement; identifying a particular one-route weak factor may require a
separate selector when the corresponding flip Choi support has rank greater
than one.
\end{corollary}

\begin{proof}
A finite sum of nonnegative odd masses vanishes exactly when every primitive
Choi support is grading even by Proposition~\ref{thm:odd-choi-occurrence}.
Products of grading-even Kraus directions remain grading even, and the stated
writers and initializations do not change parity, proving the forward
implication.  The converse follows already from words of length one.  Strict
positivity of the sum forces at least one positive primitive summand.
\end{proof}


\subsection{Active grading change and unread spectators}

Let $Z=Z^*=Z^{-1}$ be a grading on a future-visible active carrier
$\Hh_{\rm act}$, with $P_\pm=(I\pm Z)/2$.  For a completely positive map
$\Phi$ on that carrier define its active cross-sign mass by
\begin{equation}
\mu_Z^{\rm act}(\Phi)
=
\Tr\!\left[P_-\Phi(P_+)\right]
+
\Tr\!\left[P_+\Phi(P_-)\right].
\label{eq:active-cross-sign-mass}
\end{equation}
This quantity asks only whether the represented active map transfers positive
weight between the two grading sectors.

\begin{proposition}[Spectator-lift invariance of active grading change]
\label{thm:active-odd-lift-invariance}
Let $\mathcal S$ be an unread spectator, let $\sigma$ be a spectator state, and
let $\Psi$ be any trace-preserving completely positive map on $\mathcal S$.
For the product lift
\[
\widetilde\Phi=\Phi\otimes\Psi
\]
and active preparations $P_\pm\otimes\sigma$ with active Reads
$P_\mp\otimes I_{\mathcal S}$, one has
\begin{equation}
\begin{aligned}
&\Tr\!\left[(P_-\otimes I)\widetilde\Phi(P_+\otimes\sigma)\right]
+\Tr\!\left[(P_+\otimes I)\widetilde\Phi(P_-\otimes\sigma)\right]\\
&\hspace{10em}=\mu_Z^{\rm act}(\Phi).
\end{aligned}
\label{eq:active-odd-lift-invariance}
\end{equation}
Thus active grading-changing occurrence is invariant under arbitrary unread
trace-preserving spectator dynamics.  By contrast, the total grading-odd mass
of a chosen microscopic dilation need not be operationally identifiable; an
explicit continuous family with identical visible dynamics is given in
Appendix~\ref{app:microscopic-odd-lift}.
\end{proposition}

\begin{proof}
Trace preservation gives $\Tr\Psi(\sigma)=1$.  Hence
\[
\Tr\!\left[(P_-\otimes I)(\Phi\otimes\Psi)(P_+\otimes\sigma)\right]
=
\Tr[P_-\Phi(P_+)]\,\Tr\Psi(\sigma)
=
\Tr[P_-\Phi(P_+)],
\]
and the same argument applies to the reverse sign.  Adding the two identities
gives \eqref{eq:active-odd-lift-invariance}.
\end{proof}

\noindent\textit{Hidden provenance is not active matter.} A passive Read of a future-sealed hidden sector may establish microscopic
provenance, but it does not make that sector part of the active matter packet.
For hidden structure to contribute to a typed matter incidence it must return
to the active carrier through a represented router.  Such a return map is
additional physical loading data; it cannot be supplied by an unread spectator
or inferred from unchanged visible record marginals.


\subsubsection{A microscopic odd-mass ambiguity}
\label{app:microscopic-odd-lift}

\begin{proposition}[Microscopic odd mass is not fixed by visible dynamics]
\label{cth:microscopic-odd-lift}
Fix a visible carrier of dimension $d\ge1$ and take its visible primitive map
to be the identity channel.  Let an unread spectator be $\mathcal S=\C^2$ with
grading $Z_{\mathcal S}=\sigma_z$.  For $0\le\theta\le\pi/2$ put
\[
U_\theta=\cos\theta\,I+i\sin\theta\,\sigma_x,
\qquad
\Psi_\theta=\operatorname{Ad}_{U_\theta}.
\]
Every lifted map $\operatorname{id}\otimes\Psi_\theta$ induces the same
visible identity channel, and every visible multitime word is unchanged after
the spectator is discarded.  Nevertheless, for any visible grading
$Z_{\rm act}$ and the total grading
$\widetilde Z=Z_{\rm act}\otimes Z_{\mathcal S}$, the odd Choi mass of the
lift is
\begin{equation}
\mu_{\widetilde Z}
 (\operatorname{id}\otimes\Psi_\theta)
=2d\sin^2\theta.
\label{eq:microscopic-odd-lift-family}
\end{equation}
Thus operationally indistinguishable lifts can have continuously different
microscopic odd mass; by enlarging the unread spectator multiplicity this
quantity can be increased without changing the visible process.
\end{proposition}

\begin{proof}
Discarding a trace-preserving spectator leaves the visible identity channel at
every word depth.  Relative to $Z_{\mathcal S}$,
\[
(U_\theta)_{\rm ev}=\cos\theta\,I,
\qquad
(U_\theta)_{\rm odd}=i\sin\theta\,\sigma_x.
\]
The unique Kraus operator of the lifted unitary channel is
$I_d\otimes U_\theta$.  Its odd part relative to $\widetilde Z$ is
$I_d\otimes(U_\theta)_{\rm odd}$, hence
\[
\norm{I_d\otimes(U_\theta)_{\rm odd}}_{\HS}^2
=d\,\norm{i\sin\theta\,\sigma_x}_{\HS}^2
=2d\sin^2\theta.
\]
Proposition~\ref{thm:odd-choi-occurrence} identifies this squared norm with the
odd Choi mass.  The visible active mass, by contrast, remains zero by
Proposition~\ref{thm:active-odd-lift-invariance}.
\end{proof}


\subsubsection{The narrower Clifford-occurrence margin}

The quantity in \eqref{eq:odd-choi-identities} is the active cross-sign mass
of \eqref{eq:active-cross-sign-mass}; hence
Proposition~\ref{thm:active-odd-lift-invariance} shows that this primitive
occurrence test is unchanged by unread trace-preserving spectator lifts.
The Clifford margin below is a narrower test: it asks whether the represented
external Clifford action itself changes the same grading.

\begin{definition}[Clifford matter margin]
\label{def:clifford-matter}
Define
\begin{equation}
p_{\rm Cl}
=
\frac{1}{2\dim\Hh}
\sum_{\mu=0}^{3}
\norm{P_-\sigma_\mu P_+}_{\HS}^2.
\label{eq:clifford-matter}
\end{equation}
\end{definition}

\begin{proposition}[Clifford occurrence test]
\label{prop:clifford-occurrence}
One has $p_{\rm Cl}>0$ exactly when at least one represented Clifford axis
changes the protected grading.  If $p_{\rm Cl}=0$, the grading reduces the
external Clifford action and its residual action lies in the normal
multiplicity commutant.
\end{proposition}

\begin{proof}
Every summand in \eqref{eq:clifford-matter} is nonnegative.  The sum vanishes
exactly when $P_-\sigma_\mu P_+=0$ for every $\mu$, and self-adjointness then
also gives $P_+\sigma_\mu P_-=0$.  Thus $J$ commutes with the generated
external algebra.
\end{proof}


\subsection{Weak gauge-cyclic support and a rank-one counterexample}
\label{app:weak-cyclic-example}

For completeness, the range step in
Proposition~\ref{thm:weak-copy-census} can be checked directly.  For any
$x\in W_2\otimes M_H$,
\[
 \langle x,\overline J_Hx\rangle
 =\int_{\SU(W_2)}\|J_H^{1/2}(g^*\otimes I)x\|^2\,\dd g.
\]
The integrand is continuous and nonnegative and Haar measure has full
support.  The integral therefore vanishes exactly when every integrand
vanishes.  Since $\Ran J_H^{1/2}=\Ran J_H$ in finite dimension,
\[
 \Ker\overline J_H
 =\bigcap_g((g\otimes I)\Ran J_H)^\perp
 =(\mathcal K_H^{\rm occ})^\perp.
\]
Taking orthogonal complements gives the claimed gauge-cyclic range.

Let $(w_1,w_2)$ and $(m_1,m_2)$ be orthonormal bases of $W_2$ and $M_H$,
and put
\[
 \psi=\frac{w_1\otimes m_1+w_2\otimes m_2}{\sqrt2},\qquad
 J_H=|\psi\rangle\langle\psi|.
\]
Then $\rank J_H=1$ but
\[
 B_H=\frac12I_{M_H},\qquad
 \overline J_H=\frac14 I_{W_2\otimes M_H},\qquad
 \mathfrak p_H=\frac12.
\]
Proposition~\ref{thm:weak-copy-census} therefore gives
$\mathcal K_H^{\rm occ}=W_2\otimes M_H$ and $g_H^{\rm occ}=2$.
The number of positive coefficient vectors is not the number of reached
weak-doublet copies; the latter is the multiplicity of the gauge-cyclic
support.


\subsection{One-generation incidence and component structure}
\label{app:one-generation}

The commutant language also gives a compact check of how Yukawa incidence
reduces residual scalar symmetries inside one generation.

Let
\[
\mathcal F_{15}=Q\oplus u\oplus d\oplus L\oplus e
\]
with the modules of \eqref{eq:SM-types}.  Ignoring generation multiplicity for
the moment, the gauge-generated algebra is
\begin{equation}
\A_{\rm gauge}
=
M_6(\C)\oplus M_3(\C)\oplus M_3(\C)\oplus M_2(\C)\oplus\C,
\label{eq:one-gen-gauge-algebra}
\end{equation}
because the five modules are irreducible and pairwise inequivalent.  Hence
\[
\A_{\rm gauge}'\cong\C^5.
\]

\begin{proposition}[Incidence reduces the one-generation commutant]
\label{prop:one-gen-incidence}
Assume the three Higgs-contracted incidence histories
\[
Q\leftrightarrow u,\qquad
Q\leftrightarrow d,\qquad
L\leftrightarrow e
\]
are nonzero.  Then
\[
\A_{\rm gauge+inc}
=
M_{12}(\C)\oplus M_3(\C),
\qquad
\A_{\rm gauge+inc}'\cong\C^2.
\]
The two residual central directions are the quark and lepton connected
components.  Adding a neutral singlet and a nonzero $L$--neutral incidence
changes the second matrix block to $M_4(\C)$ while leaving the residual
commutant $\C^2$.
\end{proposition}

\begin{proof}
A nonzero block between two full matrix summands, together with multiplication
by endpoint matrix units, generates every rank-one map between the summands.
Products and adjoints therefore generate the full matrix algebra on each
connected component of the incidence graph.  The three displayed edges have
components $\{Q,u,d\}$ and $\{L,e\}$, of dimensions $12$ and $3$.
\end{proof}

This audit is intentionally distinct from the generation multiplicity theorem.
It describes the commutant inside one fixed generation packet.  The
rank-three generation carrier of Proposition~\ref{thm:geometry-short}
is an additional multiplicity factor reconstructed before numerical flavour
breaking is specified.


For the component ancestry hypotheses of
Proposition~\ref{thm:common-generation-factor}, the corresponding explicit
three-copy decomposition is
\[
 \mathcal H_{\rm mat}\cong
 (V_Q\oplus V_u\oplus V_d)\otimes G_q
 \oplus(V_L\oplus V_e[\oplus V_\nu])\otimes G_\ell,
 \qquad\dim G_q=\dim G_\ell=3.
\]
The neutral-singlet term requires its own represented ancestry hypotheses;
its multiplicity is not fixed by the charged anomaly equations.


\section{Stability estimates and compact-Mosco proofs}
\label{app:stability}

This appendix gives the incidence/flavour decomposition and coefficient
perturbation estimates, followed by the complete varying-space spectral
argument.  The cofinal kernel theorem in the main text does not require
the compactness hypotheses used in the latter argument.


\subsection{Type-central connectivity and conditional flavour rigidity}
\label{app:incidence-flavour-stability}

\begin{proposition}[Type-central incidence graph Laplacian]
\label{prop:type-central-graph-laplacian}
Let $P_v$ be the represented orthogonal type projections, put
$h_v=\rank P_v$, and let typed incidences satisfy
$Y_e=P_{t(e)}Y_eP_{s(e)}$.  Set
$w_e=\norm{Y_e}_{\HS}^2$, let $B$ be the oriented graph incidence matrix, and
define
\[
 H=\operatorname{diag}(h_v),
 \qquad
 \mathcal L_{\rm cen}
 =H^{-1/2}B^*\operatorname{diag}(2w_e)BH^{-1/2}.
\]
For the type-central operator $X_z=\sum_vz_vP_v$,
\begin{equation}
 \sum_e\left(
   \norm{[X_z,Y_e]}_{\HS}^2+
   \norm{[X_z,Y_e^*]}_{\HS}^2
 \right)
 =2\sum_e w_e|z_{t(e)}-z_{s(e)}|^2.
 \label{eq:type-central-graph-energy}
\end{equation}
If the positive-weight incidence graph has connected components
$C_1,\ldots,C_c$, then
\begin{equation}
 \Ker\mathcal L_{\rm cen}
 =H^{1/2}\Span\{\mathbf1_{C_1},\ldots,\mathbf1_{C_c}\},
 \label{eq:type-central-graph-kernel}
\end{equation}
and this kernel corresponds exactly to the block-scalar centre
$\Span\{P_{C_a}:1\le a\le c\}$, where
$P_C=\sum_{v\in C}P_v$.  The Hilbert--Schmidt projection of $X_z$ onto that
centre replaces the coefficients on each component by the weighted average
\[
 \bar z_C=\frac{\sum_{v\in C}h_vz_v}{\sum_{v\in C}h_v}.
\]
Consequently the first positive eigenvalue of $\mathcal L_{\rm cen}$ is the
exact type-central rigidity gap of the connected-component centre.
\end{proposition}

\begin{proof}
On a typed edge $e:s\to t$,
$[X_z,Y_e]=(z_t-z_s)Y_e$ and
$[X_z,Y_e^*]=(z_s-z_t)Y_e^*$.  Summing the squared norms proves
\eqref{eq:type-central-graph-energy}.  The energy vanishes exactly when $z$ is
constant on every positive-weight connected component, proving
\eqref{eq:type-central-graph-kernel}.  Orthogonality of the type projections
and $\norm{P_v}_{\HS}^2=h_v$ gives the weighted component-average projector and
the normalized spectral-gap statement.
\end{proof}

\begin{corollary}[Incidence--flavour factorization of the Howe certificate]
\label{cor:howe-incidence-flavour-factorization}
Assume on a structural residue branch that each represented incidence factors
as $Y_e=D_e\otimes F_e$ on its typed module and generation coefficient spaces.
On the block-diagonal multiplicity test space, decompose the tested
complement of the protected algebra into orthogonal type-central contrasts
and noncentral flavour directions.  Put $s=s(e)$, $t=t(e)$ and
$\kappa_e=\norm{D_e}_{\HS}^2$.  Retaining both the incidence and its adjoint,
define
\begin{align*}
 (\mathscr D_{\rm cen}z)_e
 &=\sqrt{\kappa_e}\bigl((z_t-z_s)F_e,\ (z_s-z_t)F_e^*\bigr),\\
 (\mathscr D_{\rm fl}X)_e
 &=\sqrt{\kappa_e}\bigl(X_tF_e-F_eX_s,\ X_sF_e^*-F_e^*X_t\bigr).
\end{align*}
Then the relative Howe Gram restricted to this test space has block form
\[
 A_{\rm cen}=\mathscr D_{\rm cen}^*\mathscr D_{\rm cen},
 \qquad
 B=\mathscr D_{\rm cen}^*\mathscr D_{\rm fl},
 \qquad
 C_{\rm fl}=\mathscr D_{\rm fl}^*\mathscr D_{\rm fl}+C_{\rm extra},
\]
where $C_{\rm extra}\succeq0$ contains the non-incidence flavour
commutators.  When $A_{\rm cen}\succ0$ on the declared type-central contrast
space,
\begin{equation}
 S_{\rm fl\mid cen}
 =C_{\rm extra}
 +\mathscr D_{\rm fl}^*(I-P_{\rm cen})\mathscr D_{\rm fl},
 \qquad
 P_{\rm cen}=\mathscr D_{\rm cen}
 (\mathscr D_{\rm cen}^*\mathscr D_{\rm cen})^{-1}
 \mathscr D_{\rm cen}^*.
 \label{eq:howe-factorized-Schur}
\end{equation}
On this test space, exact rigidity separates into incidence connectivity and
conditional flavour rigidity.  In particular, if deleting a displayed incidence edge would
create an additional type-central commutant direction excluded from the
proposed protected algebra, then
$\mathbb G_{\rm Howe}\succ0$ already forces that edge to occur.
\end{corollary}

\begin{proof}
For a type-central contrast, the commutator with $Y_e$ is the scalar endpoint
difference multiplied by $Y_e$; for a noncentral multiplicity variation it is
the coefficient defect $X_{t(e)}F_e-F_eX_{s(e)}$.  The adjoint incidence
contributes $X_{s(e)}F_e^*-F_e^*X_{t(e)}$ independently on complex operator
space.  Summing both Hilbert--Schmidt pairings gives the three blocks of the
restricted Gram.  The displayed Schur complement is the standard identity
$C-B^*A^{-1}B$ with
$P_{\rm cen}=\mathscr D_{\rm cen}A^{-1}\mathscr D_{\rm cen}^*$.  The final
claim follows from Theorem~\ref{thm:howe-certificate}: a missing edge would
leave an additional exact central zero mode outside the proposed protected
kernel.
\end{proof}


\subsection{Perturbed coefficient spaces and one-sided rank certification}
\label{app:coefficient-stability}

\begin{proposition}[Stable dominant coefficient space under incidence error]
\label{thm:incidence-slice-cluster}
Let an exact incidence $Y$ have operator--Schmidt rank $r$, smallest nonzero
Schmidt coefficient $s>0$, and largest coefficient $L$.  Suppose
$\|\widetilde Y-Y\|_{\HS}\le\varepsilon$ and let
$H_Y,H_{\widetilde Y}$ be the slice Grams of
Proposition~\ref{prop:incidence-slice-Gram}.  Then
\begin{equation}
 |\widetilde\sigma_k-\sigma_k|\le\varepsilon.
 \label{eq:incidence-Schmidt-Weyl}
\end{equation}
Hence $\varepsilon<s/2$ separates an $r$-dimensional dominant coefficient
cluster.  More quantitatively, put
\[
 \eta=(2L+\varepsilon)\varepsilon.
\]
If $\eta<s^2/2$, and
\[
 P=\mathbf1_{(0,\infty)}(H_Y),\qquad
 \widetilde P=\mathbf1_{[s^2/2,\infty)}(H_{\widetilde Y}),
\]
then $\rank\widetilde P=r$ and
\begin{equation}
 \norm{\widetilde P-P}_{\rm op}
 \le\frac{\eta}{s^2-\eta}.
 \label{eq:incidence-slice-projector}
\end{equation}
Thus a separated operator--Schmidt block yields a noise-stable coefficient
subspace even though literal exact slice rank need not be stable.
\end{proposition}

\begin{proof}
The slice map is linear in $Y$ and
$\|\mathcal S_{\widetilde Y}-\mathcal S_Y\|_{\rm op}
 \le\|\widetilde Y-Y\|_{\HS}$, so singular-value perturbation gives
\eqref{eq:incidence-Schmidt-Weyl}.  Also
\[
 \norm{H_{\widetilde Y}-H_Y}_{\rm op}
 \le(2L+\varepsilon)\varepsilon=\eta.
\]
The exact positive spectrum of $H_Y$ lies in $[s^2,\infty)$, while the
remaining spectrum is zero.  Weyl's eigenvalue bound therefore gives exactly
$r$ eigenvalues of $H_{\widetilde Y}$ above $s^2/2$ when $\eta<s^2/2$.
With $Q=I-P$, one has $QH_Y=0$ and
$H_{\widetilde Y}|_{\Ran\widetilde P}\succeq(s^2-\eta)I$, whence
\[
 Q\widetilde P
 =Q(H_{\widetilde Y}-H_Y)\widetilde P
   (H_{\widetilde Y}|_{\Ran\widetilde P})^{-1}.
\]
Taking norms and using equality of the two projector ranks proves
\eqref{eq:incidence-slice-projector}.
\end{proof}

\begin{corollary}[One-sided operator--Schmidt rank certificate]
\label{cor:incidence-one-sided-rank}
If an observed incidence $\widehat Y$ obeys
$\|\widehat Y-Y\|_{\HS}\le\varepsilon$, then
\begin{equation}
 \sigma_k(Y)\ge(\widehat\sigma_k-\varepsilon)_+,\qquad
 \rank Y\ge\#\{k:\widehat\sigma_k>\varepsilon\}.
 \label{eq:incidence-one-sided-rank}
\end{equation}
Without an independent upper-rank premise, a coefficient below the error floor
cannot certify exact absence of that Schmidt direction.
\end{corollary}

\begin{proof}
Apply \eqref{eq:incidence-Schmidt-Weyl}.  The final caveat is sharp because an
orthogonal Schmidt term of arbitrarily small norm can be added whenever the
ambient tensor spaces admit one more coefficient direction.
\end{proof}


\subsection{Compact-Mosco resolvent convergence}
\label{app:compact-mosco}

\begin{proposition}[Compact Mosco convergence implies norm-resolvent convergence]
\label{prop:compact-mosco-resolvent}
If $q_X\to q_\infty$ in the sense of
Definition~\ref{def:transported-mosco} and the family is asymptotically
compact in the sense of Definition~\ref{def:asymptotic-compactness}, then
\begin{equation}
\norm{R_X-R_\infty}_{\rm op}\longrightarrow0.
\label{eq:norm-resolvent}
\end{equation}
Consequently, if
$0\le\lambda_{1,X}\le\lambda_{2,X}\le\cdots$ denote the discrete screened
eigenvalues of $\mathscr L_X$, repeated with multiplicity, then for every
fixed index in the limiting discrete spectrum,
\begin{equation}
\lambda_{k,X}\longrightarrow\lambda_{k,\infty}.
\label{eq:eigenvalue-convergence}
\end{equation}
If an isolated spectral cluster of $\mathscr L_\infty$ is separated by a
positive gap, the corresponding transported Riesz projections converge in
operator norm and their ranks are eventually constant.
\end{proposition}

We establish the resolvent conclusion through its variational formulation
and a collective-compactness argument, and then treat the low spectrum.

Retain the embeddings $J_X:\mathscr H_X\to\mathscr H$ and the closed
nonnegative forms $q_X$ of Definition~\ref{def:transported-mosco} and Definition~\ref{def:asymptotic-compactness}.  For $f\in\mathscr H$ let
\begin{equation}
u_X=(\mathscr L_X+I)^{-1}J_X^*f.
\label{eq:resolvent-minimizer}
\end{equation}
Equivalently, $u_X$ is the unique minimizer on $D(q_X)$ of
\begin{equation}
\mathcal F_{X,f}(u)
=q_X(u)+\norm{u}^2-2\operatorname{Re}\ip{J_Xu}{f}_{\mathscr H}.
\label{eq:resolvent-functional}
\end{equation}
Writing $q_X(\cdot,\cdot)$ for the sesquilinear form obtained by polarizing
$q_X$, the Euler identity is
\begin{equation}
q_X(u_X,v)+\ip{u_X}{v}
=\ip{f}{J_Xv}_{\mathscr H},
\qquad v\in D(q_X).
\label{eq:resolvent-euler}
\end{equation}
Taking $v=u_X$ gives
\begin{equation}
\norm{u_X}^2+q_X(u_X)
\le \norm{f}\,\norm{u_X}
\le \norm{f}^2,
\qquad
\norm{u_X}\le\norm f.
\label{eq:resolvent-energy-bound}
\end{equation}

\begin{lemma}[Mosco convergence gives strong transported resolvents]
\label{lem:mosco-strong-resolvent}
Under Definition~\ref{def:transported-mosco}, for every fixed
$f\in\mathscr H$,
\begin{equation}
J_X(\mathscr L_X+I)^{-1}J_X^*f
\longrightarrow
J_\infty(\mathscr L_\infty+I)^{-1}J_\infty^*f
\label{eq:strong-resolvent}
\end{equation}
strongly in $\mathscr H$.
\end{lemma}

\begin{proof}
The estimate \eqref{eq:resolvent-energy-bound} makes $(J_Xu_X)$ bounded.
Every weakly convergent subsequence has, by the Mosco liminf inequality, a
limit $J_\infty u$ satisfying
\[
q_\infty(u)+\norm u^2-2\operatorname{Re}\ip{J_\infty u}{f}
\le\liminf_X\mathcal F_{X,f}(u_X).
\]
For every $v\in D(q_\infty)$ choose a Mosco recovery sequence $v_X$.
Minimality of $u_X$ gives
$\mathcal F_{X,f}(u_X)\le\mathcal F_{X,f}(v_X)$; passing to the limit shows
that $u$ minimizes $\mathcal F_{\infty,f}$.  Strict convexity identifies
$u=(\mathscr L_\infty+I)^{-1}J_\infty^*f$ and therefore identifies every weak
cluster point.  Comparing the minimum values and using the liminf inequality
gives convergence of $\norm{u_X}^2+q_X(u_X)$ to the limiting value.  Weak
convergence together with convergence of the Hilbert norms yields strong
convergence of $J_Xu_X$.  Thus the whole sequence converges.
\end{proof}

\begin{lemma}[Asymptotic compactness gives collective compactness]
\label{lem:collective-compactness}
Assume Definition~\ref{def:asymptotic-compactness}.  If $(f_X)$ is bounded in
$\mathscr H$, then every sequence
\[
J_X(\mathscr L_X+I)^{-1}J_X^*f_X
\]
has a strongly convergent subsequence.  Hence the transported resolvents
$(R_X)$ form a collectively compact family.  The limiting resolvent $R_\infty$
is compact.
\end{lemma}

\begin{proof}
Apply \eqref{eq:resolvent-energy-bound} with $f=f_X$.  The corresponding
resolvent vectors have uniformly bounded norms and form energies, so
asymptotic compactness gives a strongly convergent subsequence.  The same
argument for the limiting form, applied to an arbitrary bounded sequence of
right-hand sides, proves compactness of $R_\infty$.
\end{proof}

\begin{lemma}[Collectively compact self-adjoint convergence is norm convergence]
\label{lem:collective-norm}
Let $T_X$ be bounded self-adjoint operators on a Hilbert space.  If
$T_X\to T$ strongly and $(T_X)$ is collectively compact, then $T$ is compact
and
\begin{equation}
\norm{T_X-T}_{\rm op}\longrightarrow0.
\label{eq:collective-norm}
\end{equation}
\end{lemma}

\begin{proof}
Let $B_1$ be the closed unit ball and let
\[
 K=\overline{\bigcup_X T_X(B_1)}.
\]
Collective compactness means that $K$ is compact.  For every fixed
$x\in B_1$, strong convergence gives $T_Xx\to Tx$, while every $T_Xx$ belongs
to $K$; hence $Tx\in K$.  Thus $T(B_1)\subseteq K$, so $T$ is compact.  This
argument uses strong convergence only at a fixed vector and avoids any
uniformity claim on moving vectors.

If \eqref{eq:collective-norm} failed, there would be a cofinal sequence of
indices, unit vectors $x_X$, and $\varepsilon>0$ with
$\norm{(T_X-T)x_X}\ge\varepsilon$.  Passing to a subsequence, weak compactness
of the unit ball gives $x_X\rightharpoonup x$, collective compactness gives
$T_Xx_X\to y$, and compactness of $T$ gives $Tx_X\to Tx$.  For every fixed
$z$, self-adjointness and strong convergence at the fixed vector $z$ give
\[
\ip{y}{z}
=\lim_X\ip{T_Xx_X}{z}
=\lim_X\ip{x_X}{T_Xz}
=\ip{x}{Tz}
=\ip{Tx}{z}.
\]
Thus $y=Tx$, so $(T_X-T)x_X\to0$, contradicting the choice of the sequence.
\end{proof}

The strong-resolvent lemma, the collective-compactness lemma, and the
self-adjoint norm-convergence lemma prove
$\|R_X-R_\infty\|_{\rm op}\to0$.  The next lemma completes the spectral
claims of Proposition~\ref{prop:compact-mosco-resolvent}.


\subsection{Low spectrum and the alternative protected-kernel criterion}

Let $\mu_{1,X}\ge\mu_{2,X}\ge\cdots>0$ denote the nonzero eigenvalues of the
compact positive operator $R_X$, repeated with multiplicity.  On the screened
space they are related to the eigenvalues of $\mathscr L_X$ by
\begin{equation}
\mu_{k,X}=\frac{1}{1+\lambda_{k,X}}.
\label{eq:resolvent-spectrum-map}
\end{equation}

\begin{lemma}[Low-spectrum convergence]
\label{lem:low-spectrum-convergence}
If $\norm{R_X-R_\infty}\to0$, then for every fixed $k$,
\begin{equation}
\mu_{k,X}\to\mu_{k,\infty},
\qquad
\lambda_{k,X}\to\lambda_{k,\infty}
\label{eq:low-spectrum-convergence}
\end{equation}
whenever the limiting eigenvalue is finite.  If $\Sigma$ is an isolated finite
spectral cluster of $R_\infty$ and $\Gamma$ is a positively oriented contour
separating it from the remaining spectrum, then the transported Riesz
projections
\begin{equation}
\Pi_X(\Gamma)
=\frac{1}{2\pi i}\int_\Gamma(z-R_X)^{-1}\,\dd z
\label{eq:riesz-projection}
\end{equation}
converge to $\Pi_\infty(\Gamma)$ in operator norm.  Their ranks are eventually
equal.
\end{lemma}

\begin{proof}
For compact self-adjoint operators the min--max principle gives
$|\mu_{k,X}-\mu_{k,\infty}|\le\norm{R_X-R_\infty}$ for each fixed $k$.
Equation~\eqref{eq:resolvent-spectrum-map} gives the second convergence.
For the projection statement, norm convergence and the resolvent identity
imply uniform convergence of $(z-R_X)^{-1}$ on $\Gamma$; integrate around the
contour.  If two projections differ in operator norm by less than one, their
ranges have the same finite dimension, proving eventual rank stability.
\end{proof}

\begin{proposition}[Protected-kernel locking]
\label{prop:protected-kernel-locking}
Suppose each $\M_X$ is an $m$-dimensional subspace of
$\Ker\mathscr L_X$ and
$\Ker\mathscr L_\infty=\M_\infty$ has dimension $m$.  Under compact Mosco
convergence, for all sufficiently late $X$,
\[
\Ker\mathscr L_X=\M_X,
\]
the transported kernel projections converge in operator norm, and
\[
\lambda_{m+1}(\mathscr L_X)
\to\lambda_{m+1}(\mathscr L_\infty)>0.
\]
\end{proposition}

\begin{proof}
The eigenvalue $0$ of $\mathscr L_\infty$ corresponds to the eigenvalue $1$
of $R_\infty$, of multiplicity $m$.  Compactness of the limiting resolvent and
finite-dimensionality of the kernel separate this eigenvalue from the rest of
the spectrum.  Lemma~\ref{lem:low-spectrum-convergence} therefore gives a
nearby spectral cluster of rank exactly $m$ for all sufficiently late $X$.
Because $\M_X\subseteq\Ker\mathscr L_X$ already supplies $m$ exact copies of
the eigenvalue $1$ of $R_X$, that entire cluster is the kernel and no extra
zero mode can occur.  The same lemma gives projection and next-eigenvalue
convergence.
\end{proof}

This is a distinct spectral criterion: it assumes an exact limiting
kernel and uses compact-Mosco convergence to lock the late finite kernels.
By contrast, Theorem~\ref{thm:cofinal-duality} derives limiting kernel
exactness from protected transport and a uniform finite complement gap,
without compactness.


\subsection{Why protected transport is indispensable}

\begin{proposition}[A uniformly gapped protected mode can escape]
\label{cth:escaping-protected-mode}
Fixed protected dimension, transported Mosco convergence, and a uniform
finite-cutoff complement gap do not by themselves produce a nonzero protected
limiting kernel.  On $\ell^2(\N)$ let $e_n$ be the standard basis,
$\M_n=\C e_n$, $P_n=|e_n\rangle\langle e_n|$, and
\begin{equation}
 q_n[x]=\|(I-P_n)x\|^2.
 \label{eq:escaping-mode-form}
\end{equation}
Then $\Ker q_n=\M_n$ and the complementary gap equals one at every cutoff,
but $q_n$ Mosco-converges to $q_\infty[x]=\|x\|^2$, whose kernel is zero.
Thus the finite protected-basis transport in
Lemma~\ref{lem:protected-basis-transport} is a genuine hypothesis rather than
a consequence of dimension and gap information.
\end{proposition}

\begin{proof}
For every fixed $x\in\ell^2(\N)$, $\langle e_n,x\rangle\to0$, so
$q_n[x]\to\|x\|^2$; the same observation gives the Mosco liminf inequality
and constant recovery sequences.  The finite kernels are the moving lines
$\C e_n$, which converge weakly to zero and have no nonzero transported limit.
\end{proof}


\begin{thebibliography}{99}

\bibitem{ColemanMandula1967}
S.~Coleman and J.~Mandula,
\newblock All possible symmetries of the $S$ matrix,
\newblock \emph{Physical Review} \textbf{159} (1967), 1251--1256.
\newblock doi:10.1103/PhysRev.159.1251.

\bibitem{Connes1994}
A.~Connes,
\newblock \emph{Noncommutative Geometry},
\newblock Academic Press, San Diego, 1994.

\bibitem{ConnesReality1995}
A.~Connes,
\newblock Noncommutative geometry and reality,
\newblock \emph{Journal of Mathematical Physics} \textbf{36} (1995), no.~11,
6194--6231.
\newblock doi:10.1063/1.531241.

\bibitem{ConnesLott1991}
A.~Connes and J.~Lott,
\newblock Particle models and noncommutative geometry,
\newblock \emph{Nuclear Physics B -- Proceedings Supplements} \textbf{18} (1991),
29--47.
\newblock doi:10.1016/0920-5632(91)90120-4.

\bibitem{ConnesGravityMatter1996}
A.~Connes,
\newblock Gravity coupled with matter and the foundation of non-commutative
geometry,
\newblock \emph{Communications in Mathematical Physics} \textbf{182} (1996),
no.~1, 155--176.
\newblock doi:10.1007/BF02506388.

\bibitem{ChamseddineConnes1997}
A.~H.~Chamseddine and A.~Connes,
\newblock The spectral action principle,
\newblock \emph{Communications in Mathematical Physics} \textbf{186} (1997),
731--750.
\newblock doi:10.1007/s002200050126.

\bibitem{ChamseddineConnesMarcolli2007}
A.~H.~Chamseddine, A.~Connes, and M.~Marcolli,
\newblock Gravity and the standard model with neutrino mixing,
\newblock \emph{Advances in Theoretical and Mathematical Physics}
\textbf{11} (2007), no.~6, 991--1089.
\newblock doi:10.4310/ATMP.2007.v11.n6.a3.

\bibitem{ChamseddineConnes2008}
A.~H.~Chamseddine and A.~Connes,
\newblock Why the Standard Model,
\newblock \emph{Journal of Geometry and Physics} \textbf{58} (2008), no.~1,
38--47.
\newblock doi:10.1016/j.geomphys.2007.09.011.

\bibitem{vanSuijlekom2025}
W.~D.~van Suijlekom,
\newblock \emph{Noncommutative Geometry and Particle Physics},
\newblock 2nd ed., Mathematical Physics Studies, Springer, 2025.

\bibitem{Krajewski1998}
T.~Krajewski,
\newblock Classification of finite spectral triples,
\newblock \emph{Journal of Geometry and Physics} \textbf{28} (1998), 1--30.
\newblock doi:10.1016/S0393-0440(97)00068-5.

\bibitem{IochumSchuckerStephan2004}
B.~Iochum, T.~Sch\"ucker, and C.~Stephan,
\newblock On a classification of irreducible almost commutative geometries,
\newblock \emph{Journal of Mathematical Physics} \textbf{45} (2004),
5003--5041.
\newblock doi:10.1063/1.1811372.

\bibitem{vanDenDungenVanSuijlekom2012}
K.~van den Dungen and W.~D.~van Suijlekom,
\newblock Particle physics from almost commutative spacetimes,
\newblock \emph{Reviews in Mathematical Physics} \textbf{24} (2012),
1230004.
\newblock doi:10.1142/S0129055X1230004X.

\bibitem{ChamseddineConnesVanSuijlekom2013}
A.~H.~Chamseddine, A.~Connes, and W.~D.~van Suijlekom,
\newblock Beyond the spectral Standard Model: emergence of Pati--Salam
unification,
\newblock \emph{Journal of High Energy Physics} \textbf{2013} (2013),
no.~11, 132.
\newblock doi:10.1007/JHEP11(2013)132.

\bibitem{BoyleFarnsworth2018}
L.~Boyle and S.~Farnsworth,
\newblock A new algebraic structure in the standard model of particle physics,
\newblock \emph{Journal of High Energy Physics} \textbf{2018} (2018), no.~6, 071.
\newblock doi:10.1007/JHEP06(2018)071.

\bibitem{Barrett2007}
J.~W.~Barrett,
\newblock A Lorentzian version of the non-commutative geometry of the standard
model of particle physics,
\newblock \emph{Journal of Mathematical Physics} \textbf{48} (2007), 012303.
\newblock doi:10.1063/1.2408400.

\bibitem{DevastatoEtAl2018}
A.~Devastato, S.~Farnsworth, F.~Lizzi, and P.~Martinetti,
\newblock Lorentz signature and twisted spectral triples,
\newblock \emph{Journal of High Energy Physics} \textbf{2018} (2018),
no.~3, 89.
\newblock doi:10.1007/JHEP03(2018)089.

\bibitem{DAndreaDabrowski2016}
F.~D'Andrea and L.~Dabrowski,
\newblock The Standard Model in noncommutative geometry and Morita equivalence,
\newblock \emph{Journal of Noncommutative Geometry} \textbf{10} (2016),
551--578.
\newblock doi:10.4171/JNCG/242.

\bibitem{DabrowskiDAndreaSitarz2018}
L.~Dabrowski, F.~D'Andrea, and A.~Sitarz,
\newblock The Standard Model in noncommutative geometry: fundamental fermions
as internal forms,
\newblock \emph{Letters in Mathematical Physics} \textbf{108} (2018),
1323--1340.
\newblock doi:10.1007/s11005-017-1036-x.

\bibitem{DabrowskiSitarz2019}
L.~Dabrowski and A.~Sitarz,
\newblock Fermion masses, mass-mixing and the almost commutative geometry of
the Standard Model,
\newblock \emph{Journal of High Energy Physics} \textbf{2019} (2019),
no.~2, 68.
\newblock doi:10.1007/JHEP02(2019)068.

\bibitem{Howe1989}
R.~Howe,
\newblock Remarks on classical invariant theory,
\newblock \emph{Transactions of the American Mathematical Society}
\textbf{313} (1989), no.~2, 539--570.
\newblock doi:10.1090/S0002-9947-1989-0986027-X.

\bibitem{DoplicherRoberts1989}
S.~Doplicher and J.~E.~Roberts,
\newblock A new duality theory for compact groups,
\newblock \emph{Inventiones Mathematicae} \textbf{98} (1989), 157--218.
\newblock doi:10.1007/BF01388849.

\bibitem{DoplicherRoberts1990}
S.~Doplicher and J.~E.~Roberts,
\newblock Why there is a field algebra with a compact gauge group describing
the superselection structure in particle physics,
\newblock \emph{Communications in Mathematical Physics} \textbf{131} (1990),
51--107.
\newblock doi:10.1007/BF02097680.

\bibitem{Todorov2010}
I.~Todorov,
\newblock Gauge symmetry and Howe duality in 4D conformal field theory models,
\newblock \emph{Advances in Mathematical Physics} \textbf{2010} (2010),
Article ID 509538, 12 pages.
\newblock doi:10.1155/2010/509538.

\bibitem{TraylingBaylis2001}
G.~Trayling and W.~E.~Baylis,
\newblock A geometric basis for the standard-model gauge group,
\newblock \emph{Journal of Physics A: Mathematical and General} \textbf{34} (2001),
3309--3324.
\newblock doi:10.1088/0305-4470/34/15/309.

\bibitem{Daviau2017}
C.~Daviau,
\newblock Gauge group of the Standard Model in $Cl_{1,5}$,
\newblock \emph{Advances in Applied Clifford Algebras} \textbf{27} (2017), 279--290.
\newblock doi:10.1007/s00006-015-0566-5.

\bibitem{Gresnigt2020}
N.~G.~Gresnigt,
\newblock The Standard Model particle content with complete gauge symmetries from
minimal ideals of two Clifford algebras,
\newblock \emph{European Physical Journal C} \textbf{80} (2020), 583.
\newblock doi:10.1140/epjc/s10052-020-8141-1.

\bibitem{Krasnov2021}
K.~Krasnov,
\newblock SO(9) characterization of the Standard Model gauge group,
\newblock \emph{Journal of Mathematical Physics} \textbf{62} (2021), 021703.
\newblock doi:10.1063/5.0039941.

\bibitem{BaezSchwahn2026}
J.~C.~Baez and P.~Schwahn,
\newblock The Standard Model gauge group from the exceptional Jordan algebra,
\newblock arXiv:2606.15235 [math-ph], 2026.

\bibitem{FureyRoadmap2025}
N.~Furey,
\newblock An algebraic roadmap of particle theories. Part II: theoretical checkpoints,
\newblock \emph{Annalen der Physik} \textbf{537} (2025), 2400323.
\newblock doi:10.1002/andp.202400323.

\bibitem{ShaliziCrutchfield2001}
C.~R.~Shalizi and J.~P.~Crutchfield,
\newblock Computational mechanics: Pattern and prediction, structure and simplicity,
\newblock \emph{Journal of Statistical Physics} \textbf{104} (2001), 817--879.
\newblock doi:10.1023/A:1010388907793.

\bibitem{ChiribellaEtAl2009}
G.~Chiribella, G.~M.~D'Ariano, and P.~Perinotti,
\newblock Theoretical framework for quantum networks,
\newblock \emph{Physical Review A} \textbf{80} (2009), 022339.
\newblock doi:10.1103/PhysRevA.80.022339.

\bibitem{PollockEtAl2018}
F.~A.~Pollock, C.~Rodr\'iguez-Rosario, T.~Frauenheim, M.~Paternostro, and K.~Modi,
\newblock Non-Markovian quantum processes: Complete framework and efficient characterization,
\newblock \emph{Physical Review A} \textbf{97} (2018), 012127.
\newblock doi:10.1103/PhysRevA.97.012127.

\bibitem{PelissierSpacetime2026}
A.~Pelissier,
\newblock Renewal Geometry and the emergence of Lorentzian spacetime,
\newblock HAL preprint hal-05744772, 2026.
\newblock \url{https://hal.science/hal-05744772}

\bibitem{PelissierSpectral2026}
A.~Pelissier,
\newblock From Predictive Dynamics to Spectral Geometry,
\newblock HAL preprint hal-05752672, 2026.
\newblock \url{https://hal.science/hal-05752672}

\bibitem{Hucks1991}
J.~Hucks,
\newblock Global structure of the standard model, anomalies, and charge
quantization,
\newblock \emph{Physical Review D} \textbf{43} (1991), 2709--2717.
\newblock doi:10.1103/PhysRevD.43.2709.

\bibitem{Tong2017}
D.~Tong,
\newblock Line operators in the Standard Model,
\newblock \emph{Journal of High Energy Physics} \textbf{2017} (2017),
no.~7, 104.
\newblock doi:10.1007/JHEP07(2017)104.

\bibitem{GengMarshak1989}
C.~Q.~Geng and R.~E.~Marshak,
\newblock Uniqueness of quark and lepton representations in the Standard Model
from the anomalies viewpoint,
\newblock \emph{Physical Review D} \textbf{39} (1989), 693--696.
\newblock doi:10.1103/PhysRevD.39.693.

\bibitem{MinahanRamondWarner1990}
J.~A.~Minahan, P.~Ramond, and R.~C.~Warner,
\newblock Comment on anomaly cancellation in the Standard Model,
\newblock \emph{Physical Review D} \textbf{41} (1990), 715--716.
\newblock doi:10.1103/PhysRevD.41.715.

\bibitem{AlvarezGraciaBondiaMartin1995}
E.~Alvarez, J.~M.~Gracia-Bond\'ia, and C.~P.~Mart\'in,
\newblock Anomaly cancellation and the gauge group of the Standard Model in
NCG,
\newblock \emph{Physics Letters B} \textbf{364} (1995), 33--40.
\newblock doi:10.1016/0370-2693(95)01051-3.

\bibitem{Furey2018}
C.~Furey,
\newblock Three generations, two unbroken gauge symmetries, and one
eight-dimensional algebra,
\newblock \emph{Physics Letters B} \textbf{785} (2018), 84--89.
\newblock doi:10.1016/j.physletb.2018.08.032.

\bibitem{Gresnigt2026}
N.~G.~Gresnigt,
\newblock Three fermion generations and an untriplicated gauge sector with
$S_3$ family symmetry,
\newblock \emph{Physics Letters B} \textbf{880} (2026), 140782.
\newblock doi:10.1016/j.physletb.2026.140782.

\bibitem{Baez2026}
J.~C.~Baez,
\newblock Three generations in $E_7$,
\newblock arXiv:2608.06271 [math-ph], 2026.

\bibitem{Schiffler2014}
R.~Schiffler,
\newblock \emph{Quiver Representations},
\newblock CMS Books in Mathematics, Springer, Cham, 2014.

\bibitem{KuwaeShioya2003}
K.~Kuwae and T.~Shioya,
\newblock Convergence of spectral structures: a functional analytic theory and
its applications to spectral geometry,
\newblock \emph{Communications in Analysis and Geometry} \textbf{11} (2003),
599--673.
\newblock doi:10.4310/CAG.2003.v11.n4.a1.

\bibitem{BhowmickEtAl2014}
J.~Bhowmick, F.~D'Andrea, B.~Das, and L.~Dabrowski,
\newblock Quantum gauge symmetries in noncommutative geometry,
\newblock \emph{Journal of Noncommutative Geometry} \textbf{8} (2014),
433--471.
\newblock doi:10.4171/JNCG/161.

\bibitem{Maldacena1998}
J.~M.~Maldacena,
\newblock The large $N$ limit of superconformal field theories and supergravity,
\newblock \emph{Advances in Theoretical and Mathematical Physics} \textbf{2}
(1998), 231--252.
\newblock doi:10.4310/ATMP.1998.v2.n2.a1.

\bibitem{Witten1998}
E.~Witten,
\newblock Anti de Sitter space and holography,
\newblock \emph{Advances in Theoretical and Mathematical Physics} \textbf{2}
(1998), 253--290.
\newblock doi:10.4310/ATMP.1998.v2.n2.a2.

\end{thebibliography}


\end{document}
